% Traduction : la concordance des temps ; le style direct et indirect — Espagnol Tle A — Cours signé OnBuch+
% Programme officiel MINESEC. Compiler avec : tectonic E27-concordancia-tiempos-estilo-directo-indirecto.tex
\newcommand{\DOCMATIERE}{Espagnol}
\newcommand{\DOCNIVEAU}{Tle A}
\newcommand{\DOCMODULE}{Módulo 5 : Médias, communication et culture de la paix}
\newcommand{\DOCLECON}{E27}
\newcommand{\DOCTITRE}{Traduction : la concordance des temps ; le style direct et indirect}
% OnBuch+ — préambule « cours signés » ENRICHI (compilé avec tectonic / XeLaTeX)
% Système visuel « Pop » d'OnBuch+ : fond crème (jamais blanc), bordures encre
% épaisses, ombres « dures » décalées, titres Archivo Black en capitales.
% Version MATHÉMATIQUES : graphes (pgfplots/tikz), boîtes pédagogiques
% (définition, propriété, méthode, attention, le savais-tu, exercice résolu,
% exercice, TP, bilan), page de garde et signature OnBuch+ ancrée en bas.
%
% Macros attendues dans main.tex AVANT \input{preamble} :
%   \DOCMATIERE  \DOCNIVEAU  \DOCMODULE  \DOCLECON (numéro)  \DOCTITRE
\documentclass[11pt]{article}

\usepackage{fontspec}
\usepackage[spanish,french]{babel} % français = langue principale ; l'espagnol via \esp{..} / espagnol
\usepackage[a4paper,margin=2cm,top=3.4cm,bottom=2.8cm,headheight=1.4cm,footskip=1.3cm]{geometry}
\usepackage{amsmath,amssymb,amsfonts}
\usepackage{mathtools} % \xmapsto, \coloneqq... souvent utilisés par les modèles
\usepackage{cancel} % simplifications barrées (bilans de forces)
% \abs{z} (module, valeur absolue) et \norm{u} (norme), utilisés par les modèles
\providecommand{\abs}{}\let\abs\relax\DeclarePairedDelimiter{\abs}{\lvert}{\rvert}
\providecommand{\norm}{}\let\norm\relax\DeclarePairedDelimiter{\norm}{\lVert}{\rVert}
\usepackage[table]{xcolor} % [table] : fournit aussi \rowcolors (alternance de lignes)
\usepackage{pagecolor}
\usepackage{tikz}
\usetikzlibrary{arrows.meta,positioning,calc,shapes.geometric,shapes.misc,shapes.arrows,decorations.pathmorphing,decorations.pathreplacing,decorations.markings,patterns,shadows,mindmap,trees,fit,backgrounds,matrix,intersections,angles,quotes}
\usepackage{pgfplots}
\usepackage[version=4]{mhchem} % équations nucléaires \ce{^{235}_{92}U}
\usepackage[european,siunitx]{circuitikz} % schémas électriques
\pgfplotsset{compat=1.18}
\usepgfplotslibrary{fillbetween,groupplots} % aires entre courbes (calcul intégral)
\usepackage{enumitem}
\usepackage{fancyhdr}
% [most] charge toutes les bibliothèques tcolorbox (listings, minted, raster,
% documentation...) et gonfle inutilement la mémoire TeX sur un document long
% (~300 boîtes) : on ne charge que ce qui est réellement utilisé.
\usepackage[skins,breakable]{tcolorbox}
\usepackage{graphicx}
\usepackage{colortbl}
\usepackage{booktabs}
\usepackage{tabularx}
\usepackage{diagbox} % case d'en-tête diagonale des tableaux à double entrée
% Colonnes extensibles centrée / gauche / droite (C, L, R), souvent utilisées par les modèles
\newcolumntype{C}{>{\centering\arraybackslash}X}
\newcolumntype{L}{>{\raggedright\arraybackslash}X}
\newcolumntype{R}{>{\raggedleft\arraybackslash}X}
\usepackage{array}
\usepackage{multicol}
\usepackage{multirow}
\usepackage{siunitx}
% parse-numbers=false : accepte aussi \SI{2,5 \times 10^{-3}}{...}
\sisetup{locale=FR,output-decimal-marker={,},group-minimum-digits=4,parse-numbers=false}
\DeclareSIUnit\ohm{\ensuremath{\symup{\Omega}}} % U+2126 absent de la police
\DeclareSIUnit\division{div} % réglages d'oscilloscope (ms/div, V/div)
\DeclareSIUnit\tr{tr} % tours (vitesse de rotation, ex. \SI{1500}{\tr\per\minute})
\DeclareSIUnit\molar{M} % concentration molaire (ex. \SI{3}{\molar}, \SI{1}{\micro\molar})
\DeclareSIUnit\an{an} % année (le modèle confond parfois avec \year, un compteur TeX) — ex. \SI{5}{\kilo\meter\per\an}
\DeclareSIUnit\FCFA{FCFA} % franc CFA (ex. \SI{500}{\FCFA\per\kilo\gram})
\DeclareSIUnit\degreeBrix{\textdegree Brix} % teneur en sucre d'un sirop/jus (réfractométrie)
\DeclareSIUnit\month{mois} % le modèle utilise parfois \month, un compteur TeX, comme unité de durée
\usepackage{qrcode}
% \degree : commande fréquemment utilisée par les modèles pour un angle ou une
% température en dehors de \SI{}{} (non fournie par les paquets chargés).
\providecommand{\degree}{^{\circ}}
% \qed (fin de démonstration), utilisé par les modèles sans amsthm
\providecommand{\qed}{\hfill$\square$}
% \barre : double barre des valeurs interdites dans un tableau de variations (tkz-tab)
\providecommand{\barre}{\ensuremath{\|}}
% environnement proof (amsthm non chargé) utilisé par les modèles
\ifdefined\proof\else\newenvironment{proof}[1][Démonstration]{\par\noindent\textit{#1.}\ }{\qed\par}\fi
\usepackage{lastpage}
\usepackage{hyperref}
\hypersetup{hidelinks}

% ============================================================
% Palette « Pop » OnBuch+ — mêmes hex que la classe P (plus_ui.dart)
% ============================================================
\definecolor{poporange}{HTML}{F59321}
\definecolor{popdark}{HTML}{E07A0C}
\definecolor{popcream}{HTML}{FFF6E8}
\definecolor{popink}{HTML}{1C1714}
\definecolor{poppurple}{HTML}{7A5AE0}
\definecolor{popgreen}{HTML}{2E9E6A}
\definecolor{poppink}{HTML}{E0457B}
\definecolor{popblue}{HTML}{2F7DE1}
\definecolor{popgold}{HTML}{C98A12}
\definecolor{popmuted}{HTML}{8A7F76}
\definecolor{popline}{HTML}{EADFD2}
\definecolor{popgray}{HTML}{8A7F76}
\colorlet{popgrey}{popgray}
% Alias tolérants (le modèle nomme parfois les couleurs différemment)
\colorlet{popwhite}{white}
\colorlet{popblack}{popink}
\colorlet{popred}{poppink}
\colorlet{popyellow}{popgold}
% Teintes claires (fonds de tableaux/encadrés) — le modèle nomme parfois
% les couleurs avec un suffixe L (ex. popblueL) sans les avoir définies.
\colorlet{poporangeL}{poporange!20}
\colorlet{popdarkL}{popdark!20}
\colorlet{poppurpleL}{poppurple!20}
\colorlet{popgreenL}{popgreen!20}
\colorlet{poppinkL}{poppink!20}
\colorlet{popblueL}{popblue!20}
\colorlet{popgoldL}{popgold!20}
\colorlet{popredL}{popred!20}
\colorlet{popgrayL}{popgray!20}
% Teintes claires pour fonds de tableaux / zones
\colorlet{popblueL}{popblue!10!white}
\colorlet{poporangeL}{poporange!14!white}
\colorlet{popgreenL}{popgreen!12!white}
\colorlet{poppurpleL}{poppurple!12!white}
\colorlet{poppinkL}{poppink!12!white}
\colorlet{popgoldL}{popgold!14!white}

\pagecolor{popcream}

% ============================================================
% Typographie
% ============================================================
\defaultfontfeatures{Ligatures=TeX}
\setmainfont{PlusJakartaSans-Medium.ttf}[Path=../../../fonts/,BoldFont=PlusJakartaSans-Bold.ttf]
% Symboles, API et emojis absents de la police principale : repli sur DejaVu Sans (caractères actifs)
\newfontface\symfont{DejaVuSans.ttf}[Path=../../../fonts/]
\makeatletter
\newcommand\symactive[1]{\catcode#1=13 \begingroup\lccode`\~=#1 \lowercase{\endgroup\def~}{{\symfont\char#1}}}
\newcommand\symblank[1]{\catcode#1=13 \begingroup\lccode`\~=#1 \lowercase{\endgroup\def~}{}}
\newcommand\symhyph[1]{\catcode#1=13 \begingroup\lccode`\~=#1 \lowercase{\endgroup\def~}{\mbox{-}}}
\newcount\symc
\newcommand\symrange[2]{\symc=#1 \@whilenum\symc<#2 \do{\expandafter\symactive\expandafter{\the\symc}\advance\symc by 1 }}
\newcommand\symblankrange[2]{\symc=#1 \@whilenum\symc<#2 \do{\expandafter\symblank\expandafter{\the\symc}\advance\symc by 1 }}
\makeatother
\symrange{"0250}{"02FF}\symactive{"2713}\symactive{"2714}\symactive{"2717}\symactive{"2718}\symactive{"2610}\symactive{"2611}\symactive{"2612}
\symhyph{"2011}\symhyph{"2E1A}\symblankrange{"1F300}{"1FAFF}\symblank{"FE0F}
% Maths en sans-serif (Fira Math) avec chiffres et lettres droites de Plus
% Jakarta, pour que les formules se fondent dans le texte.
\usepackage[math-style=ISO,bold-style=ISO]{unicode-math}
\setmathfont{FiraMath-Regular.otf}[Path=../../../fonts/]
\setmathfont{PlusJakartaSans-Medium.ttf}[Path=../../../fonts/,range={up/num,up/{Latin,latin},\mathup}]
\usepackage{icomma} % virgule décimale sans espace en mode math
% Caractères mathématiques Unicode absents des polices de texte (Plus Jakarta,
% Archivo Black) : rendus par leur équivalent mathématique, en texte comme en
% formule — sinon ils s'affichent en carré vide (ex. « Arithmétique dans ℤ »).
\usepackage{newunicodechar}
\newunicodechar{ℝ}{\ensuremath{\mathbb{R}}}
\newunicodechar{ℤ}{\ensuremath{\mathbb{Z}}}
\newunicodechar{ℂ}{\ensuremath{\mathbb{C}}}
\newunicodechar{ℕ}{\ensuremath{\mathbb{N}}}
\newunicodechar{ℚ}{\ensuremath{\mathbb{Q}}}
\newunicodechar{𝔻}{\ensuremath{\mathbb{D}}}
\newunicodechar{α}{\ensuremath{\alpha}}
\newunicodechar{β}{\ensuremath{\beta}}
\newunicodechar{γ}{\ensuremath{\gamma}}
\newunicodechar{δ}{\ensuremath{\delta}}
\newunicodechar{ε}{\ensuremath{\varepsilon}}
\newunicodechar{θ}{\ensuremath{\theta}}
\newunicodechar{λ}{\ensuremath{\lambda}}
\newunicodechar{μ}{\ensuremath{\mu}}
\newunicodechar{σ}{\ensuremath{\sigma}}
\newunicodechar{τ}{\ensuremath{\tau}}
\newunicodechar{φ}{\ensuremath{\varphi}}
\newunicodechar{ψ}{\ensuremath{\psi}}
\newunicodechar{ω}{\ensuremath{\omega}}
\newunicodechar{Σ}{\ensuremath{\Sigma}}
% majuscules produites par \MakeUppercase dans les titres (α-aminés -> Α)
\newunicodechar{Α}{\ensuremath{\alpha}}
\newunicodechar{Β}{\ensuremath{\beta}}
\newunicodechar{Γ}{\ensuremath{\Gamma}}
\newunicodechar{Δ}{\ensuremath{\Delta}}
\newunicodechar{Ε}{\ensuremath{\varepsilon}}
\newunicodechar{Θ}{\ensuremath{\Theta}}
\newunicodechar{Λ}{\ensuremath{\Lambda}}
\newunicodechar{Μ}{\ensuremath{\mu}}
\newunicodechar{Τ}{\ensuremath{\tau}}
\newunicodechar{Φ}{\ensuremath{\Phi}}
\newunicodechar{Ψ}{\ensuremath{\Psi}}
\newunicodechar{Ω}{\ensuremath{\Omega}}
\newunicodechar{↦}{\ensuremath{\mapsto}}
\newunicodechar{∈}{\ensuremath{\in}}
\newunicodechar{∉}{\ensuremath{\notin}}
\newunicodechar{∧}{\ensuremath{\wedge}}
\newunicodechar{⇔}{\ensuremath{\Leftrightarrow}}
\newunicodechar{⟺}{\ensuremath{\Longleftrightarrow}}
\newunicodechar{⇒}{\ensuremath{\Rightarrow}}
\newunicodechar{⟹}{\ensuremath{\Longrightarrow}}
\newunicodechar{∘}{\ensuremath{\circ}}
\newunicodechar{∪}{\ensuremath{\cup}}
\newunicodechar{∩}{\ensuremath{\cap}}
\newunicodechar{≡}{\ensuremath{\equiv}}
\newunicodechar{⊂}{\ensuremath{\subset}}
\newunicodechar{⊥}{\ensuremath{\perp}}
\newunicodechar{∃}{\ensuremath{\exists}}
\newunicodechar{∀}{\ensuremath{\forall}}
\newunicodechar{∛}{\ensuremath{\sqrt[3]{\,}}}
\newunicodechar{∜}{\ensuremath{\sqrt[4]{\,}}}
\newunicodechar{⃗}{\ensuremath{{}^{\to}}}
\newunicodechar{ⁿ}{\ifmmode^{n}\else\textsuperscript{\ensuremath{n}}\fi}
\newunicodechar{ˣ}{\ifmmode^{x}\else\textsuperscript{\ensuremath{x}}\fi}
\newunicodechar{ᵇ}{\ifmmode^{b}\else\textsuperscript{\ensuremath{b}}\fi}
\newunicodechar{ʳ}{\ifmmode^{r}\else\textsuperscript{\ensuremath{r}}\fi}
\newunicodechar{ᵘ}{\ifmmode^{u}\else\textsuperscript{\ensuremath{u}}\fi}
\newunicodechar{ʸ}{\ifmmode^{y}\else\textsuperscript{\ensuremath{y}}\fi}
\newunicodechar{ᵃ}{\ifmmode^{a}\else\textsuperscript{\ensuremath{a}}\fi}
\newunicodechar{ᵉ}{\ifmmode^{e}\else\textsuperscript{\ensuremath{e}}\fi}
\newunicodechar{ᵏ}{\ifmmode^{k}\else\textsuperscript{\ensuremath{k}}\fi}
\newunicodechar{ᵖ}{\ifmmode^{p}\else\textsuperscript{\ensuremath{p}}\fi}
\newunicodechar{ᴺ}{\ifmmode^{N}\else\textsuperscript{\ensuremath{N}}\fi}
\newunicodechar{ᶜ}{\ifmmode^{c}\else\textsuperscript{\ensuremath{c}}\fi}
\newunicodechar{ʰ}{\ifmmode^{h}\else\textsuperscript{\ensuremath{h}}\fi}
\newunicodechar{ᵠ}{\ifmmode^{q}\else\textsuperscript{\ensuremath{q}}\fi}
\newunicodechar{ₙ}{\ifmmode_{n}\else\textsubscript{\ensuremath{n}}\fi}
\newunicodechar{ₐ}{\ifmmode_{a}\else\textsubscript{\ensuremath{a}}\fi}
\newunicodechar{ᵢ}{\ifmmode_{i}\else\textsubscript{\ensuremath{i}}\fi}
\newunicodechar{ᵣ}{\ifmmode_{r}\else\textsubscript{\ensuremath{r}}\fi}
\newunicodechar{ₖ}{\ifmmode_{k}\else\textsubscript{\ensuremath{k}}\fi}
\newunicodechar{ᵦ}{\ifmmode_{\beta}\else\textsubscript{\ensuremath{\beta}}\fi}
\newunicodechar{ₓ}{\ifmmode_{x}\else\textsubscript{\ensuremath{x}}\fi}
\newfontfamily\popdisplay{ArchivoBlack-Regular.ttf}[Path=../../../fonts/]
\newfontfamily\poplogo{SpaceGrotesk-Bold.ttf}[Path=../../../fonts/]
\newfontfamily\popsemi{PlusJakartaSans-SemiBold.ttf}[Path=../../../fonts/]
\newfontfamily\popxbold{PlusJakartaSans-ExtraBold.ttf}[Path=../../../fonts/]
\newfontfamily\popxboldls{PlusJakartaSans-ExtraBold.ttf}[Path=../../../fonts/,LetterSpace=3.0]

\setlist[enumerate]{leftmargin=1.7em,itemsep=5pt,topsep=4pt}
\setlist[itemize]{leftmargin=1.7em,itemsep=4pt,topsep=4pt,label={\color{poporange}\textbullet}}
\setlength{\parindent}{0pt}
\setlength{\parskip}{5pt}
\renewcommand{\arraystretch}{1.25}

% pgfplots : style Pop par défaut
\pgfplotsset{
  every axis/.append style={axis line style={popink,line width=1pt},tick style={popink},
    grid style={popline,line width=0.5pt},label style={font=\small},tick label style={font=\footnotesize}},
  popcurve/.style={poporange,line width=1.8pt,no markers,smooth},
}
% Même style disponible dans un \draw TikZ simple (hors pgfplots)
\tikzset{popcurve/.style={poporange,line width=1.8pt}}
\tikzset{popgrid/.style={popline,very thin}}
\colorlet{popcurve}{poporange} % le nom du style est parfois utilisé comme couleur

% ============================================================
% Pill / squiggle
% ============================================================
\newcommand{\poppill}[3]{%
  \tikz[baseline=-0.55ex]\node[draw=popink,line width=1.2pt,rounded corners=2.6pt,%
    fill=#2,inner xsep=6.5pt,inner ysep=3pt]{%
    \popxboldls\fontsize{7}{7}\selectfont\color{#3}\MakeUppercase{#1}};%
}
\newcommand{\popsquiggle}[1]{%
  \tikz[baseline=-0.4ex]{
    \draw[#1,line width=2.6pt,line cap=round]
      plot [smooth] coordinates {(0,0.09) (0.34,0.01) (0.68,0.21) (1.02,0.01) (1.36,0.10)};
  }%
}
% Mot-clé surligné (marqueur orange sous le texte)
\newcommand{\cle}[1]{{\popxbold\color{popdark}#1}}

% ============================================================
% En-tête / pied
% ============================================================
\pagestyle{fancy}
\fancyhf{}
\renewcommand{\headrulewidth}{0pt}
\renewcommand{\footrulewidth}{0pt}
\lhead{\raisebox{-0.15ex}{\poplogo\fontsize{13}{13}\selectfont\color{popink}OnBuch\color{poporange}+}}
\rhead{\poppill{\DOCMATIERE\ \textbullet\ \DOCNIVEAU\ \textbullet\ Leçon \DOCLECON}{white}{popink}}
\lfoot{\popxbold\fontsize{7.6}{7.6}\selectfont\color{popmuted}\MakeUppercase{Cours signé OnBuch+}\ \textbullet\ {\popsemi\DOCTITRE}}
\rfoot{\tikz[baseline=-0.6ex]\node[rounded corners=8.5pt,fill=popink,minimum height=17pt,minimum width=17pt,inner xsep=5pt,inner sep=0]{\popxbold\fontsize{8}{8}\selectfont\color{popcream}\thepage\,/\,\pageref*{LastPage}};}

% ============================================================
% Titres
% ============================================================
\newcommand{\chaptitre}[2]{%
  \begin{tcolorbox}[enhanced,colback=poporange,colframe=popink,boxrule=2.6pt,arc=11pt,
      drop shadow=popink,left=15pt,right=15pt,top=12pt,bottom=11pt,before skip=3pt,after skip=18pt]
    \poppill{#2}{popink}{popcream}\par\vspace{9pt}
    {\raggedright\hyphenpenalty=10000\exhyphenpenalty=10000\popdisplay\fontsize{22}{24}\selectfont\color{popcream}\MakeUppercase{#1}\par}
  \end{tcolorbox}
}
% Section numérotée : pastille orange avec le numéro + titre Archivo
\newcounter{popsec}
\newcommand{\coursec}[1]{%
  \stepcounter{popsec}\setcounter{popsub}{0}%
  \par\vspace{16pt}\noindent
  \begin{minipage}{\linewidth}
  \tikz[baseline=-0.7ex]\node[draw=popink,line width=1.6pt,fill=poporange,rounded corners=4pt,minimum size=22pt,inner sep=0]{\popdisplay\fontsize{12}{12}\selectfont\color{popcream}\thepopsec};%
  \hspace{8pt}{\popdisplay\fontsize{15}{17}\selectfont\color{popink}\MakeUppercase{#1}}\par
  \vspace{4pt}\noindent{\color{poporange}\rule{36pt}{3.6pt}}
  \end{minipage}\par\nopagebreak\vspace{8pt}
}
\newcounter{popsub}
\newcommand{\courssub}[1]{%
  \stepcounter{popsub}%
  \par\vspace{9pt}\noindent{\popxbold\fontsize{11.5}{13}\selectfont\color{popink}{\color{poporange}\thepopsec.\thepopsub}\hspace{6pt}#1}\par\nopagebreak\vspace{4pt}
}

% ============================================================
% Boîtes pédagogiques — même recette : fond blanc, bordure épaisse colorée,
% ombre dure encre, badge plein. Titre optionnel [#1].
% ============================================================
\tcbset{popbase/.style={breakable,enhanced,colback=white,boxrule=2.3pt,arc=9pt,
  shadow={3pt}{-3pt}{0pt}{popink},left=14pt,right=14pt,top=11pt,bottom=11pt,before skip=11pt,after skip=15pt,
  fontupper=\popsemi\fontsize{10.6}{14.5}\selectfont\color{popink},
  fontlower=\popsemi\fontsize{10.6}{14.5}\selectfont\color{popink}}}
% Coche dessinée (le glyphe ✓ n'existe pas dans les polices du document)
\newcommand{\popcheck}[1]{\tikz[baseline=-0.5ex]\draw[#1,line width=1.6pt,line cap=round,line join=round](-0.13,0.0)--(-0.04,-0.09)--(0.13,0.1);}
\providecommand{\coche}{\popcheck{popgreen}}
\providecommand{\resultat}[1]{\par\smallskip\noindent\fcolorbox{popgreen}{popgreenL}{\parbox{\dimexpr\linewidth-2\fboxsep-2\fboxrule\relax}{\textbf{Résultat.} #1}}\par\smallskip}
\newcommand{\popboxhead}[3]{\poppill{#1}{#2}{white}\ifx\relax#3\relax\else\hspace{6pt}{\popxbold\fontsize{10.6}{12}\selectfont\color{popink}#3}\fi\par\vspace{7pt}}

\newtcolorbox{aretenir}[1][]{popbase,colframe=popgreen,before upper={\popboxhead{À retenir}{popgreen}{#1}}}
% Texte en espagnol (document de lecture, dialogue, poème, extrait) : cadre
% d'étude. Un titre optionnel (source, auteur) s'affiche dans l'en-tête.
\newtcolorbox{textebox}[1][]{popbase,colframe=popblue,before upper={\popboxhead{Texto}{popblue}{#1}\begin{otherlanguage}{spanish}},after upper={\end{otherlanguage}}}
% Marques ✓ / ✗ (forme correcte / fautive) souvent utilisées par les modèles
\providecommand{\cross}{\textcolor{poppink}{\ensuremath{\times}}}
\providecommand{\xmark}{\cross}\providecommand{\crossmark}{\cross}\providecommand{\cmark}{\ensuremath{\checkmark}}
% Ligne de réponse / justification à compléter (inventée par les modèles)
\providecommand{\justification}[1]{\par\noindent\textit{Justification~:} \dotfill\par}
% \textipa (package tipa non chargé) : la police ne couvre pas l'API -> simple passage
\providecommand{\textipa}[1]{#1}
% Soulignement (ulem non chargé) : \uline, \ul
\providecommand{\uline}[1]{\underline{#1}}\providecommand{\ul}[1]{\underline{#1}}
% \glossaire{\item ...} inventée par les modèles : liste de glossaire
\providecommand{\glossaire}[1]{\begin{itemize}#1\end{itemize}}
% \alert (beamer) inventée par les modèles : texte mis en évidence
\providecommand{\alert}[1]{\textbf{\textcolor{poppink}{#1}}}
% variantes ASCII d'ordinaux écrites en macro
\providecommand{\iere}{\textsuperscript{re}}\providecommand{\ieme}{\textsuperscript{e}}\providecommand{\ere}{\textsuperscript{re}}\providecommand{\eme}{\textsuperscript{e}}\providecommand{\er}{\textsuperscript{er}}\providecommand{\ie}{\textsuperscript{e}}
% Ordinaux français écrits en macro par les modèles : 1\ière, 3\ième
\newcommand{\ière}{\textsuperscript{re}}\newcommand{\ième}{\textsuperscript{e}}
% \exemple (inventée par les modèles) : étiquette « Exemple : »
\providecommand{\exemple}{\textbf{Exemple~:}\ }
% Mot ou phrase en espagnol à l'intérieur d'un paragraphe français
\newcommand{\esp}[1]{\ifmmode\text{\textit{#1}}\else\begingroup\selectlanguage{spanish}\itshape #1\endgroup\fi}
\newtcolorbox{exemplebox}[1][]{popbase,colframe=poporange,before upper={\popboxhead{Exemple}{poporange}{#1}}}
\newtcolorbox{definition}[1][]{popbase,colframe=popblue,before upper={\popboxhead{Définition}{popblue}{#1}}}
\newtcolorbox{propriete}[1][]{popbase,colframe=poppurple,before upper={\popboxhead{Propriété}{poppurple}{#1}}}
\newtcolorbox{methode}[1][]{popbase,colframe=popgold,before upper={\popboxhead{Méthode}{popgold}{#1}}}
\newtcolorbox{attention}[1][]{popbase,colframe=poppink,before upper={\popboxhead{Attention}{poppink}{#1}}}
\newtcolorbox{experience}[1][]{popbase,colframe=popblue,colback=popblueL,before upper={\popboxhead{Expérience}{popblue}{#1}}}
% « Le savais-tu ? » : carte violette pleine (contexte, culture, Cameroun)
\newtcolorbox{savaistu}[1][]{popbase,colback=poppurple,colframe=popink,
  fontupper=\popsemi\fontsize{10.4}{14.2}\selectfont\color{white},
  before upper={\poppill{Le savais-tu\,?}{white}{poppurple}\ifx\relax#1\relax\else\hspace{6pt}{\popxbold\color{white}#1}\fi\par\vspace{7pt}}}
% Exercice résolu : énoncé en haut, \tcblower puis la solution sur fond crème
\newcounter{popexo}\newcounter{popexr}
\newtcolorbox{exoresolu}[1][]{popbase,colframe=poporange,colbacklower=popcream,
  segmentation style={popink,line width=1.2pt,dashed},
  before upper={\stepcounter{popexr}\popboxhead{Exercice résolu \thepopexr}{poporange}{#1}},
  before lower={\poppill{Solution}{popink}{popcream}\par\vspace{6pt}}}
% Exercice (non résolu en place) — la correction est dans la partie Corrigés
\newtcolorbox{exercice}[1][]{popbase,colframe=popink,
  before upper={\stepcounter{popexo}\popboxhead{Exercice \thepopexo}{popink}{#1}}}
% Corrigé
\newtcolorbox{corrige}[1][]{popbase,colframe=popgreen,colback=popgreenL,
  before upper={\popboxhead{Corrigé}{popgreen}{#1}}}
% Objectifs / prérequis / bilan
\newtcolorbox{objectifs}{popbase,colframe=popink,colback=poporangeL,
  before upper={\popboxhead{Objectifs de la leçon}{popink}{}}}
\newtcolorbox{prerequis}{popbase,colframe=popink,colback=popblueL,
  before upper={\popboxhead{Prérequis}{popblue}{}}}
\newtcolorbox{bilan}{popbase,colframe=popink,colback=white,boxrule=2.8pt,
  before upper={\popboxhead{Fiche bilan}{poporange}{L'essentiel à connaître pour l'examen}}}
% Figure encadrée avec légende
\newtcolorbox{popfigure}[1][]{enhanced,breakable=false,colback=white,colframe=popline,boxrule=1.4pt,arc=8pt,
  left=10pt,right=10pt,top=10pt,bottom=8pt,before skip=10pt,after skip=12pt,halign=center,
  lower separated=false,halign lower=center,
  fontlower=\popsemi\fontsize{9}{11.5}\selectfont\color{popmuted},
  #1}
\newcommand{\legende}[1]{\par\vspace{6pt}{\popsemi\fontsize{9}{11.5}\selectfont\color{popmuted}#1}}

% ============================================================
% Signature OnBuch+
% ============================================================
\newcommand{\poplogoinline}[1]{{\poplogo\fontsize{#1}{#1}\selectfont\color{popink}OnBuch\color{poporange}+}}
\newcommand{\signatureonbuch}{%
  \begin{tcolorbox}[enhanced,colback=popink,colframe=popink,boxrule=2.2pt,arc=10pt,
      drop shadow=poporange,left=14pt,right=14pt,top=10pt,bottom=10pt,before skip=0pt,after skip=0pt]
    \tikz[baseline=-0.7ex]\node[circle,fill=popgreen,draw=popcream,line width=1.3pt,minimum size=22pt,inner sep=0]{\popcheck{white}};%
    \hspace{9pt}\begin{minipage}[c]{0.62\linewidth}
      {\popxbold\fontsize{12}{13}\selectfont\color{popcream}Signé\ \textbullet\ {\poplogo OnBuch\color{poporange}+}}\par\vspace{2pt}
      {\raggedright\popsemi\fontsize{8}{10.5}\selectfont\color{popline}\MakeUppercase{Équipe pédagogique OnBuch+}\\\MakeUppercase{Conforme au programme officiel MINESEC}\par}
    \end{minipage}\hfill
    {\poplogo\fontsize{22}{22}\selectfont\color{popcream}OnBuch\color{poporange}+}
  \end{tcolorbox}%
}

% ============================================================
% Page de garde
% #1 titre · #2 matière · #3 niveau · #4 module · #5 n° leçon · #6 durée ·
% #7 description / objectifs (texte ou itemize)
% ============================================================
\newcommand{\pagegarde}[7]{%
  \thispagestyle{empty}
  \newgeometry{a4paper,margin=1.7cm,top=1.6cm,bottom=1.4cm}%
  \begin{tikzpicture}[remember picture,overlay]
    \fill[poporange!18] ([xshift=-3.2cm,yshift=-3.6cm]current page.north east) circle (4.6cm);
    \fill[poppurple!14] ([xshift=2.4cm,yshift=7.6cm]current page.south west) circle (3.8cm);
  \end{tikzpicture}%
  {\poplogo\fontsize{30}{30}\selectfont\color{popink}OnBuch\color{poporange}+}\hfill
  \raisebox{0.6ex}{\poppill{Cours signé \textbullet\ Premium}{popink}{popcream}}\par
  \vspace{1pt}\popsquiggle{poppurple}\par
  \vspace{8pt}
  \begin{tcolorbox}[enhanced,colback=poporange,colframe=popink,boxrule=2.8pt,arc=13pt,
      drop shadow=popink,left=18pt,right=18pt,top=12pt,bottom=13pt,after skip=10pt]
    \poppill{#2 \textbullet\ #3}{popink}{popcream}\hspace{5pt}\poppill{#4}{white}{popink}\par\vspace{8pt}
    {\popdisplay\fontsize{14}{14}\selectfont\color{popink}LEÇON #5}\par\vspace{4pt}
    {\raggedright\hyphenpenalty=10000\exhyphenpenalty=10000\popdisplay\fontsize{25}{27}\selectfont\color{popcream}\MakeUppercase{#1}\par}
    \vspace{8pt}
    {\popxbold\fontsize{9.5}{9.5}\selectfont\color{popink}\MakeUppercase{Durée conseillée : #6}}
  \end{tcolorbox}
  \begin{tcolorbox}[enhanced,colback=white,colframe=popink,boxrule=2.2pt,arc=10pt,
      drop shadow=popink,left=16pt,right=16pt,top=10pt,bottom=10pt,after skip=8pt]
    \poppill{Dans cette leçon}{poppurple}{white}\par\vspace{5pt}
    \popsemi\fontsize{9.3}{12.6}\selectfont\color{popink}#7
  \end{tcolorbox}
  \noindent\begin{minipage}[t]{0.487\linewidth}
    \begin{tcolorbox}[enhanced,colback=popgreen,colframe=popink,boxrule=2.2pt,arc=10pt,
        drop shadow=popink,left=11pt,right=11pt,top=10pt,bottom=10pt,height=3.1cm,valign=center]
      \tcbox[colback=white,colframe=popink,boxrule=1.3pt,arc=5pt,boxsep=0pt,nobeforeafter,
        left=3pt,right=3pt,top=3pt,bottom=3pt,valign=center]{\qrcode[height=1.9cm]{https://play.google.com/store/apps/details?id=com.luvvix.onbuch&hl=fr}}%
      \hspace{8pt}\begin{minipage}[c]{0.5\linewidth}
        {\popdisplay\fontsize{11}{12.5}\selectfont\color{white}\MakeUppercase{Télécharge OnBuch}}\par\vspace{4pt}
        {\raggedright\popsemi\fontsize{8.4}{11}\selectfont\color{white}Gratuit \textbullet\ Google Play\\Tuteur IA, annales, concours, résultats\par}
      \end{minipage}
    \end{tcolorbox}
  \end{minipage}\hfill
  \begin{minipage}[t]{0.487\linewidth}
    \begin{tcolorbox}[enhanced,colback=poppurple,colframe=popink,boxrule=2.2pt,arc=10pt,
        drop shadow=popink,left=11pt,right=11pt,top=10pt,bottom=10pt,height=3.1cm,valign=center]
      \tikz[baseline=-0.7ex]\node[circle,fill=white,draw=popink,line width=1.3pt,minimum size=1.6cm,inner sep=0]{\popdisplay\fontsize{24}{24}\selectfont\color{poppurple}+};%
      \hspace{8pt}\begin{minipage}[c]{0.6\linewidth}
        {\popdisplay\fontsize{11}{12.5}\selectfont\color{white}\MakeUppercase{Passe à OnBuch+}}\par\vspace{4pt}
        {\raggedright\popsemi\fontsize{8.4}{11}\selectfont\color{white}Tous les cours signés, Tuteur IA illimité, fiches PDF — onglet OnBuch+ de l'app\par}
      \end{minipage}
    \end{tcolorbox}
  \end{minipage}
  \vspace{8pt}
  \popcontenu
  \vfill
  \signatureonbuch
  \restoregeometry
  \newpage
}

% Rangée « Ce cours contient » (page de garde)
\newcommand{\poptile}[3]{% #1 couleur #2 gros texte #3 légende
  \begin{tcolorbox}[enhanced,colback=#1,colframe=popink,boxrule=2pt,arc=9pt,shadow={2.5pt}{-2.5pt}{0pt}{popink},
      left=6pt,right=6pt,top=9pt,bottom=9pt,halign=center,valign=center,height=1.75cm,width=\linewidth,nobeforeafter]
    {\popdisplay\fontsize{12.5}{13}\selectfont\color{white}\MakeUppercase{#2}}\par\vspace{3pt}
    {\centering\popsemi\fontsize{7.8}{9.5}\selectfont\color{white}#3\par}
  \end{tcolorbox}}
\newcommand{\popcontenu}{%
  \par\noindent{\popxboldls\fontsize{7.5}{7.5}\selectfont\color{popmuted}CE COURS SIGNÉ CONTIENT}\par\vspace{7pt}
  \noindent\begin{minipage}[t]{0.235\linewidth}\poptile{popblue}{Cours}{complet, illustré, pas à pas}\end{minipage}\hfill
  \begin{minipage}[t]{0.235\linewidth}\poptile{poporange}{Méthodes}{et exercices résolus}\end{minipage}\hfill
  \begin{minipage}[t]{0.235\linewidth}\poptile{poppink}{Exercices}{type examen + corrigés}\end{minipage}\hfill
  \begin{minipage}[t]{0.235\linewidth}\poptile{popgreen}{Fiche bilan}{l'essentiel à retenir}\end{minipage}\par}

% Sommaire
\newtcolorbox{sommaire}{enhanced,colback=white,colframe=popink,boxrule=2.2pt,arc=10pt,
  drop shadow=popink,left=18pt,right=18pt,top=13pt,bottom=13pt,before skip=6pt,after skip=20pt,
  before upper={\poppill{Sommaire}{poppurple}{white}\par\vspace{9pt}\popsemi\fontsize{10.6}{14.5}\selectfont\color{popink}}}

% Signature de fin de cours — poussée tout en bas de la dernière page
\newcommand{\findecours}{%
  \par\vfill\nopagebreak
  \begin{center}\popsquiggle{poporange}\end{center}\vspace{4pt}
  \signatureonbuch
}
\AtBeginDocument{\def\llbracket{\mathopen{[\mkern-3.5mu[}}\def\rrbracket{\mathclose{]\mkern-3.5mu]}}} % intervalles d'entiers (glyphe ⟦ absent de la police)
% Glyphes absents de Fira Math : redéfinis à partir de glyphes disponibles
\AtBeginDocument{%
  \def\setminus{\mathbin{\backslash}}\let\smallsetminus\setminus
  \def\vdots{\mathord{\vbox{\baselineskip3.2pt\lineskiplimit0pt\kern4pt\hbox{.}\hbox{.}\hbox{.}}}}%
  \def\ddots{\mathinner{\mkern1mu\raise6.5pt\hbox{.}\mkern2mu\raise3.6pt\hbox{.}\mkern2mu\raise0.7pt\hbox{.}\mkern1mu}}%
  \def\triangle{\mathord{\scalebox{0.9}{$\symup{\Delta}$}}}%
  \def\nabla{\mathord{\raisebox{\depth}{\scalebox{1}[-1]{$\symup{\Delta}$}}}}%
  \def\checkmark{\popcheck{popdark}}%
  \def\star{\mathbin{\ast}}\def\bigstar{\mathord{\ast}}%
  \def\diamond{\mathbin{\scalebox{0.6}{$\Diamond$}}}%
  \def\bigcup{\mathop{\mathchoice{\raisebox{-0.3em}{\scalebox{1.7}{$\cup$}}}{\scalebox{1.25}{$\cup$}}{\cup}{\cup}}\displaylimits}%
  \def\bigcap{\mathop{\mathchoice{\raisebox{-0.3em}{\scalebox{1.7}{$\cap$}}}{\scalebox{1.25}{$\cap$}}{\cap}{\cap}}\displaylimits}%
  \def\lesseqqgtr{\mathrel{\lessgtr}}\def\bigoplus{\mathop{\oplus}\displaylimits}%
}
\providecommand{\ssi}{si et seulement si} % abréviation fréquente
\AtBeginDocument{\providecommand{\wideparen}{\overparen}} % arc de cercle

%% FIGFIX-BEGIN v4 — correctifs globaux des figures (docs/onbuch-generation/tools/figfix)
\usepackage{adjustbox}\usepackage{etoolbox}
% toute figure TikZ/pgfplots plus large que la ligne est réduite à la largeur disponible
\BeforeBeginEnvironment{tikzpicture}{\begin{adjustbox}{max width=\linewidth}}
\AfterEndEnvironment{tikzpicture}{\end{adjustbox}}
% légendes pgfplots lisibles par-dessus les courbes
\pgfplotsset{every axis/.append style={legend style={fill=white,fill opacity=0.92,text opacity=1,draw=black!15,inner sep=3pt}}}
% étiquettes d'arêtes (syntaxe edge["..."]) avec fond blanc
\tikzset{every edge quotes/.append style={fill=white,inner sep=1.2pt}}
%% FIGFIX-END
\begin{document}

\pagegarde{Traduction : la concordance des temps ; le style direct et indirect}{Espagnol}{Tle A}{Módulo 5 : Médias, communication et culture de la paix}{E27}{2 h}{Cette leçon de traduction entraîne les élèves de Terminale A à maîtriser la concordance des temps à l'indicatif et au subjonctif ainsi que la transposition du style direct au style indirect au passé. À travers l'observation d'exemples authentiques, des tableaux de correspondances et des exercices de thème et de version, les apprenants acquièrent les automatismes exigés à l'épreuve d'espagnol du Baccalauréat.
\vspace{4pt}
\begin{multicols}{2}\small
\begin{itemize}
\item À la fin de cette leçon, je sais identifier le temps de la subordonnée exigé par le temps de la principale (concordance à l'indicatif).
\item À la fin de cette leçon, je sais appliquer la concordance des temps au subjonctif (présent vs imparfait).
\item À la fin de cette leçon, je sais transposer un discours direct en discours indirect introduit par un verbe au présent ou au futur.
\item À la fin de cette leçon, je sais transposer un discours direct en discours indirect introduit par un verbe au passé (déclaration, question, ordre).
\item À la fin de cette leçon, je sais employer correctement les verbes introducteurs decir, preguntar, afirmar, prometer avec leurs constructions.
\item À la fin de cette leçon, je sais transformer les marqueurs de temps et de lieu (hoy → aquel día, aquí → allí, mañana → al día siguiente).
\end{itemize}\end{multicols}}

\begin{sommaire}
\begin{multicols}{2}
\begin{enumerate}
\item Observation et découverte
\item La concordance des temps à l'indicatif
\item La concordance des temps au subjonctif
\item Transposition du style direct au style indirect au passé
\item Les verbes introducteurs et leurs constructions
\item Comparaison français/espagnol, pièges et entraînement
\item Activité d'intégration
\item Méthodes et exercices résolus
\item Exercices
\item Corrigés des exercices
\item Fiche bilan
\end{enumerate}
\end{multicols}
\end{sommaire}

\begin{objectifs}
\begin{itemize}
\item À la fin de cette leçon, je sais identifier le temps de la subordonnée exigé par le temps de la principale (concordance à l'indicatif).
\item À la fin de cette leçon, je sais appliquer la concordance des temps au subjonctif (présent vs imparfait).
\item À la fin de cette leçon, je sais transposer un discours direct en discours indirect introduit par un verbe au présent ou au futur.
\item À la fin de cette leçon, je sais transposer un discours direct en discours indirect introduit par un verbe au passé (déclaration, question, ordre).
\item À la fin de cette leçon, je sais employer correctement les verbes introducteurs decir, preguntar, afirmar, prometer avec leurs constructions.
\item À la fin de cette leçon, je sais transformer les marqueurs de temps et de lieu (hoy → aquel día, aquí → allí, mañana → al día siguiente).
\item À la fin de cette leçon, je sais comparer les systèmes français et espagnol du discours rapporté et éviter les calques.
\item À la fin de cette leçon, je sais traduire en thème et en version des phrases au style indirect passé.
\end{itemize}
\end{objectifs}

\begin{prerequis}
\begin{itemize}
\item Les temps de l'indicatif : présent, imparfait, passé simple, futur, conditionnel (Première)
\item Le subjonctif présent et ses emplois principaux (Première)
\item Le subjonctif imparfait : formation (hubiera/hubiese) (début de Terminale)
\item Les pronoms personnels et possessifs et leur adaptation dans le discours
\item La phrase complexe : proposition principale et subordonnée introduite par que
\end{itemize}
\end{prerequis}

\newpage

\coursec{Observation et découverte}

\courssub{Situation d'accroche : du micro du journaliste à la une du journal}

Imaginez la scène. Nous sommes à Yaoundé, en pleine campagne électorale. Sous le soleil du quartier Nlongkak, un journaliste de la CRTV tend son micro à un député en visite dans la circonscription. Le député, sourire aux lèvres, lève le doigt et déclare solennellement :

\esp{« Prometo que construiremos un puente antes del próximo año. »}

« Je promets que nous construirons un pont avant l'année prochaine. »

Le lendemain matin, dans les colonnes du journal, la phrase a changé de visage :

\esp{El diputado afirmó que construirían un puente antes del año siguiente.}

« Le député a affirmé qu'ils construiraient un pont avant l'année suivante. »

Regardez bien ce qui s'est passé. Le verbe \esp{prometo} (je promets) a disparu au profit d'\esp{afirmó} (il a affirmé). Le futur \esp{construiremos} (nous construirons) est devenu \esp{construirían} (ils construiraient). Le mot \esp{próximo} (prochain) s'est mué en \esp{siguiente} (suivant). Les guillemets ont sauté. Rien de tout cela n'est un hasard : c'est une mécanique de précision, réglée comme une horloge.

\begin{definition}[La problématique de la leçon]
Pourquoi \esp{prometo} devient-il \esp{afirmó} ? Pourquoi \esp{construiremos} devient-il \esp{construirían} ? Quelle loi commande ces transformations ? Répondre à ces questions, c'est découvrir la \cle{concordance des temps} et le \cle{style indirect}, deux outils indispensables pour comprendre la presse hispanophone, commenter un texte journalistique et réussir l'épreuve de traduction du Baccalauréat.
\end{definition}

\courssub{Lecture : un même fait, deux présentations}

Observons maintenant un extrait de presse original, rédigé dans le style des dépêches hispanophones. Lisez-le d'abord pour le sens, puis relisez-le en traçant d'une couleur les paroles citées et d'une autre les paroles rapportées.

\begin{textebox}[Dépêche de presse — Declaraciones del ministro]
El ministro de Obras Públicas visitó ayer la ciudad de Bata y anunció ante los periodistas un ambicioso programa de infraestructuras.

En rueda de prensa, el ministro declaró: « Inauguraremos el hospital en diciembre y hemos contratado a cien médicos. Estoy convencido de que este proyecto cambiará la vida de los ciudadanos. »

Más tarde, el portavoz del Gobierno resumió las palabras del ministro en su cuenta oficial. Según el comunicado, el ministro declaró que inaugurarían el hospital en diciembre y que habían contratado a cien médicos. Afirmó también que estaba convencido de que aquel proyecto cambiaría la vida de los ciudadanos.

Los periodistas presentes le preguntaron si el Gobierno pensaba construir más escuelas en la región. El ministro respondió que sí y prometió que anunciaría nuevas medidas la semana siguiente. Finalmente, pidió a los medios de comunicación que informaran con rigor y que no publicaran noticias falsas.

La prensa local acogió estas declaraciones con interés, aunque algunos ciudadanos recordaron que otras promesas anteriores todavía no se habían cumplido.
\end{textebox}

\textbf{Glossaire.} \esp{obras públicas} : travaux publics ; \esp{rueda de prensa} : conférence de presse ; \esp{portavoz} : porte-parole ; \esp{comunicado} : communiqué ; \esp{medios de comunicación} : médias ; \esp{con rigor} : avec rigueur ; \esp{noticias falsas} : fausses informations ; \esp{acogió} : a accueilli ; \esp{cumplir (una promesa)} : tenir (une promesse).

\textbf{Questions de compréhension.}
\begin{enumerate}
\item Quelles sont les deux promesses concrètes faites par le ministre ?
\item Relevez dans le texte toutes les formes verbales qui introduisent une parole rapportée.
\item Comparez : \esp{Inauguraremos} (style direct) et \esp{inaugurarían} (style indirect). Quel temps verbal est employé dans chaque cas ?
\item Comment la question \esp{¿El Gobierno piensa construir más escuelas?} est-elle rapportée dans le texte ? Quel mot la introduit ?
\end{enumerate}

\courssub{Repérage : les verbes introducteurs}

Premier réflexe du traducteur : repérer le \cle{verbe introducteur}, c'est-à-dire le verbe qui ouvre la porte du discours rapporté. Dans notre texte, nous avons relevé :

\begin{center}
\begin{tabularx}{\linewidth}{|X|X|X|}
\hline
\rowcolor{popblueL} \textbf{Verbo introductor} & \textbf{Français} & \textbf{Type de parole rapportée} \\
\hline
\esp{declarar} & déclarer & affirmation (déclarative) \\
\hline
\esp{afirmar} & affirmer & affirmation (déclarative) \\
\hline
\esp{prometer} & promettre & affirmation (déclarative) \\
\hline
\esp{preguntar} & demander & question (interrogative) \\
\hline
\esp{responder / contestar} & répondre & affirmation (déclarative) \\
\hline
\esp{pedir} & demander (prier) & ordre, souhait (impérative) \\
\hline
\end{tabularx}
\end{center}

\begin{attention}[Faux amis et pièges de traduction]
\esp{Preguntar} signifie « demander » au sens de « poser une question » ; « demander » au sens de « réclamer, prier de faire » se traduit par \esp{pedir}. Ne traduisez jamais « Il a demandé une augmentation » par \esp{preguntó un aumento} : dites \esp{pidió un aumento}. Autre piège : \esp{decir que} se construit toujours avec \esp{que}, même quand le français omet « que » : « Il a dit qu'il viendrait » = \esp{Dijo que vendría.} En espagnol, \esp{que} n'est \textbf{jamais} omis.
\end{attention}

\courssub{Comparaison des deux versions : qu'est-ce qui change ?}

Plaçons côte à côte les paroles du ministre et leur version rapportée.

\begin{popfigure}
\begin{tikzpicture}[
  box/.style={draw=popblue, thick, rounded corners, fill=popblueL, text width=5.6cm, align=center, font=\small, inner sep=6pt},
  lab/.style={font=\bfseries\small, text=popdark},
  arr/.style={-{Stealth[length=3mm]}, very thick, poporange}
]
\node[lab] at (-3.2,4.6) {ESTILO DIRECTO};
\node[lab] at (3.2,4.6) {ESTILO INDIRECTO};
\node[box, align=center] (a1) at (-3.2,3.4){«Inaugur\textcolor{poporange}{a}remos el hospital»\\ (futuro)};
\node[box, align=center] (b1) at (3.2,3.4){declaró que inaugur\textcolor{poporange}{ían}\\ (condicional)};
\node[box, align=center] (a2) at (-3.2,1.6){«hemos contrat\textcolor{poporange}{a}do a cien médicos»\\ (pret. perfecto)};
\node[box, align=center] (b2) at (3.2,1.6){que hab\textcolor{poporange}{ían} contratado\\ (pret. pluscuamperfecto)};
\node[box, align=center] (a3) at (-3.2,-0.2){«cambiar\textcolor{poporange}{á} la vida»\\ (futuro)};
\node[box, align=center] (b3) at (3.2,-0.2){que cambiar\textcolor{poporange}{ía} la vida\\ (condicional)};
\draw[arr] (a1) -- (b1);
\draw[arr] (a2) -- (b2);
\draw[arr] (a3) -- (b3);
\node[align=center, font=\small\itshape, text=popmuted] at (0,-1.4) {El verbo introductor está en pasado (declaró) : los tiempos retroceden.};
\end{tikzpicture}
\legende{Du style direct au style indirect : quand le verbe introducteur est au passé, chaque temps « recule » d'un cran.}
\end{popfigure}

Le constat est limpide : dès que le verbe introducteur est au \cle{passé} (\esp{declaró}, \esp{afirmó}, \esp{prometió}, \esp{preguntó}), les temps de la parole rapportée effectuent un \cle{recul} systématique :

\begin{center}
\begin{tabularx}{\linewidth}{|X|X|X|}
\hline
\rowcolor{popblueL} \textbf{Estilo directo} & \textbf{Estilo indirecto (pasado)} & \textbf{Français} \\
\hline
presente : \esp{construyo} & imperfecto : \esp{construía} & je construis → je construisais \\
\hline
futuro : \esp{construiré} & condicional : \esp{construiría} & je construirai → je construirais \\
\hline
pret. perfecto : \esp{he construido} & pluscuamperfecto : \esp{había construido} & j'ai construit → j'avais construit \\
\hline
pret. indefinido : \esp{construí} & pluscuamperfecto : \esp{había construido} & je construisis → j'avais construit \\
\hline
presente de subjuntivo : \esp{construya} & imperfecto de subjuntivo : \esp{construyera} & que je construise → que je construisisse \\
\hline
\end{tabularx}
\end{center}

\begin{propriete}[Règle dégagée : le temps du verbe introducteur commande]
Le temps du \cle{verbe introducteur} détermine le temps de la \cle{subordonnée} :
\begin{itemize}
\item Verbe introducteur au \textbf{présent} ou au \textbf{futur} (\esp{dice}, \esp{dirá}) : les temps de la parole rapportée \textbf{ne changent pas}. \esp{Dice: «Vendré» → Dice que vendrá.} « Il dit : "Je viendrai" → Il dit qu'il viendra. »
\item Verbe introducteur au \textbf{passé} (\esp{dijo}, \esp{ha dicho}, \esp{decía}) : les temps de la parole rapportée \textbf{reculent} selon le tableau ci-dessus. \esp{Dijo: «Vendré» → Dijo que vendría.} « Il a dit qu'il viendrait. »
\end{itemize}
C'est exactement la même logique qu'en français (« il dit qu'il viendra » / « il a dit qu'il viendrait ») : le conditionnel espagnol joue le rôle du « futur dans le passé », comme en français.
\end{propriete}

\courssub{La concordance des temps au subjonctif}

Lorsque la phrase directe exprime un souhait, un doute, une nécessité ou un ordre (subjonctif présent), le passage au style indirect au passé impose l'\cle{imparfait du subjonctif} (les deux formes \esp{-ra} / \esp{-se} sont correctes, \esp{-ra} étant plus fréquente à l'oral).

\begin{center}
\begin{tabularx}{\linewidth}{|X|X|X|}
\hline
\rowcolor{poppurpleL} \textbf{Estilo directo (Subj. Présent)} & \textbf{Estilo indirecto (Subj. Imparfait)} & \textbf{Français} \\
\hline
\esp{Quiero que vengas.} & \esp{Quería que vinieras / vinieses.} & Je veux que tu viennes → Je voulais que tu vinsses. \\
\hline
\esp{Es necesario que estudie.} & \esp{Fue necesario que estudiara / estudiase.} & Il faut qu'il étudie → Il fallait qu'il étudiat. \\
\hline
\esp{¡No digas mentiras!} & \esp{Me pidió que no dijera / dijese mentiras.} & Ne mens pas → Il m'a demandé de ne pas mentir. \\
\hline
\end{tabularx}
\end{center}

\begin{aretenir}[Les deux formes de l'imparfait du subjonctif]
Pour tout verbe, l'imparfait du subjonctif a deux séries de terminaisons, interchangeables :
\begin{itemize}
\item Série \esp{-ra} : \esp{cantara, comiera, viviera, hiciera, fuera, fuera} (irréguliers : \esp{ser/ir} → \esp{fuera}; \esp{estar} → \esp{estuviera}; \esp{haber} → \esp{hubiera}; \esp{poder} → \esp{pudiera}; \esp{querer} → \esp{quisiera}; \esp{saber} → \esp{supiera}; \esp{tener} → \esp{tuviera}; \esp{venir} → \esp{viniera}).
\item Série \esp{-se} : \esp{cantase, comiese, viviese, hiciese, fuese, fuese}.
\end{itemize}
\emph{Conseil :} Apprenez la série \esp{-ra} en priorité ; reconnaissez la série \esp{-se} à la lecture.
\end{aretenir}

\courssub{Les trois types de phrases rapportées}

Toute parole rapportée appartient à l'un de ces trois types. C'est la nature de la phrase d'origine qui détermine le connecteur à employer.

\begin{popfigure}
\begin{tikzpicture}[mindmap, grow cyclic,
  every node/.style={concept, font=\small, align=center},
  root concept/.append style={concept color=popblue, font=\small\bfseries},
  level 1/.append style={level distance=3.4cm, sibling angle=120},
  level 1 concept/.append style={font=\small}]
\node[root concept, align=center]{DISCURSO\\ INDIRECTO}
  child[concept color=popgreen] { node[align=center]{Declarativa\\ \esp{que} + indicativo\\ \esp{Dijo que venía.}} }
  child[concept color=poporange] { node[align=center]{Interrogativa\\ \esp{si} / palabra interrogativa\\ \esp{Preguntó si venía.}} }
  child[concept color=poppurple] { node[align=center]{Imperativa\\ infinitivo / \esp{que} + subjuntivo\\ \esp{Pidió que viniera.}} };
\end{tikzpicture}
\legende{Carte mentale : les trois visages du discours indirect.}
\end{popfigure}

\begin{definition}[Les trois types de phrases rapportées]
\begin{enumerate}
\item \textbf{Phrase déclarative} : elle est rapportée par \esp{que} + indicatif. \esp{«Estoy cansado» → Dice que está cansado.} « Il dit qu'il est fatigué. »
\item \textbf{Phrase interrogative} : question totale (réponse oui/non) rapportée par \esp{si} ; question partielle rapportée par le mot interrogatif (\esp{qué, quién, dónde, cuándo, cómo, por qué}), qui \textbf{garde son accent}. \esp{«¿Vienes?» → Pregunta si vengo.} « Il demande si je viens. » \esp{«¿Dónde vives?» → Pregunta dónde vivo.} « Il demande où j'habite. »
\item \textbf{Phrase impérative} (ordre, conseil, prière) : rapportée par l'infinitif après \esp{pedir, mandar, aconsejar, prohibir} (\esp{Me pidió ayudarle} : « Il m'a demandé de l'aider ») ou par \esp{que} + subjonctif (\esp{Me pidió que le ayudara}).
\end{enumerate}
\end{definition}

\begin{methode}[Transposer du style direct au style indirect : 4 étapes]
\begin{enumerate}
\item \textbf{Identifier le temps du verbe introducteur} : présent/futur (pas de changement) ou passé (recul des temps).
\item \textbf{Choisir le connecteur} : \esp{que} pour une déclaration, \esp{si} ou mot interrogatif pour une question, infinitif ou \esp{que} + subjonctif pour un ordre.
\item \textbf{Adapter les temps verbaux} : appliquer le tableau de recul (futur → conditionnel, présent → imparfait, parfait → plus-que-parfait, subjonctif présent → subjonctif imparfait).
\item \textbf{Adapter personnes, possessifs et marqueurs spatio-temporels} :
\begin{center}
\begin{tabular}{|l|l|}
\hline
\textbf{Direct} & \textbf{Indirect (passé)} \\
\hline
\esp{yo / tú} & \esp{él / ella} \\
\hline
\esp{mi / tu} & \esp{su} \\
\hline
\esp{hoy} & \esp{aquel día} \\
\hline
\esp{mañana} & \esp{al día siguiente / al día siguiente} \\
\hline
\esp{ayer} & \esp{el día anterior / la víspera} \\
\hline
\esp{aquí} & \esp{allí} \\
\hline
\esp{este / esta} & \esp{aquel / aquella} \\
\hline
\esp{el próximo año / la semana próxima} & \esp{el año siguiente / la semana siguiente} \\
\hline
\end{tabular}
\end{center}
\end{enumerate}
\end{methode}

\begin{exoresolu}[Transposition guidée]
Énoncé. Transposez au style indirect, le verbe introducteur étant au passé : \esp{La alcaldesa prometió: «Construiremos un mercado nuevo el próximo año y ya hemos elegido el terreno».}
\tcblower
Solution détaillée. Étape 1 : \esp{prometió} est au prétérit simple (passé) → les temps vont reculer. Étape 2 : phrase déclarative → connecteur \esp{que}. Étape 3 : \esp{construiremos} (futur) → \esp{construirían} (conditionnel) ; \esp{hemos elegido} (parfait) → \esp{habían elegido} (plus-que-parfait). Étape 4 : \esp{el próximo año} → \esp{el año siguiente} ; \esp{ya} peut être conservé. Résultat : \esp{La alcaldesa prometió que construirían un mercado nuevo el año siguiente y que ya habían elegido el terreno.} Traduction : « La mairesse a promis qu'ils construiraient un marché nouveau l'année suivante et qu'ils avaient déjà choisi le terrain. »
\end{exoresolu}

\begin{attention}[Les trois erreurs à ne jamais commettre]
\begin{itemize}
\item \textbf{Conserver guillemets et deux-points} : le style indirect les supprime obligatoirement. On n'écrit jamais \esp{Dijo: que vendría.}
\item \textbf{Oublier \esp{que}} : contrairement au français familier (« il a dit il viendrait »), l'espagnol exige toujours \esp{que} : \esp{Dijo que vendría.}
\item \textbf{Oublier le recul des temps} : après un verbe introducteur au passé, écrire \esp{Dijo que vendrá} (au lieu de \esp{vendría}) est une faute de concordance, sanctionnée à l'épreuve de traduction.
\end{itemize}
\end{attention}

\begin{savaistu}[Le discours rapporté, colonne vertébrale de la presse]
Ouvrez n'importe quel journal hispanophone — de Madrid à Malabo, de Mexico à Yaoundé pour la presse en langue espagnole — et comptez : une grande partie des phrases d'un article politique rapporte des paroles. En Guinée équatoriale, seul pays d'Afrique subsaharienne hispanophone, les dépêches officielles utilisent massivement le style indirect : \esp{El Gobierno anunció que...}, \esp{El ministro afirmó que...} Maîtriser la concordance des temps, c'est donc posséder la clé de lecture de l'actualité du monde hispanique — et un atout décisif pour l'épreuve de traduction du Bac.
\end{savaistu}

\begin{exercice}[À vous de jouer]
Transposez au style indirect avec le verbe introducteur indiqué au passé :
\begin{enumerate}
\item \esp{El profesor dijo: «El examen será difícil».}
\item \esp{María preguntó: «¿Has visto el partido de los Leones Indomables?».}
\item \esp{El padre pidió a su hijo: «Estudia más y apaga el teléfono».}
\end{enumerate}
\end{exercice}

\begin{corrige}[Exercice « À vous de jouer »]
\begin{enumerate}
\item \esp{El profesor dijo que el examen sería difícil.} « Le professeur a dit que l'examen serait difficile. » (futur → conditionnel)
\item \esp{María preguntó si había visto el partido de los Leones Indomables.} « María a demandé si j'avais vu le match des Lions Indomptables. » (question totale → \esp{si} ; parfait → plus-que-parfait)
\item \esp{El padre pidió a su hijo que estudiara más y que apagara el teléfono.} « Le père a demandé à son fils d'étudier davantage et d'éteindre le téléphone. » (impératif → \esp{que} + imparfait du subjonctif)
\end{enumerate}
\end{corrige}

\coursec{La concordance des temps à l'indicatif}

Dans la section précédente, nous avons observé un dialogue et un article de presse : les mêmes paroles apparaissaient tantôt au \cle{style direct} (entre guillemets), tantôt au \cle{style indirect} (introduites par \esp{que}). Nous avons remarqué que les temps verbaux « reculaient » d'un cran lorsque le verbe introducteur était au passé. Il est temps maintenant de formuler rigoureusement cette mécanique, qui est au cœur de toute traduction de presse et de tout commentaire de texte en Terminale.

\courssub{Observation : un fait divers rapporté par la presse}

Lisons d'abord cette dépêche originale, rédigée dans le style des journaux hispanophones. Repérez les verbes introducteurs et les temps des subordonnées.

\begin{textebox}[Dépêche de presse — Yaoundé]
El ministro de Comunicación anunció ayer en rueda de prensa que el gobierno lanzará una campaña nacional de educación mediática en todos los liceos del país. Afirmó que los jóvenes pasan demasiado tiempo en las redes sociales y que muchas veces no verifican la información que reciben. «La desinformación es un peligro para la paz social», declaró, y prometió que el programa estaría funcionando antes del próximo curso escolar.

Por su parte, la directora de un instituto de Douala explicó que sus alumnos ya habían participado en un taller sobre los medios de comunicación y añadió que los resultados habían sido muy positivos. «Los estudiantes aprenden a distinguir las noticias falsas», dijo, y aseguró que continuaría la experiencia durante todo el año.

Finalmente, un portavoz de los profesores recordó que en 2010 el país ya había intentado una reforma parecida, pero reconoció que entonces faltaban medios materiales. Prometió que los docentes colaborarían activamente y pidió que la formación fuera obligatoria para todos.
\end{textebox}

\textbf{Glossaire.} \esp{rueda de prensa} : conférence de presse ; \esp{lanzar} : lancer ; \esp{mediático} : médiatique ; \esp{verificar} : vérifier ; \esp{desinformación} : désinformation ; \esp{taller} : atelier ; \esp{noticias falsas} : fausses informations ; \esp{portavoz} : porte-parole ; \esp{docentes} : enseignants ; \esp{formación} : formation ; \esp{medios materiales} : moyens matériels.

\textbf{Questions d'observation.}
\begin{enumerate}
\item Soulignez les verbes introducteurs (\esp{anunció, afirmó, prometió, explicó, dijo, aseguró, recordó, reconoció, pidió}). À quel temps sont-ils ?
\item Classez les temps des subordonnées en deux groupes : ceux qui expriment le présent ou le passé (\esp{pasan, verifican, habían participado, habían sido, faltaban}), et ceux qui expriment le futur rapporté (\esp{lanzará, estaría, continuaría, colaborarían}).
\item Comparez \esp{«La desinformación \textbf{es} un peligro»} (style direct) avec ce qu'aurait écrit le journaliste au style indirect : \esp{declaró que la desinformación \textbf{era} un peligro}. Que constatez-vous ?
\item Remarquez enfin \esp{pidió que la formación \textbf{fuera} obligatoria} : quel mode est employé après \esp{pedir que} ? Nous y reviendrons dans la section sur le subjonctif.
\end{enumerate}

\courssub{Le principe général : le « recul » des temps}

\begin{propriete}[Principe de la concordance des temps à l'indicatif]
Le temps de la subordonnée dépend du temps du \cle{verbe introducteur} :
\begin{itemize}
\item \textbf{Verbe introducteur au présent, au futur ou au passé composé} (\esp{dice, dirá, ha dicho}) : le temps de la subordonnée est \textbf{inchangé} par rapport au discours direct. On peut employer présent, passé ou futur selon le sens.
\item \textbf{Verbe introducteur au passé} (\esp{dijo, decía, había dicho} : passé simple, imparfait, plus-que-parfait) : les temps de la subordonnée \textbf{reculent d'un cran} vers le passé :
\begin{itemize}
\item présent $\rightarrow$ \textbf{imparfait} ;
\item passé composé / parfait $\rightarrow$ \textbf{plus-que-parfait} ;
\item futur $\rightarrow$ \textbf{conditionnel} ;
\item futur antérieur $\rightarrow$ \textbf{conditionnel passé}.
\end{itemize}
\end{itemize}
\end{propriete}

\begin{exemplebox}[Le même énoncé, deux points de vue]
\begin{itemize}
\item \esp{Dice: «Vengo mañana.»} — Il dit : « Je viens demain. »\\
$\rightarrow$ \esp{Dice que viene mañana.} — Il dit qu'il vient demain. (présent conservé)
\item \esp{Dijo: «Vengo mañana.»} — Il a dit : « Je viens demain. »\\
$\rightarrow$ \esp{Dijo que venía al día siguiente.} — Il a dit qu'il venait le lendemain. (présent $\rightarrow$ imparfait)
\item \esp{Prometió: «Lo haré.»} — Il a promis : « Je le ferai. »\\
$\rightarrow$ \esp{Prometió que lo haría.} — Il a promis qu'il le ferait. (futur $\rightarrow$ conditionnel)
\item \esp{Afirmó: «Ya lo he terminado.»} — Il a affirmé : « Je l'ai déjà terminé. »\\
$\rightarrow$ \esp{Afirmó que ya lo había terminado.} — Il a affirmé qu'il l'avait déjà terminé. (passé composé $\rightarrow$ plus-que-parfait)
\end{itemize}
\end{exemplebox}

\courssub{Le tableau des correspondances}

Voici le tableau à apprendre par cœur. Il fonctionne exactement comme en français, ce qui est une excellente nouvelle pour les francophones : la mécanique est la même, seules les formes changent.

\begin{center}
\begin{tabularx}{\linewidth}{|X|X|X|}
\hline
\rowcolor{popblueL} Discours direct & Discours indirect (verbe au passé) & Exemple espagnol \\
\hline
Présent (\esp{vengo}) & Imparfait (\esp{venía}) & \esp{Dijo que venía.} — Il a dit qu'il venait. \\
\hline
Passé composé (\esp{he venido}) / parfait & Plus-que-parfait (\esp{había venido}) & \esp{Dijo que había venido.} — Il a dit qu'il était venu. \\
\hline
Imparfait (\esp{venía}) & Imparfait (\esp{venía}) — inchangé & \esp{Dijo que venía todos los días.} — Il a dit qu'il venait tous les jours. \\
\hline
Passé simple (\esp{vino}) & Passé simple ou plus-que-parfait (\esp{vino / había venido}) & \esp{Dijo que vino el día anterior.} — Il a dit qu'il était venu la veille. \\
\hline
Futur (\esp{vendré}) & Conditionnel (\esp{vendría}) & \esp{Dijo que vendría.} — Il a dit qu'il viendrait. \\
\hline
Futur antérieur (\esp{habré venido}) & Conditionnel passé (\esp{habría venido}) & \esp{Dijo que habría venido antes del lunes.} — Il a dit qu'il serait venu avant lundi. \\
\hline
\end{tabularx}
\end{center}

\begin{aretenir}[Les temps qui ne bougent pas]
L'\textbf{imparfait} et le \textbf{passé simple} de la subordonnée \cle{se maintiennent} au style indirect passé : ils sont déjà des temps du passé, ils ne peuvent pas reculer davantage.\\
\esp{«Yo vivía en Douala.»} $\rightarrow$ \esp{Dijo que vivía en Douala.} — Il a dit qu'il vivait à Douala.\\
\esp{«Ayer llegué tarde.»} $\rightarrow$ \esp{Confesó que llegó tarde el día anterior.} — Il a avoué qu'il était arrivé en retard la veille.
\end{aretenir}

\courssub{La frise du « recul » des temps}

La figure suivante visualise la mécanique : quand le verbe introducteur passe au passé, chaque temps de la subordonnée glisse d'un cran vers la gauche (vers le passé).

\begin{popfigure}
\begin{center}
\begin{tikzpicture}[font=\small]
% Axe du temps
\draw[-{Stealth[length=3mm]}, thick, popdark] (-0.5,0) -- (11,0) node[right]{tiempo};
\foreach \x/\lab in {1/{plus-que-parfait}, 3.5/{imparfait / passé simple}, 6/{présent}, 8.5/{futur}} {
  \fill[popdark] (\x,0) circle (2pt);
  \node[below, popdark, align=center] at (\x,-0.1) {\lab};
}
% Flèches de recul
\draw[-{Stealth[length=2.5mm]}, very thick, poporange] (6,0.9) to[bend right=25] (3.5,0.9);
\node[above, poporange] at (4.75,1.35) {présent $\rightarrow$ imparfait};
\draw[-{Stealth[length=2.5mm]}, very thick, poppurple] (8.5,1.9) to[bend right=25] (6,1.9);
\node[above, poppurple] at (7.25,2.35) {futur $\rightarrow$ conditionnel};
\draw[-{Stealth[length=2.5mm]}, very thick, popgreen] (3.5,2.9) to[bend right=25] (1,2.9);
\node[above, popgreen] at (2.25,3.35) {passé composé $\rightarrow$ plus-que-parfait};
\node[align=center, popblue] at (9.6,2.9) {\textbf{Verbe introducteur}\\ \textbf{au passé :}\\ recul d'un cran};
\end{tikzpicture}
\end{center}
\legende{Frise des temps : le recul d'un cran quand la principale est au passé.}
\end{popfigure}

\courssub{Le conditionnel, « futur dans le passé »}

\begin{definition}[Futur dans le passé]
Quand on rapporte au passé une parole qui contenait un futur, ce futur devient \textbf{conditionnel}. Le conditionnel espagnol (\esp{-ría, -rías, -ría, -ríamos, -ríais, -rían}) est alors un \cle{temps} de l'indicatif, appelé « futur dans le passé », exactement comme en français : \esp{Prometió que vendría.} — Il a promis qu'il viendrait.
\end{definition}

\begin{exemplebox}[Le futur rapporté au passé]
\begin{itemize}
\item \esp{El ministro anunció que la campaña \textbf{comenzaría} en enero.} — Le ministre a annoncé que la campagne commencerait en janvier.
\item \esp{Los profesores prometieron que \textbf{colaborarían} activamente.} — Les professeurs ont promis qu'ils collaboreraient activement.
\item \esp{Me escribió que me \textbf{llamaría} al día siguiente.} — Il m'a écrit qu'il m'appellerait le lendemain.
\item \esp{Pensábamos que el examen \textbf{sería} más difícil.} — Nous pensions que l'examen serait plus difficile.
\end{itemize}
\end{exemplebox}

\begin{attention}[Le conditionnel n'est pas toujours hypothétique]
Les francophones ont tendance à traduire tout conditionnel par « peut-être / éventuellement ». Erreur ! Dans \esp{dijo que vendría}, le conditionnel est un \textbf{temps} (futur dans le passé) : il faut traduire « il a dit qu'il \emph{viendrait} », et non « il viendrait peut-être ». De même, dans la presse, \esp{según fuentes oficiales, el acuerdo se firmaría mañana} signifie « l'accord serait signé demain » (futur rapporté), pas une simple hypothèse. Le contexte et le verbe introducteur au passé sont vos indices.
\end{attention}

\courssub{Comparaison avec le français}

La bonne nouvelle : le français et l'espagnol fonctionnent de la même manière. La traduction est donc presque mécanique, à condition de connaître les formes espagnoles.

\begin{center}
\begin{tabularx}{\linewidth}{|X|X|X|}
\hline
\rowcolor{popblueL} Français & Español & Remarque \\
\hline
Il dit qu'il viendra. & \esp{Dice que vendrá.} & Présent + futur : pas de changement. \\
\hline
Il a dit qu'il viendrait. & \esp{Dijo que vendría.} & Futur $\rightarrow$ conditionnel dans les deux langues. \\
\hline
Il a dit qu'il venait. & \esp{Dijo que venía.} & Présent $\rightarrow$ imparfait. \\
\hline
Il a dit qu'il était venu. & \esp{Dijo que había venido.} & Passé composé $\rightarrow$ plus-que-parfait. \\
\hline
Il a dit qu'il serait venu avant lundi. & \esp{Dijo que habría venido antes del lunes.} & Futur antérieur $\rightarrow$ conditionnel passé. \\
\hline
\end{tabularx}
\end{center}

\courssub{L'exception : les vérités générales et les faits encore valables}

\begin{propriete}[Exception à la concordance]
Quand la subordonnée exprime une \cle{vérité générale}, une loi scientifique ou un fait \textbf{encore valable au moment où l'on parle}, le présent peut (et souvent doit) se maintenir, même après un verbe au passé :\\
\esp{El profesor dijo que la tierra \textbf{gira} alrededor del sol.} — Le professeur a dit que la terre tourne autour du soleil.\\
\esp{Anunció que el museo \textbf{está} cerrado los lunes.} — Il a annoncé que le musée est fermé le lundi (c'est encore le cas).
\end{propriete}

Le français connaît la même liberté : « Galilée affirma que la Terre \emph{tourne} » est aussi correct que « tournait ». En espagnol comme en français, le présent insiste sur la validité actuelle, l'imparfait se contente de rapporter la parole.

\begin{exoresolu}[Transposition au style indirect passé]
Énoncé. Transposez au style indirect en commençant par \esp{La periodista afirmó que...} : \esp{«Los alumnos de Terminale leen la prensa todos los días y han mejorado mucho su vocabulario. El próximo mes participarán en un concurso nacional.»}
\tcblower
Solution détaillée. Le verbe introducteur \esp{afirmó} est au passé simple : chaque temps recule d'un cran.
\begin{enumerate}
\item \esp{leen} (présent) $\rightarrow$ \esp{leían} (imparfait) ;
\item \esp{han mejorado} (passé composé) $\rightarrow$ \esp{habían mejorado} (plus-que-parfait) ;
\item \esp{participarán} (futur) $\rightarrow$ \esp{participarían} (conditionnel, futur dans le passé).
\end{enumerate}
On adapte aussi les marqueurs temporels : \esp{el próximo mes} $\rightarrow$ \esp{al mes siguiente}.\\
\textbf{Résultat :} \esp{La periodista afirmó que los alumnos de Terminale leían la prensa todos los días y que habían mejorado mucho su vocabulario. Añadió que al mes siguiente participarían en un concurso nacional.} — La journaliste a affirmé que les élèves de Terminale lisaient la presse tous les jours et qu'ils avaient beaucoup amélioré leur vocabulaire. Elle a ajouté que le mois suivant ils participeraient à un concours national.
\end{exoresolu}

\begin{attention}[Pièges fréquents des francophones]
\begin{itemize}
\item \textbf{Oublier le recul} : \esp{Dijo que viene mañana} pour « il a dit qu'il venait le lendemain » est fautif si l'événement est passé ; il faut \esp{dijo que venía al día siguiente}.
\item \textbf{Calquer « demain » et « hier »} : au style indirect passé, \esp{mañana} devient \esp{al día siguiente}, \esp{ayer} devient \esp{el día anterior / la víspera}, \esp{hoy} devient \esp{aquel día}.
\item \textbf{Confondre imparfait et conditionnel} : \esp{venía} (il venait) et \esp{vendría} (il viendrait) n'ont ni le même sens ni la même fonction. Vérifiez toujours le temps du discours direct d'origine.
\item \textbf{Accorder mal le plus-que-parfait} : l'auxiliaire \esp{haber} se conjugue, mais le participe est invariable : \esp{ellos habían venido}, jamais \esp{habían venidos}.
\end{itemize}
\end{attention}

\begin{savaistu}[Le style indirect, langue de la presse hispanophone]
Dans la presse hispanophone, le style indirect est la norme absolue pour rapporter les déclarations politiques : ouvrez \emph{El País} (Espagne) ou \emph{Página/12} (Argentine) et vous lirez à chaque page \esp{el presidente anunció que..., la ministra aseguró que..., el portavoz prometió que...} suivi d'imparfaits, de plus-que-parfaits et de conditionnels. Maîtriser la concordance des temps, c'est donc pouvoir lire ces journaux sans dictionnaire — et c'est aussi une compétence directement évaluée au baccalauréat dans les épreuves de traduction et de commentaire de texte.
\end{savaistu}

\coursec{La concordance des temps au subjonctif}

Après avoir maîtrisé la concordance des temps à l'indicatif (section précédente), nous abordons maintenant son pendant indispensable : la concordance au subjonctif. En espagnol, le subjonctif n'est pas une option stylistique, c'est une \cle{obligation grammaticale} dès que la proposition principale exprime une subjectivité (volonté, émotion, doute, nécessité). La règle de concordance y est encore plus rigoureuse qu'à l'indicatif : le temps de la subordonnée dépend \emph{exclusivement} du temps de la principale.

\courssub{Les contextes déclencheurs du subjonctif : rappel synthétique}

Avant d'appliquer la concordance, il faut identifier les structures qui \cle{imposent} le subjonctif. On retient l'acronyme \cle{VEDON} (Volonté, Émotion, Doute, Ordre, Nécessité) ou, plus complet, \cle{QUESO} (Question, Urgence, Émotion, Subjectivité, Opinion). En style indirect, les verbes introducteurs de ces catégories (voir section 5) basculent la subordonnée au subjonctif.

\begin{definition}[Propositions subordonnées au subjonctif]
Le subjonctif s'emploie dans une proposition subordonnée (généralement introduite par \esp{que}) lorsque le verbe de la proposition principale exprime :
\begin{itemize}
    \item \textbf{Volonté, désir, ordre, conseil :} \esp{querer, desear, preferir, mandar, ordenar, aconsejar, insistir en, pedir, prohibir, permitir.}
    \item \textbf{Sentiment, émotion, valeur :} \esp{alegrar(se), entristecer, sorprender, extrañar, molestar, gustar, importar, parecer bien/mal, es bueno/importante/necesario que.}
    \item \textbf{Doute, négation, incertitude :} \esp{dudar, no creer, no pensar, no estar seguro, es dudoso, es improbable, negar.}
    \item \textbf{But, finalité :} \esp{para que, a fin de que, con el objeto de que, de modo que (finalité).}
    \item \textbf{Condition, concession, temps (antériorité future) :} \esp{aunque (concession), con tal que, siempre que, antes de que, hasta que, cuando (futur).}
\end{itemize}
\end{definition}

\begin{textebox}[Dialogue : Projet d'échange scolaire]
— \esp{¿Quieres que organicemos un viaje a Guinea Ecuatorial este verano?}, pregunta la profesora.
— \esp{Me parece excelente idea, pero dudo que los padres permitan que vayamos sin acompañantes}, responde Kevin.
— \esp{Es necesario que preparemos un dossier convincente. Insisto en que todos participéis.}
— \esp{Espero que el director apruebe el proyecto antes de que terminen las clases.}
\end{textebox}

\emph{Traduction :} — « Veux-tu que nous organisions un voyage en Guinée équatoriale cet été ? » demande la professeure. — « C'est une excellente idée, mais je doute que les parents permettent que nous y allions sans accompagnateurs », répond Kevin. — « Il est nécessaire que nous préparions un dossier convaincant. J'insiste pour que tous vous participiez. » — « J'espère que le directeur approuvera le projet avant que les cours ne se terminent. »

\begin{definition}[Glossaire du dialogue]
\begin{itemize}
    \item \esp{organicemos} : organisions (subj. présent, 1re pers. pl. de \esp{organizar})
    \item \esp{permitan} : permettent (subj. présent, 3e pers. pl. de \esp{permitir})
    \item \esp{vayamos} : allions (subj. présent, 1re pers. pl. de \esp{ir}, irrégulier)
    \item \esp{preparemos} : préparions (subj. présent, 1re pers. pl. de \esp{preparar})
    \item \esp{participéis} : participiez (subj. présent, 2e pers. pl. de \esp{participar})
    \item \esp{apruebe} : approuve (subj. présent, 3e pers. sg. de \esp{aprobar})
    \item \esp{terminen} : se terminent (subj. présent, 3e pers. pl. de \esp{terminar})
\end{itemize}
\end{definition}

\courssub{La règle fondamentale de concordance : présent/futur vs passé/conditionnel}

C'est le cœur du mécanisme. Le temps du subjonctif dans la subordonnée est dicté par le temps du verbe introducteur dans la principale.

\begin{propriete}[Règle d'or de la concordance au subjonctif]
\begin{itemize}
    \item \textbf{Principale au PRÉSENT (ou FUTUR, ou IMPÉRATIF)} $\rightarrow$ Subordonnée au \cle{SUBJONCTIF PRÉSENT} (antériorité : subjonctif passé).
    \item \textbf{Principale au PASSÉ (imparfait, passé simple, passé composé, plus-que-parfait) ou au CONDITIONNEL} $\rightarrow$ Subordonnée au \cle{SUBJONCTIF IMPARFAIT} (antériorité : subjonctif plus-que-parfait).
\end{itemize}
\end{propriete}

\begin{exemplebox}[Concordance : Présent $\rightarrow$ Présent / Passé $\rightarrow$ Imparfait]
\textbf{Simultanéité / Postériorité :}
\begin{itemize}
    \item \esp{Quiero que vengas.} (Je veux que tu viennes.) $\rightarrow$ \esp{Quería que vinieras.} (Je voulais que tu viennes.)
    \item \esp{Espero que llueva.} (J'espère qu'il pleuve.) $\rightarrow$ \esp{Esperaba que lloviera.} (J'espérais qu'il pleuve.)
    \item \esp{Es necesario que estudies.} (Il faut que tu étudies.) $\rightarrow$ \esp{Fue necesario que estudiaras.} (Il a fallu que tu étudies.)
\end{itemize}
\textbf{Antériorité (subjonctif passé / plus-que-parfait) :}
\begin{itemize}
    \item \esp{Me alegro de que hayas venido.} (Je suis content que tu sois venu.) $\rightarrow$ \esp{Me alegré de que hubieras venido.} (J'ai été content que tu fusses venu / que tu eusses venu.)
    \item \esp{Dudo que lo hayan hecho.} (Je doute qu'ils l'aient fait.) $\rightarrow$ \esp{Dudaba que lo hubieran hecho.} (Je doutais qu'ils l'eussent fait.)
\end{itemize}
\end{exemplebox}

\courssub{Formation du subjonctif imparfait : la méthode infaillible}

Le subjonctif imparfait (aussi appelé « pretérito imperfecto de subjuntivo ») a deux séries de terminaisons : \cle{-ra} (la plus courante) et \cle{-se} (plus littéraire, fréquente en Espagne). \textbf{La clé unique pour ne jamais se tromper, même avec les irréguliers,} est de partir de la 3\textsuperscript{e} personne du pluriel du \cle{passé simple (pretérito indefinido)}.

\begin{methode}[Former le subjonctif imparfait en 3 étapes]
\begin{enumerate}
    \item Conjuguer le verbe à la \textbf{3\textsuperscript{e} personne du pluriel du passé simple} (\esp{ellos/ellas/ustedes}).
    \item \textbf{Supprimer la terminaison -ron} : on obtient le \cle{radical imparfait du subjonctif}.
    \item \textbf{Ajouter les terminaisons} :
    \begin{center}
    \begin{tabular}{|c|c|c|c|c|c|}\hline
    \rowcolor{popblueL} Personne & -ar (hablar) & -er/-ir (comer/vivir) & -ar (-se) & -er/-ir (-se) \\ \hline
    yo & habl\textbf{ara} & com\textbf{iera} / viv\textbf{iera} & habl\textbf{ase} & com\textbf{iese} / viv\textbf{iese} \\ \hline
    tú & habl\textbf{aras} & com\textbf{ieras} / viv\textbf{ieras} & habl\textbf{ases} & com\textbf{ieses} / viv\textbf{ieses} \\ \hline
    él/ella/Ud. & habl\textbf{ara} & com\textbf{iera} / viv\textbf{iera} & habl\textbf{ase} & com\textbf{iese} / viv\textbf{iese} \\ \hline
    nosotros & habl\textbf{áramos} & com\textbf{iéramos} / viv\textbf{iéramos} & habl\textbf{ásemos} & com\textbf{iésemos} / viv\textbf{iésemos} \\ \hline
    vosotros & habl\textbf{arais} & com\textbf{ierais} / viv\textbf{ierais} & habl\textbf{aseis} & com\textbf{ieseis} / viv\textbf{ieseis} \\ \hline
    ellos/Uds. & habl\textbf{aran} & com\textbf{ieran} / viv\textbf{ieran} & habl\textbf{asen} & com\textbf{iesen} / viv\textbf{iesen} \\ \hline
    \end{tabular}
    \end{center}
\end{enumerate}
\end{methode}

\begin{aretenir}[Pourquoi cette méthode est magique : les irréguliers]
Le radical du subjonctif imparfait \textbf{est identique} au radical du passé simple (sans le \esp{-ron}). Tous les verbes irréguliers au passé simple le restent \emph{exactement de la même façon} au subjonctif imparfait. Pas de nouvelle liste à apprendre !

\begin{center}
\begin{tabularx}{\linewidth}{|X|X|X|X|}\hline
\rowcolor{popblueL} Infinitif & Passé simple (3\textsuperscript{e} pl.) & Radical subj. imparf. & Subj. imparf. (yo) \\ \hline
\esp{tener} & \esp{tuvie\textbf{ron}} & \esp{tuvier-} & \esp{tuviera / tuviese} \\ \hline
\esp{estar} & \esp{estuvie\textbf{ron}} & \esp{estuvier-} & \esp{estuviera / estuviese} \\ \hline
\esp{andar} & \esp{anduvie\textbf{ron}} & \esp{anduvier-} & \esp{anduviera / anduviese} \\ \hline
\esp{poder} & \esp{pudie\textbf{ron}} & \esp{pudier-} & \esp{pudiera / pudiese} \\ \hline
\esp{poner} & \esp{pusie\textbf{ron}} & \esp{pusier-} & \esp{pusiera / pusiese} \\ \hline
\esp{saber} & \esp{supie\textbf{ron}} & \esp{supier-} & \esp{supiera / supiese} \\ \hline
\esp{querer} & \esp{quisie\textbf{ron}} & \esp{quisier-} & \esp{quisiera / quisiese} \\ \hline
\esp{hacer} & \esp{hicie\textbf{ron}} & \esp{hicier-} & \esp{hiciera / hiciese} \\ \hline
\esp{venir} & \esp{vinie\textbf{ron}} & \esp{vinier-} & \esp{viniera / viniese} \\ \hline
\esp{decir} & \esp{dijeron} & \esp{dijer-} & \esp{dijera / dijese} \\ \hline
\esp{traer} & \esp{trajieron} & \esp{trajier-} & \esp{trajera / trajese} \\ \hline
\esp{ir / ser} & \esp{fueron} & \esp{fuer-} & \esp{fuera / fuese} \\ \hline
\end{tabularx}
\end{center}
\end{aretenir}

\begin{exemplebox}[Application de la méthode sur verbes irréguliers]
\begin{itemize}
    \item \esp{El profesor pidió que los alumnos \textbf{tuvieran} paciencia.} (Le prof a demandé que les élèves \textbf{aient} de la patience. $\leftarrow$ \esp{tuvieron} $\rightarrow$ \esp{tuvieran})
    \item \esp{No creí que él \textbf{pudiera} resolverlo.} (Je ne croyais pas qu'il \textbf{puisse} le résoudre. $\leftarrow$ \esp{pudieron} $\rightarrow$ \textbf{pudiera})
    \item \esp{Quisiera (yo) hacer una pregunta.} (Je \textbf{voudrais} poser une question. \emph{Formule de politesse, subj. imparfait de \esp{querer} pour adoucir le présent.})
\end{itemize}
\end{exemplebox}

\courssub{Les temps composés du subjonctif : passé et plus-que-parfait}

Pour exprimer l'\cle{antériorité} par rapport à la principale, on utilise les formes composées.

\begin{propriete}[Subjonctif passé et plus-que-parfait]
\begin{itemize}
    \item \textbf{Subjonctif passé (pretérito perfecto de subjuntif) :} \esp{haya} + \cle{participe passé}.\\
    S'emploie si la principale est au \textbf{présent, futur, impératif} (ou présent parfait \esp{he pedido que...}).
    \item \textbf{Subjonctif plus-que-parfait (pretérito pluscuamperfecto de subjuntif) :} \esp{hubiera / hubiese} + \cle{participe passé}.\\
    S'emploie si la principale est au \textbf{passé (quelconque) ou conditionnel}.
\end{itemize}
\end{propriete}

\begin{exemplebox}[Antériorité : Subjonctif passé vs Plus-que-parfait]
\begin{center}
\begin{tabularx}{\linewidth}{|X|X|}\hline
\rowcolor{popblueL} Principale au présent/futur $\rightarrow$ Subj. PASSÉ & Principale au passé/cond. $\rightarrow$ Subj. PLUS-QUE-PARFAIT \\ \hline
\esp{Me alegro de que \textbf{hayas aprobado}.} & \esp{Me alegré de que \textbf{hubieras aprobado}.} \\
(Je suis content que tu \textbf{aies réussi}.) & (J'ai été content que tu \textbf{eusses réussi}.) \\ \hline
\esp{Espero que \textbf{hayan llegado} bien.} & \esp{Esperaba que \textbf{hubieran llegado} bien.} \\
(J'espère qu'ils \textbf{soient arrivés} bien.) & (J'espérais qu'ils \textbf{fussent arrivés} bien.) \\ \hline
\esp{Dudo que \textbf{haya terminado} ya.} & \esp{Dudaba que \textbf{hubiera terminado} ya.} \\
(Je doute qu'il \textbf{ait fini} déjà.) & (Je doutais qu'il \textbf{eût fini} déjà.) \\ \hline
\end{tabularx}
\end{center}
\end{exemplebox}

\courssub{Le style indirect impératif : infinitif ou subjonctif ?}

C'est un point crucial pour la traduction (thème/version). Quand on rapporte un ordre, un conseil ou une interdiction (verbes : \esp{mandar, ordenar, pedir, aconsejar, prohibir, permitir, rogar}), l'espagnol offre deux constructions selon l'identité des sujets.

\begin{propriete}[Style indirect impératif : choix de la construction]
\begin{enumerate}
    \item \textbf{Sujet de l'ordre DIFFÉRENT du sujet de la principale} $\rightarrow$ \cle{Subjonctif imparfait} (si principale au passé) introduit par \esp{que}.
    \item \textbf{Sujet de l'ordre IDENTIQUE au sujet de la principale} (ou indéterminé/impersonnel) $\rightarrow$ \cle{Infinitif} (souvent précédé de préposition \esp{de} ou \esp{a} selon le verbe).
\end{enumerate}
\end{propriete}

\begin{exemplebox}[Ordre rapporté : les deux constructions]
\textbf{Verbe introducteur au passé (\esp{mandó, pidió, aconsejó}) :}
\begin{itemize}
    \item \textit{Sujets différents :} \esp{El padre mandó que \textbf{se callaran} los niños.} (Le père a ordonné que les enfants \textbf{se taisent}.) $\leftarrow$ Subj. imparf. \esp{callaran} (sujet \esp{los niños} $\neq$ \esp{el padre}).
    \item \textit{Même sujet :} \esp{El padre mandó \textbf{callarse}.} (Le père a ordonné \textbf{de se taire} / a ordonné \textbf{de se taire lui-même}.) $\leftarrow$ Infinitif \esp{callarse}.
    \item \textit{Sujet indéterminé (impersonnel) :} \esp{Se prohibió \textbf{fumar}.} (Il a été interdit \textbf{de fumer}.) $\leftarrow$ Infinitif \esp{fumar}.
    \item \textit{Avec pronom objet (me, te, le, nos, os, les) :} \esp{Me pidió que \textbf{me callara}.} (Il m'a demandé de me taire.) $\leftarrow$ Subj. imparf. \textbf{MAIS} \esp{Me pidió \textbf{callarme}.} (Il m'a demandé de me taire.) $\leftarrow$ Infinitif pronominal.
\end{itemize}
\end{exemplebox}

\begin{attention}[Piège majeur pour francophones : le calque du subjonctif présent]
En français, après un verbe au passé, on garde souvent le subjonctif présent : « Il voulait que je \textbf{vienne}. »
\emph{En espagnol, c'est \textbf{interdit}.} La concordance impose l'imparfait du subjonctif.
\begin{center}
\begin{tabular}{|l|l|}\hline
\rowcolor{poporangeL} \textbf{Français} & \textbf{Espagnol correct} \\ \hline
Il voulait que je vienne. & \esp{Quería que yo \textbf{viniera} / \textbf{viniese}.} \\ \hline
Il a ordonné qu'on parte. & \esp{Ordenó que \textbf{saliéramos} / \textbf{saliésemos}.} \\ \hline
J'aimerais que tu sois là. & \esp{Me gustaría que \textbf{estuvieras} / \textbf{estuvieses} ahí.} \\ \hline
\end{tabular}
\end{center}
\cle{Erreur fréquente :} \esp{*Quería que venga} \textcolor{red}{✗} (calque du français « voulait que je vienne »). \textbf{Correct :} \esp{Quería que viniera} \textcolor{popgreen}{✓}.
\end{attention}

\begin{popfigure}
\begin{tikzpicture}[
    node distance=1.2cm and 1.5cm,
    decision/.style={diamond, draw=popblue, thick, fill=popblueL, text width=3.5cm, align=center, inner sep=4pt},
    process/.style={rectangle, draw=popgreen, thick, fill=popgreenL, text width=3cm, align=center, rounded corners, inner sep=4pt},
    start/.style={ellipse, draw=poporange, thick, fill=poporangeL, text width=2.5cm, align=center, inner sep=4pt},
    arrow/.style={-Stealth, thick, draw=popdark}
]
\node[start] (start) {Verbe introducteur};
\node[decision, below of=start] (temps) {Principale au \textbf{Présent / Futur / Impératif} ?};
\node[process, below left=1.5cm and 0.5cm of temps, align=center] (pres){\textbf{SUBJONCTIF PRÉSENT}\\(simultanéité/postériorité)\\\esp{Quiero que vengas.}\\\textbf{SUBJONCTIF PASSÉ}\\(antériorité)\\\esp{Me alegro que hayas venido.}};
\node[process, below right=1.5cm and 0.5cm of temps, align=center] (past){\textbf{SUBJONCTIF IMPARFAIT}\\(simultanéité/postériorité)\\\esp{Quería que vinieras.}\\\textbf{SUBJONCTIF PLUS-QUE-PARFAIT}\\(antériorité)\\\esp{Me alegré que hubieras venido.}};
\draw[arrow] (start) -- (temps);
\draw[arrow] (temps) -- node[left, pos=0.6, fill=popcream, rounded corners, inner sep=2pt] {\textbf{OUI}} (pres);
\draw[arrow] (temps) -- node[right, pos=0.6, fill=popcream, rounded corners, inner sep=2pt, align=center]{\textbf{NON}\\(Passé / Conditionnel)} (past);
\end{tikzpicture}
\legende{Organigramme de décision : choisir le temps du subjonctif selon le temps de la principale.}
\end{popfigure}

\courssub{Tableau récapitulatif de conjugaison : \esp{hablar, tener, ir}}

Ce tableau juxtapose les deux séries du subjonctif imparfait (-ra / -se) pour un verbe régulier (-ar), un verbe irrégulier radical (-er) et le verbe le plus irrégulier (\esp{ir/ser}).

\begin{center}
\begin{tabular}{|l|c|c|c|c|}\hline
\rowcolor{popblueL}
\textbf{Forme} & \textbf{Subj. Présent} & \textbf{Subj. Imparfait (-ra)} & \textbf{Subj. Imparfait (-se)} & \textbf{Participe} \\ \hline
\esp{hablar} (reg.) & \esp{hable, hables, hable, hablemos, habléis, hablen} & \esp{hablar, hablaras, hablara, habláramos, hablarais, hablaran} & \esp{hablase, hablases, hablase, hablásemos, hablaseis, hablasen} & \esp{hablado} \\ \hline
\esp{tener} (irr.) & \esp{tenga, tengas, tenga, tengamos, tengáis, tengan} & \esp{tuviera, tuvieras, tuviera, tuviéramos, tuvierais, tuvieran} & \esp{tuviese, tuvieses, tuviese, tuviésemos, tuvieseis, tuviesen} & \esp{tenido} \\ \hline
\esp{ir} (sup. irr.) & \esp{vaya, vayas, vaya, vayamos, vayáis, vayan} & \esp{fuera, fueras, fuera, fuéramos, fuerais, fueran} & \esp{fuese, fueses, fuese, fuésemos, fuesen} & \esp{ido} \\ \hline
\end{tabular}
\end{center}

\begin{savaistu}[Les deux séries -ra / -se : histoire et usage]
Les deux formes coexistent depuis le Moyen Âge. La série en \cle{-ra} vient du \textbf{plus-que-parfait de l'indicatif latin} (\esp{cantav$\bar{\text{e}}$ram} $\rightarrow$ \esp{cantara}) ; la série en \cle{-se} vient du \textbf{plus-que-parfait du subjonctif latin} (\esp{cantavissem} $\rightarrow$ \esp{cantase}).
\begin{itemize}
    \item En \textbf{Espagne}, la série \cle{-se} reste très vivante à l'écrit comme à l'oral (surtout au centre et au nord).
    \item En \textbf{Amérique latine} (et en Guinée équatoriale), la série \cle{-ra} est \emph{quasi exclusive} à l'oral ; \cle{-se} apparaît dans le registre très formel ou juridique.
    \item \textbf{Conseil d'examen :} Maîtrisez la série \cle{-ra} (plus simple, universelle) ; reconnaissez la série \cle{-se} à la lecture. N'hésitez pas à utiliser \cle{-ra} dans vos productions écrites et orales.
\end{itemize}
\end{savaistu}

\courssub{Exercice résolu : concordance complète et style indirect impératif}

\begin{exoresolu}[Transposition au style indirect passé]
\textbf{Énoncé :} Transposez les phrases suivantes au style indirect en remplaçant le verbe introducteur par le temps indiqué entre parenthèses. Respectez la concordance des temps et choisissez la construction correcte (infinitif ou subjonctif) pour les ordres.

\begin{enumerate}
    \item \esp{« Venid temprano mañana », dice el entrenador. (El entrenador \textbf{ordenó}...)}
    \item \esp{« No toquéis eso », dice la madre a sus hijos. (La madre \textbf{prohibió}...)}
    \item \esp{« ¿Puedes ayudarme? », me pregunta mi amigo. (Mi amigo \textbf{me preguntó}...)}
    \item \esp{« Llegaré a las ocho », promete Juan. (Juan \textbf{prometió}...)}
    \item \esp{« ¡Ojalá nieva en Navidad! », dice la niña. (La niña \textbf{expresó el deseo}...)}
\end{enumerate}

\tcblower

\textbf{Solution détaillée :}

\begin{enumerate}
    \item \textbf{Ordre, sujets différents} (\esp{entrenador} vs \emph{vosotros} implicite dans \esp{venid}) $\rightarrow$ \cle{Subjonctif imparfait}.
    \esp{El entrenador ordenó que \textbf{vinieran} temprano al día siguiente.}
    \textit{(Note : \esp{mañana} $\rightarrow$ \esp{al día siguiente} ; \esp{venid} (2\textsuperscript{e} pl. impératif) $\rightarrow$ \esp{vinieran} (3\textsuperscript{e} pl. subj. imparf. de \esp{venir})).}

    \item \textbf{Interdiction, sujets différents} (\esp{madre} vs \emph{hijos}) $\rightarrow$ \cle{Subjonctif imparfait}.
    \esp{La madre prohibió a sus hijos que \textbf{tocaran} eso.}
    \textit{(Construction \esp{prohibir a alguien que + subj.} ; \esp{toquéis} (2\textsuperscript{e} pl. impératif/présent) $\rightarrow$ \esp{tocaran} (3\textsuperscript{e} pl. subj. imparf.)).}

    \item \textbf{Question indirecte, doute} $\rightarrow$ \cle{Subjonctif imparfait} (principale au passé).
    \esp{Mi amigo me preguntó si \textbf{podría} ayudarle.}
    \textit{(Verbe \esp{poder} : passé simple \esp{pudieron} $\rightarrow$ radical \esp{pudier-} $\rightarrow$ \esp{pudiera/pudiese}. Ici conditionnel \esp{podría} est aussi possible car \esp{preguntar si} + conditionnel = question indirecte sur capacité future. Mais \esp{subj. imparfait} exprime le doute : \esp{me preguntó si yo pudiera ayudarle} = « il m'a demandé si je \emph{pourrais} / \emph{étais capable de} l'aider »).}
    \textit{Choix retenu :} \esp{me preguntó si \textbf{podría} ayudarle} (standard pour capacité future rapportée) \textbf{OU} \esp{me preguntó si \textbf{pudiera} ayudarle} (doute sur la capacité). Les deux sont corrects selon la nuance.

    \item \textbf{Promesse, même sujet} (\esp{Juan} promet \esp{que él llegará}) $\rightarrow$ \cle{Infinitif} (simultanéité future rapportée au passé). \esp{Prometer} + infinitif.
    \esp{Juan prometió \textbf{llegar} a las ocho.}
    \textit{(Sujet identique $\rightarrow$ infinitif. \esp{llegaré} (futur) $\rightarrow$ \esp{llegar} (infinitif)).}

    \item \textbf{Souhait (\esp{ojalá})} $\rightarrow$ Subjonctif. Principale au passé (\esp{expresó}) $\rightarrow$ \cle{Subjonctif imparfait}.
    \esp{La niña expresó el deseo de que \textbf{nevara} en Navidad.}
    \textit{(\esp{Ojalá} + subj. présent $\rightarrow$ \esp{ojalá nieve}. Au passé : \esp{expresó el deseo de que} + subj. imparfait \esp{nevara}).}
\end{enumerate}
\end{exoresolu}

\courssub{Entraînement guidé : à vous de jouer}

\begin{exercice}[Concordance au subjonctif : transformation et traduction]
\textbf{A.} Mettez le verbe entre parenthèses au temps du subjonctif correct (présent, imparfait -ra, passé, plus-que-parfait) selon le temps de la principale.
\begin{enumerate}
    \item Es importante que tú \_\_\_\_\_\_\_\_ (estudiar) para el examen. (Présent)
    \item Era importante que tú \_\_\_\_\_\_\_\_ (estudiar) para el examen. (Imparfait)
    \item Me alegro de que vosotros \_\_\_\_\_\_\_\_ (aprobar) la prueba. (Passé)
    \item Me alegré de que vosotros \_\_\_\_\_\_\_\_ (aprobar) la prueba. (Plus-que-parfait)
    \item Dudo que ellos \_\_\_\_\_\_\_\_ (tener) razón. (Présent)
    \item Dudaba que ellos \_\_\_\_\_\_\_\_ (tener) razón. (Imparfait)
    \item El profesor pide que nosotros \_\_\_\_\_\_\_\_ (callar) ahora. (Présent)
    \item El profesor pidió que nosotros \_\_\_\_\_\_\_\_ (callar) en ese momento. (Imparfait)
\end{enumerate}

\textbf{B.} Traduisez en espagnol (style indirect passé).
\begin{enumerate}
    \item Le directeur a ordonné que les élèves sortissent. (\esp{ordenar que + subj.})
    \item Le directeur a ordonné de sortir. (\esp{ordenar + inf.})
    \item J'ai douté qu'il fût sincère. (\esp{dudar que + subj. plus-que-parfait})
    \item Nous voulions que tu viennes avec nous au Cameroun. (\esp{querer que + subj. imparfait})
    \item Elle m'a demandé de ne pas le dire. (\esp{pedir + inf. pronominal})
\end{enumerate}
\end{exercice}

\begin{corrige}[Exercice Concordance au subjonctif]
\textbf{A.}
\begin{enumerate}
    \item \esp{estudies} (subj. présent, 2\textsuperscript{e} sg.)
    \item \esp{estudiaras / estudiases} (subj. imparfait, 2\textsuperscript{e} sg.)
    \item \esp{hayáis aprobado} (subj. passé, 2\textsuperscript{e} pl. \emph{vosotros})
    \item \esp{hubierais / hubieseis aprobado} (subj. plus-que-parfait, 2\textsuperscript{e} pl.)
    \item \esp{tengan} (subj. présent, 3\textsuperscript{e} pl. irrégulier \esp{tener} $\rightarrow$ \esp{teng-})
    \item \esp{tuvieran / tuviesen} (subj. imparfait, 3\textsuperscript{e} pl. radical \esp{tuvier-})
    \item \esp{callemos} (subj. présent, 1\textsuperscript{e} pl. \esp{callar})
    \item \esp{calláramos / callásemos} (subj. imparfait, 1\textsuperscript{e} pl. ; accent sur \esp{á} / \esp{á} maintenu)
\end{enumerate}

\textbf{B.}
\begin{enumerate}
    \item \esp{El director ordenó que los alumnos \textbf{salieran} / \textbf{saliesen}.} (Sujets différents $\rightarrow$ subj. imparfait).
    \item \esp{El director ordenó \textbf{salir}.} (Sujet indéterminé / même autorité $\rightarrow$ infinitif).
    \item \esp{Dudé / Dudaba que él \textbf{hubiera sido} / \textbf{hubiese sido} sincero.} (Antériorité + principale passé $\rightarrow$ subj. plus-que-parfait de \esp{ser} : \esp{hubiera sido}).
    \item \esp{Queríamos / Quería que \textbf{vinieras} / \textbf{vinieses} con nosotros a Camerún.} (\esp{venir} $\rightarrow$ passé simple \esp{vinieron} $\rightarrow$ \esp{vinier-}).
    \item \esp{Me pidió que \textbf{no lo dijera} / \textbf{no lo dijese}.} (Sujets différents : \esp{ella} vs \esp{yo} $\rightarrow$ subj. imparfait) \textbf{OU} \esp{Me pidió \textbf{no decírselo} / \textbf{no decirlo}.} (Même sujet \esp{yo} $\rightarrow$ infinitif pronominal ; placement pronom : \esp{decírselo} / \esp{no decirlo}).
\end{enumerate}
\end{corrige}

\begin{attention}[Check-list avant l'épreuve (Thème / Version)]
\begin{itemize}
    \item \textbf{Verbe principal au passé/conditionnel ?} $\rightarrow$ \cle{Subjonctif imparfait (-ra)} obligatoire en subordonnée (sauf antériorité : plus-que-parfait).
    \item \textbf{Ordre / Conseil / Interdiction rapporté ?} $\rightarrow$ \cle{Sujets différents = Subjonctif imparfait (que + verbe)} ; \cle{Même sujet = Infinitif}.
    \item \textbf{Verbes irréguliers au passé simple ?} $\rightarrow$ \cle{Même radical au subj. imparfait} (\esp{tuvo $\rightarrow$ tuviera, fue $\rightarrow$ fuera, dijo $\rightarrow$ dijera}).
    \item \textbf{\esp{Ojalá} / \esp{Quizás} / \esp{Tal vez} au passé ?} $\rightarrow$ \cle{Subjonctif imparfait} (\esp{Ojalá viniera}).
    \item \textbf{Antériorité ?} $\rightarrow$ \cle{Subj. passé (haya + part.)} si principale présent/futur ; \cle{Subj. plus-que-parfait (hubiera + part.)} si principale passé/cond.
    \item \textbf{Accents !} \esp{habl\textbf{á}ramos, comi\textbf{é}ramos, vivi\textbf{é}ramos} (toujours sur la voyelle radicale de la 1\textsuperscript{e} pl.). \esp{fu\textbf{é}ramos}.
\end{itemize}
\end{attention}

\coursec{Transposition du style direct au style indirect au passé}

Dans la section précédente, nous avons vu comment le mode subjonctif obéit à la concordance des temps. Nous allons maintenant mettre cette mécanique au service d'une opération centrale en traduction : le \cle{style indirect} (ou discours rapporté, \esp{estilo indirecto}). Rapporter les paroles de quelqu'un, c'est exactement ce que font chaque jour les journalistes — et c'est une compétence exigée à l'examen, tant en thème qu'en version.

\courssub{Observation : un journaliste rapporte les paroles}

\begin{textebox}[Crónica de un partido en Yaoundé]
Ayer, después del partido, el entrenador habló con los periodistas en el estadio. «Estoy muy contento con el resultado —declaró—. Mis jugadores han trabajado mucho esta semana y mañana descansarán.» Luego añadió que estaba orgulloso de su equipo y que los jugadores descansarían al día siguiente. Una periodista le preguntó si el portero estaba lesionado. El entrenador respondió que no sabía nada todavía y que el médico lo examinaría esa noche. Antes de irse, el capitán gritó a los aficionados: «¡Venid aquí mañana!» Los periodistas escribieron que el capitán había invitado a los aficionados a ir allí al día siguiente. Finalmente, el entrenador ordenó a los jugadores que volvieran al vestuario y les pidió que no hablaran con la prensa.
\end{textebox}

\textbf{Glossaire :} \esp{el entrenador} = l'entraîneur ; \esp{el portero} = le gardien de but ; \esp{lesionado} = blessé ; \esp{los aficionados} = les supporters ; \esp{el vestuario} = le vestiaire ; \esp{la prensa} = la presse.

\textbf{Questions de compréhension :}
\begin{enumerate}
\item Quelles paroles sont rapportées au style direct ? Au style indirect ?
\item Quels temps espagnols apparaissent dans les paroles rapportées ?
\item Que sont devenus \esp{mañana}, \esp{aquí} et \esp{mis} dans le discours rapporté ?
\end{enumerate}

\textbf{Observation guidée.} Comparons deux versions :

\esp{Declaró: «Estoy muy contento.»} (style direct : paroles exactes, entre guillemets)

\esp{Declaró que estaba muy contento.} (style indirect : paroles rapportées, introduites par \esp{que})

On remarque trois changements simultanés : le verbe passe du présent (\esp{estoy}) à l'imparfait (\esp{estaba}) ; la personne passe de la première à la troisième (\esp{estoy} $\to$ \esp{estaba}, il) ; les marqueurs de temps et de lieu reculent (\esp{mañana} $\to$ \esp{al día siguiente}, \esp{aquí} $\to$ \esp{allí}). C'est exactement le même phénomène qu'en français : « Je suis content », a-t-il déclaré $\to$ Il a déclaré qu'il \emph{était} content.

\courssub{La phrase déclarative : que + indicatif}

\begin{propriete}[Règle fondamentale]
Au style indirect, une phrase \cle{déclarative} est introduite par la conjonction \esp{que} (jamais omise en espagnol soigné) suivie d'un verbe à l'\cle{indicatif}, dont le temps obéit à la concordance des temps vue dans les sections précédentes :

\begin{center}
\begin{tabular}{|l|l|}
\hline
\rowcolor{popblueL} \textbf{Estilo directo} & \textbf{Estilo indirecto (verbe au passé)} \\
\hline
presente : \esp{trabajo} & imperfecto : \esp{trabajaba} \\
\hline
pretérito perfecto : \esp{he trabajado} & pluscuamperfecto : \esp{había trabajado} \\
\hline
pretérito indefinido : \esp{trabajé} & imperfecto / pluscuamperfecto : \esp{trabajaba} / \esp{había trabajado} \\
\hline
futuro : \esp{trabajaré} & condicional : \esp{trabajaría} \\
\hline
\end{tabular}
\end{center}
\end{propriete}

\begin{exemplebox}[Exemples traduits]
\esp{«Trabajo mucho», dijo.} $\to$ \esp{Dijo que trabajaba mucho.} — « Je travaille beaucoup », a-t-il dit $\to$ Il a dit qu'il travaillait beaucoup.

\esp{«He visto al profesor», afirmó Ana.} $\to$ \esp{Ana afirmó que había visto al profesor.} — Ana a affirmé qu'elle avait vu le professeur.

\esp{«Mañana iremos a Douala», prometieron.} $\to$ \esp{Prometieron que irían a Douala al día siguiente.} — Ils ont promis qu'ils iraient à Douala le lendemain.
\end{exemplebox}

\begin{attention}[Ne pas omettre « que »]
En français familier on entend parfois « il a dit il venait », mais en espagnol comme en français soutenu, la conjonction est obligatoire : \esp{Dijo que venía.} De plus, \esp{que} ne prend jamais d'accent ici : on écrit \esp{que} (conjonction), à ne pas confondre avec \esp{qué} (interrogatif).
\end{attention}

\courssub{La phrase interrogative totale : si + indicatif}

Une \cle{question totale} (à laquelle on répond oui ou non) est rapportée au moyen de \esp{si}, exactement comme le français « si ».

\begin{exemplebox}[Exemple de référence]
\esp{Preguntó: «¿Vienes a la fiesta?»} $\to$ \esp{Preguntó si venía a la fiesta.} — Il a demandé : « Tu viens à la fête ? » $\to$ Il a demandé si je venais à la fête.
\end{exemplebox}

\begin{propriete}[Après « si » interrogatif : indicatif, jamais subjonctif]
Après \esp{si} introduisant une question indirecte, l'espagnol exige l'\cle{indicatif}, \textbf{jamais le subjonctif}. On dit \esp{Preguntó si vendría.} (Il a demandé s'il viendrait.) et non \emph{*preguntó si viniera}. C'est une différence capitale avec le \esp{si} conditionnel, qui lui rejette le futur et le conditionnel : ne confondez pas \esp{si vendría} (question indirecte : « s'il viendrait ») et \esp{si viniera} (condition : « s'il venait »).
\end{propriete}

\begin{exemplebox}[Autres exemples]
\esp{«¿Has terminado?», me preguntó.} $\to$ \esp{Me preguntó si había terminado.} — Il m'a demandé si j'avais terminé.

\esp{«¿Estáis cansados?», preguntó la profesora.} $\to$ \esp{La profesora preguntó si estábamos cansados.} — La professeure a demandé si nous étions fatigués.
\end{exemplebox}

\courssub{La phrase interrogative partielle : mot interrogatif conservé}

Quand la question directe contient un mot interrogatif (\esp{qué}, \esp{quién}, \esp{cuándo}, \esp{dónde}, \esp{por qué}, \esp{cuánto}, \esp{cómo}), ce mot est \cle{conservé} tel quel dans la question indirecte, suivi du verbe à l'indicatif.

\begin{exemplebox}[Exemple de référence]
\esp{«¿Dónde vives?»} $\to$ \esp{Me preguntó dónde vivía.} — « Où habites-tu ? » $\to$ Il m'a demandé où j'habitais.
\end{exemplebox}

\begin{attention}[L'accent écrit est conservé]
Dans la question indirecte, le mot interrogatif \textbf{garde son accent écrit} (\esp{dónde}, \esp{cuándo}, \esp{qué}, \esp{quién}) même s'il n'y a plus de points d'interrogation. C'est ce qui le distingue du relatif : comparez \esp{Me preguntó dónde vivía.} (question indirecte : où j'habitais) et \esp{la ciudad donde vivía} (relatif : la ville où j'habitais). Erreur fréquente de francophone : écrire \emph{*Me preguntó donde vivía} sans accent.
\end{attention}

\begin{propriete}[Ordre des mots dans la question indirecte]
L'espagnol conserve dans la question indirecte l'ordre fréquent \cle{verbe + sujet} : \esp{Me preguntó dónde vivía Juan.} (Il m'a demandé où habitait Juan.) Le français impose l'ordre sujet--verbe (« où Juan habitait ») : en version, rétablissez l'ordre français ; en thème, vous pouvez placer le sujet après le verbe, ce qui est la tournure la plus naturelle en espagnol.
\end{propriete}

\begin{exemplebox}[Série d'exemples]
\esp{«¿Qué quieres?»} $\to$ \esp{Me preguntó qué quería.} — Il m'a demandé ce que je voulais.

\esp{«¿Cuándo llegará el tren?»} $\to$ \esp{Preguntó cuándo llegaría el tren.} — Il a demandé quand le train arriverait.

\esp{«¿Por qué no viniste?»} $\to$ \esp{Me preguntó por qué no había venido.} — Il m'a demandé pourquoi je n'étais pas venu.

\esp{«¿Quién ha roto la ventana?»} $\to$ \esp{Preguntó quién había roto la ventana.} — Il a demandé qui avait cassé la fenêtre.
\end{exemplebox}

\courssub{La phrase impérative : infinitif ou que + subjonctif imparfait}

Pour rapporter un \cle{ordre} ou une \cle{prière}, l'espagnol dispose de deux constructions :

\begin{propriete}[Deux constructions pour l'impératif rapporté]
\begin{enumerate}
\item \cle{Verbe introducteur + infinitif} : \esp{ordenar}, \esp{mandar}, \esp{pedir}, \esp{decir} + complément de personne + infinitif. \esp{Le ordenó cerrar la puerta.} (Il lui ordonna de fermer la porte.)
\item \cle{Verbe introducteur + que + subjonctif imparfait} : \esp{Ordenó que cerraran la puerta.} (Il ordonna qu'on ferme la porte / qu'ils ferment la porte.) Le subjonctif imparfait est exigé parce que le verbe principal est au passé (concordance des temps).
\end{enumerate}
\end{propriete}

\begin{attention}[« Decir » : indicatif ou subjonctif ?]
Le verbe \esp{decir} change de construction selon le sens : \esp{decir que} + \textbf{indicatif} rapporte une \emph{information} (\esp{Dijo que estaba cansado.} = Il a dit qu'il était fatigué), tandis que \esp{decir a alguien que} + \textbf{subjonctif} rapporte un \emph{ordre} (\esp{Le dijo que se callara.} = Il lui dit de se taire). Le français « dire à quelqu'un de + infinitif » se traduit donc par \esp{decir a alguien que} + subjonctif ou par \esp{decir a alguien} + infinitif : \esp{Le dijo que se callara} / \esp{Le dijo callarse.} Ne traduisez jamais « il lui a dit de partir » par \emph{*le dijo que se iba}.
\end{attention}

\begin{methode}[Choisir la bonne construction pour rapporter un ordre]
\begin{enumerate}
\item Repérez l'impératif dans le style direct : \esp{«Cierra la puerta», dijo.}
\item Identifiez le destinataire de l'ordre. S'il est clair et unique, préférez l'\textbf{infinitif}, plus léger : \esp{Le ordenó cerrar la puerta.} (Il lui ordonna de fermer la porte.)
\item Si le destinataire est collectif, implicite ou si vous voulez insister sur l'ordre, utilisez \esp{que} + subjonctif imparfait : \esp{Dijo que cerrara la puerta.} (Il dit qu'elle ferme la porte.)
\item Pour un ordre \textbf{négatif}, l'infinitif reste possible : \esp{Le pidió no hablar.} ; avec \esp{que} : \esp{Le pidió que no hablara.} (Il lui demanda de ne pas parler.)
\end{enumerate}
\end{methode}

\begin{exemplebox}[Ordres rapportés]
\esp{«¡Venid aquí mañana!», gritó.} $\to$ \esp{Gritó que fueran allí al día siguiente.} — Il cria qu'ils aillent là-bas le lendemain.

\esp{«No toquéis nada», ordenó el policía.} $\to$ \esp{El policía ordenó no tocar nada} / \esp{ordenó que no tocarais nada.} — Le policier ordonna de ne rien toucher.

\esp{«Ayúdame, por favor», me pidió.} $\to$ \esp{Me pidió que la ayudara.} — Elle m'a demandé de l'aider.
\end{exemplebox}

\courssub{L'adaptation des personnes et des marqueurs}

Rapporter un discours, c'est changer de point de vue : le « je » de l'énonciation directe devient « il » ou « elle », et tous les repères de personne, de temps et de lieu se déplacent avec lui.

\begin{propriete}[Adaptation des personnes, possessifs et démonstratifs]
\begin{itemize}
\item \cle{Personnes} : \esp{yo} $\to$ \esp{él/ella} ; \esp{tú} $\to$ \esp{yo} ou \esp{él} selon qui rapporte ; les désinences verbales suivent : \esp{«Tengo prisa»} $\to$ \esp{Dijo que tenía prisa.} (Il a dit qu'il était pressé.)
\item \cle{Possessifs} : \esp{mi} $\to$ \esp{su} ; \esp{nuestro} $\to$ \esp{su} / \esp{nuestro} selon le contexte : \esp{«Mi hermano está enfermo»} $\to$ \esp{Dijo que su hermano estaba enfermo.} (Il a dit que son frère était malade.)
\item \cle{Démonstratifs} : le proche devient lointain : \esp{este} $\to$ \esp{ese} ou \esp{aquel} : \esp{«Este libro es mío»} $\to$ \esp{Dijo que aquel libro era suyo.} (Il a dit que ce livre-là était à lui.)
\item \cle{Verbes de déplacement} : le point de vue change aussi pour \esp{venir} et \esp{traer}, qui deviennent souvent \esp{ir} et \esp{llevar} : \esp{«Vendré aquí»} $\to$ \esp{Dijo que iría allí.} (Il a dit qu'il irait là-bas.)
\end{itemize}
\end{propriete}

\begin{propriete}[Transformation des marqueurs temporels et spatiaux]
Quand le verbe introducteur est au passé, les repères de temps et de lieu du locuteur reculent :

\begin{center}
\begin{tabularx}{\linewidth}{|X|X|X|}
\hline
\rowcolor{popblueL} \textbf{Estilo directo} & \textbf{Estilo indirecto} & \textbf{Français} \\
\hline
\esp{hoy} & \esp{aquel día} & aujourd'hui $\to$ ce jour-là \\
\hline
\esp{ayer} & \esp{el día anterior} & hier $\to$ la veille \\
\hline
\esp{mañana} (demain) & \esp{al día siguiente} & demain $\to$ le lendemain \\
\hline
\esp{ahora} & \esp{entonces} & maintenant $\to$ alors \\
\hline
\esp{la semana próxima} & \esp{la semana siguiente} & la semaine prochaine $\to$ la semaine suivante \\
\hline
\esp{aquí} & \esp{allí} & ici $\to$ là(-bas) \\
\hline
\esp{este lugar} & \esp{aquel lugar} & cet endroit $\to$ cet endroit-là \\
\hline
\end{tabularx}
\end{center}
\end{propriete}

\begin{popfigure}
\begin{center}
\begin{tikzpicture}
\draw[thick, popline] (0,0) rectangle (7.4,4.6);
\draw[fill=popblue, popblue] (0,3.9) rectangle (3.7,4.6);
\draw[fill=poporange, poporange] (3.7,3.9) rectangle (7.4,4.6);
\node[text=white, font=\bfseries] at (1.85,4.25) {Estilo directo};
\node[text=white, font=\bfseries] at (5.55,4.25) {Estilo indirecto};
\node[font=\itshape] at (1.85,3.4) {hoy};
\node[font=\itshape] at (5.55,3.4) {aquel día};
\node[font=\itshape] at (1.85,2.75) {ayer};
\node[font=\itshape] at (5.55,2.75) {el día anterior};
\node[font=\itshape] at (1.85,2.1) {mañana};
\node[font=\itshape] at (5.55,2.1) {al día siguiente};
\node[font=\itshape] at (1.85,1.45) {ahora};
\node[font=\itshape] at (5.55,1.45) {entonces};
\node[font=\itshape] at (1.85,0.8) {aquí};
\node[font=\itshape] at (5.55,0.8) {allí};
\node[font=\itshape] at (1.85,0.25) {este lugar};
\node[font=\itshape] at (5.55,0.25) {aquel lugar};
\foreach \y in {0.55,1.2,1.85,2.5,3.15,3.9} {\draw[popline] (0,\y) -- (7.4,\y);}
\draw[popline] (3.7,0) -- (3.7,4.6);
\foreach \y in {3.4,2.75,2.1,1.45,0.8,0.25} {\draw[-{Stealth}, thick, popdark] (2.9,\y) -- (4.2,\y);}
\end{tikzpicture}
\end{center}
\legende{Tableau des marqueurs temporels et spatiaux : du style direct au style indirect.}
\end{popfigure}

\begin{savaistu}[« Mañana » : demain ou matin ?]
Le mot \esp{mañana} a deux sens : « demain » (adverbe de temps) et « matin » (nom féminin : \esp{la mañana}). Dans le discours rapporté, les deux se transforment différemment : « demain » devient \esp{al día siguiente} (\esp{Dijo que vendría al día siguiente.} = Il a dit qu'il viendrait le lendemain), tandis que « le matin » devient \esp{por la mañana} ou \esp{aquella mañana} (\esp{Dijo que trabajaba por la mañana.} = Il a dit qu'il travaillait le matin). Seul le contexte permet de trancher : méfiez-vous de ce piège en version !
\end{savaistu}

\courssub{Schéma de synthèse : choisir le connecteur et le mode}

\begin{popfigure}
\begin{center}
\begin{tikzpicture}[
  grow=right,
  level 1/.style={sibling distance=2.6cm, level distance=3.4cm},
  level 2/.style={sibling distance=1.6cm, level distance=3.2cm},
  every node/.style={draw, rounded corners, align=center, font=\small, inner sep=4pt}
]
\node[fill=popblue, text=white, font=\bfseries, align=center]{Frase en\\estilo directo}
  child { node[fill=popgreenL, align=center]{imperativa\\(orden, ruego)}
    child { node[fill=popgreen!40, align=center]{infinitivo\\o \emph{que} +\\subjuntivo imperfecto} }
  }
  child { node[fill=poporangeL] {interrogativa}
    child { node[fill=poporange!40, align=center]{total : \emph{si}\\+ indicativo} }
    child { node[fill=poporange!40, yshift=-1.1cm, align=center]{parcial : \emph{qué, dónde...}\\+ indicativo} }
  }
  child { node[fill=poppurpleL] {declarativa}
    child { node[fill=poppurple!40, align=center]{\emph{que}\\+ indicativo} }
  };
\end{tikzpicture}
\end{center}
\legende{Type de phrase $\to$ connecteur $\to$ mode verbal au style indirect.}
\end{popfigure}

\begin{aretenir}[L'essentiel de la transposition]
\begin{itemize}
\item Déclarative : \esp{que} + indicatif (concordance des temps).
\item Interrogative totale : \esp{si} + indicatif, \textbf{jamais} de subjonctif.
\item Interrogative partielle : mot interrogatif \textbf{accentué} + indicatif.
\item Impérative : infinitif (\esp{Le ordenó salir.}) ou \esp{que} + subjonctif imparfait (\esp{Ordenó que salieran.}).
\item Adapter personnes, possessifs (\esp{mi} $\to$ \esp{su}), démonstratifs (\esp{este} $\to$ \esp{aquel}) et marqueurs (\esp{hoy} $\to$ \esp{aquel día}, \esp{aquí} $\to$ \esp{allí}).
\end{itemize}
\end{aretenir}

\courssub{Entraînement}

\begin{exoresolu}[Transposition guidée]
Énoncé. Transposez au style indirect en commençant par \esp{El periodista dijo que...} / \esp{preguntó...} / \esp{ordenó...} :

\esp{«He visitado hoy el mercado de Mvog-Mbi —dijo el periodista—. Los comerciantes están contentos. ¿Vendréis mañana aquí otra vez? —preguntó a los clientes—. ¡No aparquéis delante de la entrada! —ordenó el policía.»}

\tcblower

Solution détaillée.
\begin{enumerate}
\item \esp{«He visitado hoy el mercado»} : déclarative $\to$ \esp{que} + pluscuamperfecto ; \esp{hoy} $\to$ \esp{aquel día}. \esp{El periodista dijo que había visitado aquel día el mercado de Mvog-Mbi.} (Le journaliste a dit qu'il avait visité ce jour-là le marché de Mvog-Mbi.)
\item \esp{«Los comerciantes están contentos»} : présent $\to$ imparfait. \esp{Añadió que los comerciantes estaban contentos.} (Il a ajouté que les commerçants étaient contents.)
\item \esp{«¿Vendréis mañana aquí otra vez?»} : question totale $\to$ \esp{si} + conditionnel (futur rapporté) ; \esp{mañana} $\to$ \esp{al día siguiente} ; \esp{aquí} $\to$ \esp{allí} ; \esp{venir} $\to$ \esp{volver} car le point de vue change. \esp{Preguntó a los clientes si volverían allí al día siguiente.} (Il a demandé aux clients s'ils reviendraient là le lendemain.)
\item \esp{«¡No aparquéis delante de la entrada!»} : ordre négatif $\to$ infinitif ou \esp{que} + subjonctif imparfait. \esp{El policía ordenó no aparcar delante de la entrada} ou \esp{ordenó que no aparcaran delante de la entrada.} (Le policier a ordonné de ne pas se garer devant l'entrée.)
\end{enumerate}
\end{exoresolu}

\begin{exercice}[À vous de transposer]
Transposez les phrases suivantes au style indirect (verbe introducteur au passé) :
\begin{enumerate}
\item \esp{«Estoy cansado», dijo Pedro.}
\item \esp{«¿Has visto mi mochila?», me preguntó Ana.}
\item \esp{«¿Dónde compraste este teléfono?», le pregunté.}
\item \esp{«Mañana iremos al estadio», prometieron los alumnos.}
\item \esp{«¡Abrid los libros!», ordenó la profesora.}
\item \esp{«No puedo venir hoy aquí», explicó mi hermano.}
\end{enumerate}
\end{exercice}

\begin{corrige}[Exercice « À vous de transposer »]
\begin{enumerate}
\item \esp{Pedro dijo que estaba cansado.} (Pedro a dit qu'il était fatigué.) Présent $\to$ imparfait ; \esp{yo} $\to$ \esp{él}.
\item \esp{Ana me preguntó si había visto su mochila.} (Ana m'a demandé si j'avais vu son sac.) Question totale : \esp{si} + pluscuamperfecto ; \esp{mi} $\to$ \esp{su}.
\item \esp{Le pregunté dónde había comprado ese/aquel teléfono.} (Je lui ai demandé où il avait acheté ce téléphone.) Accent conservé sur \esp{dónde} ; \esp{este} $\to$ \esp{ese/aquel}.
\item \esp{Los alumnos prometieron que irían al estadio al día siguiente.} (Les élèves ont promis qu'ils iraient au stade le lendemain.) Futur $\to$ conditionnel ; \esp{mañana} $\to$ \esp{al día siguiente}.
\item \esp{La profesora ordenó abrir los libros} / \esp{ordenó que abriéramos los libros.} (La professeure a ordonné d'ouvrir les livres.) Infinitif ou \esp{que} + subjonctif imparfait.
\item \esp{Mi hermano explicó que no podía ir allí aquel día.} (Mon frère a expliqué qu'il ne pouvait pas venir là ce jour-là.) \esp{hoy} $\to$ \esp{aquel día} ; \esp{aquí} $\to$ \esp{allí} ; \esp{venir} se rapporte comme \esp{ir} car le point de vue change.
\end{enumerate}
\end{corrige}

\begin{attention}[Les trois pièges du francophone]
\begin{enumerate}
\item \textbf{Subjonctif après \esp{si} interrogatif} : calque du doute français. On dit \esp{Preguntó si estaba cansado.} (Il a demandé s'il était fatigué), jamais \emph{*si estuviera}.
\item \textbf{Oubli de l'accent interrogatif} : \esp{Me preguntó cuándo salía el autobús.} et non \emph{*cuando}.
\item \textbf{Marqueurs non adaptés} : laisser \esp{hoy}, \esp{aquí}, \esp{mañana} dans un discours rapporté au passé est une faute de logique énonciative, en espagnol comme en français : dites \esp{aquel día}, \esp{allí}, \esp{al día siguiente}.
\end{enumerate}
\end{attention}

La maîtrise de ces mécanismes vous permet maintenant d'aborder la section suivante : les verbes introducteurs (\esp{decir}, \esp{preguntar}, \esp{afirmar}, \esp{prometer}) et leurs constructions propres, qui donnent au discours rapporté toute sa nuance.

\coursec{Les verbes introducteurs et leurs constructions}

Après avoir étudié la transposition du style direct au style indirect au passé (section 4), nous devons maintenant examiner de près le \textbf{moteur} de cette transposition : le \cle{verbe introducteur}. C'est lui qui dicte le choix du mode (indicatif ou subjonctif), la structure de la subordonnée (avec \esp{que}, \esp{si}, infinitif) et parfois même le temps du verbe rapporté. Maîtriser les constructions de chaque famille de verbes introducteurs, c'est s'assurer de ne jamais commettre d'erreur de concordance ni de structure lors de l'épreuve de thème/version du Baccalauréat.

\courssub{Qu'est-ce qu'un verbe introducteur ?}

\begin{definition}[Verbe introducteur]
Un \textbf{verbe introducteur} (\emph{verbo introductorio}) est un verbe de parole (\esp{decir, afirmar, preguntar}), de pensée (\esp{creer, pensar, opinar}), de volonté (\esp{querer, ordenar, pedir}) ou de sentiment (\esp{alegrarse, temer, dudar}) qui introduit une proposition subordonnée rapportant les paroles, les pensées ou les sentiments d'autrui. Il détermine :
\begin{itemize}
    \item Le \textbf{mode} de la subordonnée (indicatif ou subjonctif).
    \item La \textbf{conjonction} de liaison (\esp{que}, \esp{si}, préposition + infinitif).
    \item La \textbf{concordance des temps} si le verbe introducteur est au passé.
\end{itemize}
\end{definition}

\begin{exemplebox}[Exemples de verbes introducteurs au présent et au passé]
\esp{Juan dice: «Llegaré mañana.»} $\rightarrow$ \esp{Juan \textbf{dice} que \textbf{llegará} mañana.} (présent $\rightarrow$ pas de changement de temps) \\
\esp{Juan dijo: «Llegaré mañana.»} $\rightarrow$ \esp{Juan \textbf{dijo} que \textbf{llegaría} mañana.} (passé $\rightarrow$ futur $\rightarrow$ conditionnel)
\end{exemplebox}

\courssub{Les verbes de déclaration et d'affirmation : l'indicatif}

Les verbes qui expriment une \textbf{certitude}, une \textbf{information objective} ou une \textbf{déclaration factuelle} appellent l'\textbf{indicatif} dans la subordonnée introduite par \esp{que}. C'est le cas le plus fréquent en style indirect déclaratif.

\begin{propriete}[Verbes d'assertion $\rightarrow$ Indicatif]
\begin{tabular}{|l|l|l|}
\hline
\textbf{Verbe (infinitif)} & \textbf{Sens} & \textbf{Construction} \\ \hline
\esp{decir} & dire & \esp{decir que} + indicatif \\ \hline
\esp{afirmar} & affirmer & \esp{afirmar que} + indicatif \\ \hline
\esp{declarar} & déclarer & \esp{declarar que} + indicatif \\ \hline
\esp{asegurar} & assurer & \esp{asegurar que} + indicatif \\ \hline
\esp{anunciar} & annoncer & \esp{anunciar que} + indicatif \\ \hline
\esp{explicar} & expliquer & \esp{explicar que} + indicatif \\ \hline
\esp{contestar / responder} & répondre & \esp{contestar que} + indicatif \\ \hline
\esp{manifestar} & manifester & \esp{manifestar que} + indicatif \\ \hline
\esp{reconocer} & reconnaître & \esp{reconocer que} + indicatif \\ \hline
\esp{confirmar} & confirmer & \esp{confirmar que} + indicatif \\ \hline
\end{tabular}
\end{propriete}

\begin{exemplebox}[Déclarations au style indirect passé — Indicatif]
\esp{El ministro \textbf{afirmó} que la economía \textbf{crecía} un 3\%.} \\
« Le ministre a affirmé que l'économie croissait de 3\%. » (imparfait $\leftarrow$ présent) \\[4pt]
\esp{La portavoz \textbf{anunció} que el tratado \textbf{había sido} firmado.} \\
« La porte-parole a annoncé que le traité avait été signé. » (plus-que-parfait $\leftarrow$ passé composé) \\[4pt]
\esp{El testigo \textbf{declaró} que \textbf{no vio} al acusado.} \\
« Le témoin a déclaré qu'il n'avait pas vu l'accusé. » (prétérit $\leftarrow$ prétérit) \\[4pt]
\esp{Mi hermano \textbf{explicó} que \textbf{no podía} venir.} \\
« Mon frère a expliqué qu'il ne pouvait pas venir. » (imparfait $\leftarrow$ présent) \\[4pt]
\esp{El profesor \textbf{contestó} que el examen \textbf{sería} el viernes.} \\
« Le prof a répondu que l'examen serait vendredi. » (conditionnel $\leftarrow$ futur)
\end{exemplebox}

\begin{attention}[Piège francophone : \esp{decir} vs \esp{dire à quelqu'un}]
En français, \emph{dire} s'emploie souvent avec un COI : \emph{Il a dit \textbf{à Paul} qu'il viendrait}. En espagnol, \esp{decir} prend un COD (la chose dite) et un COI (le destinataire) \textbf{sans préposition} :
\begin{itemize}
    \item \esp{Le \textbf{dijo a Pablo} que vendría.} (correct)
    \item \esp{Le \textbf{dijo que} vendría.} (correct, COD = \esp{que}-proposition)
    \item \textbf{Jamais} : \esp{*Dijo a que vendría.}
\end{itemize}
De même, \esp{preguntar} ne prend pas de préposition devant la personne : \esp{Le pregunté a María...} (Je demandai \textbf{à} Marie...).
\end{attention}

\courssub{Les verbes de question : \esp{preguntar} et \esp{preguntarse}}

Les verbes interrogatifs introduisent une \textbf{question indirecte}. La liaison se fait par \esp{si} (question totale) ou par le \textbf{mot interrogatif} (question partielle) : \esp{qué, quién, cuándo, cómo, por qué, cuánto, cuál...}. Le mode reste l'\textbf{indicatif} si la question porte sur un fait réel ; le \textbf{subjonctif} apparaît si le verbe introducteur exprime un doute (ex. \esp{dudar si}, \esp{no saber si}).

\begin{propriete}[\esp{preguntar} et \esp{preguntarse} : constructions]
\begin{itemize}
    \item \textbf{Question totale (oui/non)} : \esp{preguntar / preguntarse + si} + indicatif (ou subjonctif si doute).
    \item \textbf{Question partielle} : \esp{preguntar / preguntarse + mot interrogatif} + indicatif (ou subjonctif si doute).
    \item \textbf{Sujet} : si le sujet de la question est différent de celui du verbe introducteur, il s'exprime explicitement en espagnol (pas d'inversion sujet/verbe comme en français).
\end{itemize}
\end{propriete}

\begin{exemplebox}[Questions indirectes — Style indirect passé]
\esp{Me \textbf{preguntó} \textbf{si} \textbf{quería} café.} \\
« Il m'a demandé \textbf{si} je voulais du café. » \\[4pt]
\esp{El periodista \textbf{preguntó} \textbf{cuándo} \textbf{llegaría} el presidente.} \\
« Le journaliste a demandé \textbf{quand} le président arriverait. » \\[4pt]
\esp{Nos \textbf{preguntamos} \textbf{por qué} \textbf{no habían} venido.} \\
« Nous nous sommes demandé \textbf{pourquoi} ils n'étaient pas venus. » \\[4pt]
\esp{El turista \textbf{preguntó} \textbf{dónde} \textbf{estaba} el mercado.} \\
« Le touriste a demandé \textbf{où} était le marché. »
\end{exemplebox}

\begin{aretenir}[Accent sur les mots interrogatifs en question indirecte]
En espagnol, les mots interrogatifs gardent \textbf{toujours leur accent graphique} en question indirecte : \esp{qué, quién, cómo, cuándo, dónde, por qué, cuánto, cuál}. C'est une différence majeure avec le français (\emph{où, quand, pourquoi, combien} sans accent).
\end{aretenir}

\courssub{Les verbes de volonté, d'ordre et de promesse : subjonctif ou infinitif}

Cette famille est cruciale pour le Baccalauréat. Les verbes exprimant une \textbf{volonté}, un \textbf{ordre}, une \textbf{permission}, une \textbf{interdiction} ou un \textbf{engagement} (promesse, serment) exigent le \textbf{subjonctif} dans la subordonnée \esp{que} \textbf{si les deux sujets sont différents}. Si le sujet est le même, on utilise l'\textbf{infinitif} (souvent précédé d'une préposition).

\begin{propriete}[Verbes de volonté/ordre/promesse : alternance Subjonctif / Infinitif]
\begin{tabular}{|l|l|l|l|}
\hline
\textbf{Verbe} & \textbf{Sens} & \textbf{Sujets différents $\rightarrow$ Subjonctif} & \textbf{Même sujet $\rightarrow$ Infinitif} \\ \hline
\esp{querer} & vouloir & \esp{Quiero que vengas.} & \esp{Quiero venir.} \\ \hline
\esp{desear} & désirer & \esp{Deseo que ganes.} & \esp{Deseo ganar.} \\ \hline
\esp{mandar} & ordonner & \esp{Mandó que salieran.} & \esp{Mandó salir.} \\ \hline
\esp{ordenar} & ordonner & \esp{Ordenó que callaran.} & \esp{Ordenó callar.} \\ \hline
\esp{pedir} & demander & \esp{Pedí que me ayudaran.} & \esp{Pedí ayudarme.} \\ \hline
\esp{prohibir} & interdire & \esp{Prohibió que entráramos.} & \esp{Prohibió entrar.} \\ \hline
\esp{permitir} & permettre & \esp{Permitió que entráramos.} & \esp{Permitió entrar.} \\ \hline
\esp{aconsejar} & conseiller & \esp{Aconsejó que estudiara.} & \esp{Aconsejó estudiar.} \\ \hline
\esp{prometer} & promettre & \esp{Prometió que lo haría.} & \esp{Prometió hacerlo.} \\ \hline
\esp{jurar} & jurer & \esp{Juró que diría la verdad.} & \esp{Juró decir la verdad.} \\ \hline
\end{tabular}
\end{propriete}

\begin{exemplebox}[Volonté, ordre, promesse — Style indirect passé]
\esp{El capitán \textbf{ordenó} que los soldados \textbf{avanzaran}.} \\
« Le capitaine a ordonné que les soldats avancent. » (subjonctif imparfait : sujets différents) \\[4pt]
\esp{Mi padre me \textbf{pidió} que \textbf{llamara} a la policía.} \\
« Mon père m'a demandé d'appeler la police. » (subjonctif imparfait) \\[4pt]
\esp{El director \textbf{prohibió} \textbf{fumar} en el patio.} \\
« Le directeur a interdit de fumer dans la cour. » (infinitif : sujet unique = règle générale) \\[4pt]
\esp{Mi novio me \textbf{prometió} que \textbf{me llamaría} cada día.} \\
« Mon petit ami m'a promis qu'il m'appellerait chaque jour. » (conditionnel dans \esp{que}-proposition car \esp{prometer} + futur/conditionnel) \\[4pt]
\esp{La jefa \textbf{aconsejó} \textbf{estudiar} más.} \\
« La patronne a conseillé d'étudier plus. » (infinitif : même sujet implicite)
\end{exemplebox}

\begin{attention}[Le piège \esp{pedir} — Jamais \esp{*pedir de}]
C'est l'erreur n°1 des francophones par calque sur \emph{demander \textbf{de}}.
\begin{itemize}
    \item \textbf{Correct} : \esp{Pedí \textbf{que} me ayudaran.} (subjonctif, sujets différents)
    \item \textbf{Correct} : \esp{Pedí \textbf{ayudarme}.} (infinitif, même sujet)
    \item \textbf{INCORRECT} : \esp{*Pedí de que me ayudaran / *Pedí de ayudarme.}
\end{itemize}
De même : \esp{ordenar que / ordenar + infinitif} (jamais \esp{*ordenar de}), \esp{prohibir que / prohibir + infinitif}.
\end{attention}

\courssub{Les verbes de sentiment et de doute : le subjonctif}

Les verbes exprimant un \textbf{sentiment subjectif} (joie, tristesse, peur, regret, surprise) ou un \textbf{doute / une négation de certitude} déclenchent systématiquement le \textbf{subjonctif} dans la subordonnée \esp{que}, quel que soit le temps du verbe introducteur (présent ou passé).

\begin{propriete}[Verbes de sentiment et de doute $\rightarrow$ Subjonctif]
\begin{tabular}{|l|l|l|}
\hline
\textbf{Famille} & \textbf{Verbes principaux} & \textbf{Construction} \\ \hline
\multicolumn{3}{|c|}{\textbf{Sentiment / Émotion}} \\ \hline
Joie & \esp{alegrarse de que}, \esp{gustar que} & + subjonctif \\ \hline
Tristesse & \esp{sentir que}, \esp{lamentar que} & + subjonctif \\ \hline
Peur / Inquiétude & \esp{temer que}, \esp{preocuparse de que} & + subjonctif \\ \hline
Surprise & \esp{sorprender que}, \esp{extrañar que} & + subjonctif \\ \hline
Regret & \esp{arrepentirse de que} & + subjonctif \\ \hline
\multicolumn{3}{|c|}{\textbf{Doute / Négation de certitude}} \\ \hline
Doute & \esp{dudar que}, \esp{no creer que}, \esp{no estar seguro de que} & + subjonctif \\ \hline
Négation & \esp{negar que}, \esp{no pensar que}, \esp{no suponer que} & + subjonctif \\ \hline
Improbable & \esp{es improbable que}, \esp{es dudoso que} & + subjonctif \\ \hline
\end{tabular}
\end{propriete}

\begin{exemplebox}[Sentiment et doute — Style indirect passé]
\esp{Me \textbf{alegré} de que \textbf{hubieras} aprobado.} \\
« Je me suis réjoui que tu aies réussi. » (subjonctif plus-que-parfait : antériorité) \\[4pt]
\esp{El entrenador \textbf{temía} que el equipo \textbf{perdiera}.} \\
« L'entraîneur craignait que l'équipe perde. » (subjonctif imparfait : simultanéité) \\[4pt]
\esp{No \textbf{creí} que \textbf{fuera} verdad.} \\
« Je n'ai pas cru que ce fût vrai. » (subjonctif imparfait) \\[4pt]
\esp{Los padres \textbf{sintieron} que sus hijos \textbf{se fueran}.} \\
« Les parents ont regretté que leurs enfants partent. » (subjonctif imparfait) \\[4pt]
\esp{Dudábamos que el proyecto \textbf{saliera} adelante.} \\
« Nous doutions que le projet aboutisse. » (subjonctif imparfait)
\end{exemplebox}

\courssub{Tableau récapitulatif des quatre grandes familles}

\begin{center}
\begin{tabularx}{\linewidth}{|X|X|X|X|}
\hline
\rowcolor{popblueL}
\textbf{Famille} & \textbf{Verbes types} & \textbf{Mode de la subordonnée} & \textbf{Liaison} \\ \hline
\textbf{1. Déclaration / Affirmation / Connaissance} & \esp{decir, afirmar, declarar, asegurar, anunciar, explicar, contestar, responder, manifestar, reconocer, confirmar, creer, pensar, opinar, saber} & \textbf{INDICATIF} & \esp{que} \\ \hline
\textbf{2. Question} & \esp{preguntar, preguntarse, indagar, investigar, querer saber} & \textbf{INDICATIF} (fait réel) \newline \textbf{SUBJONCTIF} (doute) & \esp{si} / mot interrogatif (\esp{qué, cuándo...}) \\ \hline
\textbf{3. Volonté / Ordre / Permission / Promesse} & \esp{querer, desear, mandar, ordenar, pedir, prohibir, permitir, aconsejar, recomendar, prometer, jurar, insistir en que, sugerir que} & \textbf{SUBJONCTIF} (sujets différents) \newline \textbf{INFINITIF} (même sujet) & \esp{que} / infinitif (+ prép. selon verbe) \\ \hline
\textbf{4. Sentiment / Doute / Émotion} & \esp{alegrarse de que, sentir que, temer que, dudar que, no creer que, extrañar que, sorprender que, lamentar que, es posible/improbable que} & \textbf{SUBJONCTIF} & \esp{que} / \esp{de que} \\ \hline
\end{tabularx}
\end{center}

\begin{popfigure}
\begin{tikzpicture}[
    mindmap,
    concept color=popblue,
    level 1/.style={level distance=4.5cm, sibling angle=90},
    level 2/.style={level distance=3.2cm, sibling angle=45},
    every node/.style={concept, execute at begin node=\small},
    text=white,
    font=\small
]
\node[concept] {VERBOS INTRODUCTORES}
    child[concept color=popgreen] {
        node {DECLARACIÓN}
        child { node {afirmar} }
        child { node {declarar} }
        child { node {asegurar} }
        child { node {anunciar} }
        child { node {explicar} }
        child { node {contestar} }
        child { node {creer} }
        child { node {pensar} }
        edge from parent[draw=popgreen!60, thick]
    }
    child[concept color=poporange] {
        node {PREGUNTA}
        child { node {preguntar} }
        child { node {preguntarse} }
        child { node {indagar} }
        child { node {querer saber} }
        edge from parent[draw=poporange!60, thick]
    }
    child[concept color=poppurple] {
        node {VOLUNTAD / ORDEN / PROMESA}
        child { node {querer} }
        child { node {mandar} }
        child { node {ordenar} }
        child { node {pedir} }
        child { node {prohibir} }
        child { node {permitir} }
        child { node {aconsejar} }
        child { node {prometer} }
        child { node {jurar} }
        edge from parent[draw=poppurple!60, thick]
    }
    child[concept color=poppink] {
        node {SENTIMIENTO / DUDA}
        child { node {alegrarse} }
        child { node {sentir} }
        child { node {temer} }
        child { node {dudar} }
        child { node {no creer} }
        child { node {extrañar} }
        child { node {sorprender} }
        child { node {lamentar} }
        edge from parent[draw=poppink!60, thick]
    };
\node[below=0.2cm of current bounding box.south, font=\small\bfseries, text=popdark] {Famille $\rightarrow$ Mode : Déclaration/Question = Indicatif ; Volonté/Sentiment/Doute = Subjonctif (sujets diff.)};
\end{tikzpicture}
\legende{Carte mentale : les quatre familles de verbes introducteurs et le mode verbal associé.}
\end{popfigure}

\courssub{Le cas particulier de \esp{decir} : déclaration ou ordre}

\esp{Decir} est le verbe introducteur le plus fréquent et le plus \textbf{polysémique}. Sa construction change radicalement selon qu'il introduit une \textbf{information} (déclaration) ou un \textbf{ordre/une demande} (volonté).

\begin{propriete}[\esp{decir} : deux constructions, deux sens]
\begin{tabular}{|l|l|l|l|}
\hline
\textbf{Sens} & \textbf{Exemple style direct} & \textbf{Style indirect} & \textbf{Mode} \\ \hline
\textbf{Déclaration} & \esp{«Llegaré tarde», dijo Juan.} & \esp{Juan \textbf{dijo que llegaría} tarde.} & \textbf{Indicatif} \\ \hline
\textbf{Ordre / Demande} & \esp{«Llegad temprano», dijo el jefe.} & \esp{El jefe \textbf{dijo que llegaran} temprano.} & \textbf{Subjonctif} \\ \hline
\textbf{Ordre (infinitif)} & \esp{«Llegad temprano», les dijo.} & \esp{Les \textbf{dijo que llegaran} temprano.} & \textbf{Subjonctif} \\ \hline
\textbf{Conseil} & \esp{«Estudia más», me dijo.} & \esp{Me \textbf{dijo que estudiara} más.} & \textbf{Subjonctif} \\ \hline
\end{tabular}
\end{propriete}

\begin{attention}[Comment distinguer déclaration et ordre avec \esp{decir} ?]
Regardez le \textbf{contenu} de la phrase rapportée :
\begin{itemize}
    \item Si c'est une \textbf{information factuelle} (description, narration, opinion) $\rightarrow$ \esp{decir que} + \textbf{Indicatif}.
    \item Si c'est une \textbf{injction} (ordre, conseil, demande, prière, interdiction) adressée à quelqu'un $\rightarrow$ \esp{decir que} + \textbf{Subjonctif} (sujets différents) ou \esp{decir + infinitif} (même sujet).
\end{itemize}
\textbf{Astuce} : Remplacez mentalement \esp{dijo} par \esp{ordenó} ou \esp{aconsejó}. Si la phrase a du sens, c'est un ordre/conseil $\rightarrow$ Subjonctif.
\end{attention}

\courssub{Texte d'entraînement : Une conférence de presse à Yaoundé}

\begin{textebox}[Conférence de presse — Ministère de la Communication, Yaoundé]
El Ministro de Comunicación, René Emmanuel Sadi, \textbf{compareció} ayer ante la prensa para \textbf{informar} sobre las nuevas medidas de regulación de los medios digitales. \textbf{Declaró} que el gobierno \textbf{garantiza} la libertad de expresión, pero \textbf{advirtió} que no \textbf{tolerará} la propagación de «noticias falsas que \textbf{ponen en peligro} la paz social».

Ante las preguntas de los periodistas, el ministro \textbf{respondió} que el proyecto de ley \textbf{estaba} en fase de consulta y que \textbf{se presentarían} enmiendas antes de su adopción. \textbf{Aseguró} que \textbf{no había} intención de censurar, sino de «proteger a los ciudadanos de la desinformación».

Una reportera le \textbf{preguntó} \textbf{si} el texto \textbf{preveía} sanciones penales. El ministro \textbf{contestó} que \textbf{sí}, que \textbf{habría} multas y penas de prisión para los reincidentes. Otra periodista quiso \textbf{saber} \textbf{cómo} \textbf{se definiría} el concepto de «noticia falsa». Sadi \textbf{explicó} que un comité técnico \textbf{lo determinaría} caso por caso.

La oposición \textbf{ha criticado} el proyecto. Su portavoz \textbf{afirmó} que la ley \textbf{es} una «herramienta de control político» y \textbf{pidió} que \textbf{se retirara} inmediatamente. \textbf{Advirtió} que \textbf{movilizarían} a la ciudadanía si el Parlamento \textbf{aprobaba} el texto sin consenso. El Primer Ministro, por su parte, \textbf{prometió} que \textbf{se abriría} un debate nacional inclusivo.
\end{textebox}

\begin{center}
\begin{tabularx}{\linewidth}{|X|X|}
\hline
\rowcolor{popblueL} \textbf{Espagnol} & \textbf{Français} \\ \hline
\esp{compareció} & il a comparu / s'est présenté (prétérit de \esp{comparecer}) \\ \hline
\esp{informar} & informer \\ \hline
\esp{advirtió} & il a averti / mis en garde (prétérit de \esp{advertir}) \\ \hline
\esp{ponen en peligro} & mettent en danger (présent indicatif, fait actuel) \\ \hline
\esp{estaba} & était (imparfait indicatif) \\ \hline
\esp{se presentarían} & seraient présentées (conditionnel $\leftarrow$ futur) \\ \hline
\esp{no había} & il n'y avait pas (imparfait de \esp{hay}) \\ \hline
\esp{preveía} & prévoyait (imparfait) \\ \hline
\esp{habría} & il y aurait (conditionnel de \esp{hay}) \\ \hline
\esp{quiso saber} & a voulu savoir (prétérit de \esp{querer} + infinitif) \\ \hline
\esp{se definiría} & se définirait (conditionnel passif réfléchi) \\ \hline
\esp{lo determinaría} & le déterminerait (conditionnel) \\ \hline
\esp{ha criticado} & a critiqué (passé composé) \\ \hline
\esp{herramienta de control político} & outil de contrôle politique \\ \hline
\esp{pidió que se retirara} & a demandé qu'il soit retiré (subjonctif imparfait \esp{retirara} $\leftarrow$ \esp{pedir que} + passé) \\ \hline
\esp{movilizarían} & ils mobiliseraient (conditionnel) \\ \hline
\esp{aprobaba} & approuvait (imparfait \textbf{indicatif} : \esp{si} + indicatif pour hypothèse réelle) \\ \hline
\esp{prometió que se abriría} & a promis qu'il s'ouvrirait (conditionnel) \\ \hline
\end{tabularx}
\end{center}

\textbf{Questions de compréhension et d'analyse grammaticale :}
\begin{enumerate}
    \item Relevez dans le texte \textbf{6 verbes introducteurs} différents et, pour chacun, indiquez : l'infinitif, le temps employé, la famille (déclaration, question, volonté, sentiment/doute), le mode de la subordonnée introduite.
    \item Pourquoi \esp{advirtió que no tolerará} utilise-t-il le futur de l'indicatif (\esp{tolerará}) et non le subjonctif ? \textit{Indice : nature de l'avertissement (constat de fait futur).}
    \item Analysez la structure de \esp{pidió que se retirara}. Pourquoi \esp{retirara} est-il au subjonctif imparfait ?
    \item Traduisez en français les deux phrases contenant \esp{advirtió} (ministro) et \esp{advirtió} (porte-parole opposition). Notez la différence de construction.
    \item Transformez la déclaration du ministre \esp{«El proyecto de ley está en fase de consulta»} en style indirect au passé avec le verbe introducteur \esp{manifestó}.
\end{enumerate}

\courssub{Exercice résolu : Thème / Version sur les verbes introducteurs}

\begin{exoresolu}[Thème : Transposition au style indirect passé]
\textbf{Énoncé :} Transformez les phrases suivantes du style direct au style indirect en utilisant le verbe introducteur indiqué au prétérit. Respectez la concordance des temps et le mode requis.

1. \esp{«Voy a estudiar medicina», dijo Ana.} (verbe : \esp{afirmar}) \\
2. \esp{«¿Vendrás a la fiesta?», me preguntó Luis.} (verbe : \esp{preguntar}) \\
3. \esp{«Termina el trabajo antes del viernes», ordenó el jefe.} (verbe : \esp{ordenar}) \\
4. \esp{«Te ayudaré con los deberes», me prometió mi hermano.} (verbe : \esp{prometer}) \\
5. \esp{«No creo que sea buena idea», opinó Pedro.} (verbe : \esp{dudar}) \\
6. \esp{«¡Qué bien que hayas venido!», me dijo mi abuela.} (verbe : \esp{alegrarse de}) \\

\tcblower
\textbf{Solution détaillée :}

1. \esp{Ana \textbf{afirmó} que \textbf{iba a estudiar} medicina.} \\
   \textit{Analyse :} \esp{afirmar} = déclaration $\rightarrow$ indicatif. Futur périphrastique \esp{voy a estudiar} $\rightarrow$ \esp{iba a estudiar} (imparfait de \esp{ir} + infinitif).

2. \esp{Luis me \textbf{preguntó} \textbf{si} \textbf{vendría} a la fiesta.} \\
   \textit{Analyse :} Question totale $\rightarrow$ \esp{si}. \esp{preguntar} = question fait réel $\rightarrow$ indicatif. Futur \esp{vendrás} $\rightarrow$ conditionnel \esp{vendría}. COI \esp{me} placé avant le verbe conjugué.

3. \esp{El jefe \textbf{ordenó} que \textbf{terminara} el trabajo antes del viernes.} \\
   \textit{Analyse :} \esp{ordenar} = ordre $\rightarrow$ subjonctif (sujets différents : jefe / él/ella impliqué dans \esp{termina}). Prétérit \esp{termina} $\rightarrow$ subjonctif imparfait \esp{terminara}.

4. \esp{Mi hermano me \textbf{prometió} que \textbf{me ayudaría} con los deberes.} \\
   \textit{Analyse :} \esp{prometer} = promesse $\rightarrow$ \esp{que} + conditionnel (ou infinitif \esp{ayudarme} si même sujet, mais ici sujets différents : hermano / yo). Futur \esp{ayudaré} $\rightarrow$ conditionnel \esp{ayudaría}.

5. \esp{Pedro \textbf{dudó} que \textbf{fuera} buena idea.} \\
   \textit{Analyse :} \esp{dudar} = doute $\rightarrow$ subjonctif. Présent \esp{es} $\rightarrow$ subjonctif imparfait \esp{fuera}. \textit{Attention :} \esp{no creo} $\rightarrow$ subjonctif, mais le verbe introducteur ici est \esp{dudó} (passé).

6. \esp{Mi abuela se \textbf{alegró} de que \textbf{hubiera} venido.} \\
   \textit{Analyse :} \esp{alegrarse de que} = sentiment/joie $\rightarrow$ subjonctif. Antériorité (\esp{hayas venido} = passé composé) $\rightarrow$ subjonctif plus-que-parfait \esp{hubiera venido}.
\end{exoresolu}

\courssub{Exercice d'application}

\begin{exercice}[Entraînement BAC — Thème et Version]
\textbf{A. Version (Espagnol $\rightarrow$ Français) :} Traduisez en français les phrases suivantes en soignant la concordance des temps.
\begin{enumerate}
    \item \esp{El presidente anunció que firmaría el decreto mañana.}
    \item \esp{Nos preguntamos por qué no habían llegado aún.}
    \item \esp{El profesor prohibió usar el móvil en clase.}
    \item \esp{Mi madre me pidió que le llamara al llegar.}
    \item \esp{Temían que el río se desbordara.}
    \item \esp{El testigo negó que hubiera visto al acusado.}
\end{enumerate}

\textbf{B. Thème (Français $\rightarrow$ Espagnol) :} Traduisez en espagnol en choisissant le bon verbe introducteur et le bon mode.
\begin{enumerate}
    \item Le ministre a affirmé que le taux de chômage avait baissé. (\esp{afirmar})
    \item Le journaliste a demandé quand commencerait la réunion. (\esp{preguntar})
    \item Le capitaine a ordonné aux soldats de ne pas tirer. (\esp{ordenar} + subjonctif)
    \item Mon ami m'a promis qu'il viendrait me voir. (\esp{prometer} + conditionnel)
    \item Nous doutions qu'il dise la vérité. (\esp{dudar} + subjonctif imparfait)
    \item Elle s'est réjouie que nous ayons réussi. (\esp{alegrarse de} + subjonctif plus-que-parfait)
\end{enumerate}
\end{exercice}

\begin{corrige}[Exercice d'application]
\textbf{A. Version :}
1. Le président a annoncé qu'il signerait le décret demain. \\
2. Nous nous demandions pourquoi ils n'étaient pas encore arrivés. \\
3. Le professeur a interdit d'utiliser le portable en classe. (infinitif : règle générale, sujet unique) \\
4. Ma mère m'a demandé de l'appeler à mon arrivée. (subjonctif imparfait \esp{llamara} : antériorité / simultanéité future dans le passé) \\
5. Ils craignaient que le fleuve ne déborde. (subjonctif imparfait) \\
6. Le témoin a nié avoir vu l'accusé. (subjonctif plus-que-parfait \esp{hubiera visto} $\rightarrow$ antériorité)

\textbf{B. Thème :}
1. \esp{El ministro afirmó que la tasa de paro había bajado.} (plus-que-parfait indicatif : antériorité) \\
2. \esp{El periodista preguntó cuándo comenzaría la reunión.} (conditionnel : futur dans le passé) \\
3. \esp{El capitán ordenó que los soldados no dispararan.} (subjonctif imparfait : ordre, sujets différents) \\
4. \esp{Mi amigo me prometió que vendría a verme.} (conditionnel : promesse, futur dans le passé) \\
5. \esp{Dudábamos que dijera la verdad.} (subjonctif imparfait \esp{dijera} : doute au passé) \\
6. \esp{Se alegró de que hubiéramos tenido éxito.} (subjonctif plus-que-parfait : joie pour fait antérieur)
\end{corrige}

\courssub{Erreurs fréquentes des francophones : synthèse des pièges}

\begin{attention}[Top 5 des erreurs à ne plus commettre]
\begin{enumerate}
    \item \textbf{\esp{*Pedí de que...} / \esp{*Ordené de que...}} $\rightarrow$ \textbf{Jamais de préposition \esp{de}} après \esp{pedir, ordenar, mandar, prohibir, permitir, aconsejar} devant \esp{que} + subjonctif. \textit{Correct :} \esp{Pedí que viniera.}
    \item \textbf{\esp{*Dijo a que...}} $\rightarrow$ \esp{Decir} n'introduit pas \esp{que} par \esp{a}. \textit{Correct :} \esp{Le dijo que...} (COI \esp{le} = lui) ou \esp{Dijo a Juan que...}
    \item \textbf{Oubli de l'accent sur les interrogatifs indirects} : \esp{*pregunto donde esta} $\rightarrow$ \textbf{\esp{pregunto \textbf{dónde} está}}. L'accent marque la question indirecte.
    \item \textbf{Indicatif au lieu de subjonctif après volonté/sentiment} : \esp{*Quiero que vienes} $\rightarrow$ \textbf{\esp{Quiero que \textbf{vengas}}}. \esp{*Me alegro que estás aquí} $\rightarrow$ \textbf{\esp{Me alegro de que \textbf{estés} aquí}}.
    \item \textbf{Mauvaise concordance avec \esp{prometer/jurar}} : \esp{*Prometió que lo hace} $\rightarrow$ \textbf{\esp{Prometió que lo \textbf{haría}}} (conditionnel) ou \textbf{\esp{Prometió \textbf{hacerlo}}} (infinitif).
\end{enumerate}
\end{attention}

\begin{methode}[Comment choisir le bon mode en 3 secondes — Méthode BAC]
\begin{enumerate}
    \item \textbf{Identifiez le verbe introducteur} et sa \textbf{famille} (Déclaration ? Question ? Volonté/Ordre ? Sentiment/Doute ?).
    \item \textbf{Vérifiez les sujets} : \textit{Même sujet} $\rightarrow$ \textbf{Infinitif} (souvent avec préposition). \textit{Sujets différents} $\rightarrow$ \textbf{Subjonctif} (sauf Déclaration/Question fait réel = Indicatif).
    \item \textbf{Appliquez la concordance} si le verbe introducteur est au passé : Présent/Futur $\rightarrow$ Imparfait/Conditionnel ; Passé composé/Prétérit $\rightarrow$ Plus-que-parfait subjonctif.
\end{enumerate}
\end{methode}

\begin{savaistu}[Le \esp{que} « expletif » ou « neutre » en espagnol ?]
Contrairement au français où \emph{que} est obligatoire pour subordonner, l'espagnol \textbf{ne tolère pas} le \esp{que} dit « expletif » après les verbes de crainte (\esp{temer, tener miedo}) en style \textbf{positif}.
\begin{itemize}
    \item \esp{Temo \textbf{que} no venga.} (négatif $\rightarrow$ \esp{que} + subjonctif = je crains qu'il ne vienne pas)
    \item \esp{Temo \textbf{que} venga.} (positif $\rightarrow$ \esp{que} + subjonctif = je crains qu'il vienne) \textbf{Correct en espagnol standard.}
    \item \textit{Note :} En français, \emph{je crains qu'il vienne} (subjonctif expletif \emph{ne} \emph{vienne}) ; en espagnol, pas de \esp{no} expletif, le subjonctif seul porte la nuance de crainte.
\end{itemize}
Autre particularité : \esp{evitar que} + subjonctif (empêcher que), \esp{impedir que} + subjonctif.
\end{savaistu}

\begin{experience}[Activité orale : « La conférence de presse improvisée »]
\textbf{Objectif :} Produire du style indirect oralement en chaîne.
\textbf{Déroulement :}
\begin{enumerate}
    \item L'enseignant donne un thème d'actualité camerounaise (ex. : « La CAN 2025 », « La réforme du bac », « Le prix du cacao »).
    \item Un élève (le « ministre ») fait une déclaration courte en style direct : \esp{«Vamos a construir tres estadios nuevos.»}
    \item L'élève suivant (le « journaliste ») rapporte au style indirect avec un verbe introducteur imposé par le prof : \esp{El ministro \textbf{anunció} que \textbf{iban a construir} tres estadios nuevos.}
    \item Un troisième élève (l'« opposant ») réagit avec un verbe de volonté/doute : \esp{La oposición \textbf{duda} que \textbf{se construyan}.}
    \item On continue la chaîne en changeant de verbe introducteur à chaque tour (\esp{afirmar, preguntar, ordenar, prometer, temer, alegrarse}...).
\end{enumerate}
\textbf{Variante écrite :} Chaque élève écrit une phrase en style direct sur un papier, les papiers circulent, chacun transforme en style indirect avec un verbe tiré au sort.
\end{experience}

\coursec{Comparaison français/espagnol, pièges et entraînement}

Nous savons désormais choisir le bon verbe introducteur (\esp{decir}, \esp{preguntar}, \esp{afirmar}, \esp{prometer}) et le construire correctement. Il reste à franchir le dernier obstacle : celui qui fait échouer le plus grand nombre de candidats francophones au baccalauréat. Dans cette section, nous allons comparer systématiquement le français et l'espagnol, repérer les pièges classiques, puis nous entraîner en version et en thème.

\courssub{Points communs entre le français et l'espagnol}

Bonne nouvelle : les deux langues fonctionnent de manière très proche. Quand le verbe introducteur est au passé, les deux langues opèrent un \cle{recul des temps} et transforment les \cle{marqueurs de temps et de lieu}.

\begin{propriete}[Le recul des temps : une logique commune]
Au style indirect passé, le français et l'espagnol reculent les temps de la même façon :
\begin{itemize}
\item présent → imparfait : \esp{« Vengo » → Dijo que venía.} (« Il a dit qu'il vient / qu'il venait ») ;
\item passé composé / prétérit → plus-que-parfait : \esp{« He visto el accidente » → Dijo que había visto el accidente.} (« Il a dit qu'il a vu / qu'il avait vu l'accident ») ;
\item futur → conditionnel présent : \esp{« Vendré » → Dijo que vendría.} (« Il a dit qu'il viendra / qu'il viendrait »).
\end{itemize}
La correspondance futur → conditionnel est \textbf{directe} entre les deux langues : \esp{vendría} = « viendrait ». C'est un point d'appui précieux pour les francophones.
\end{propriete}

\begin{propriete}[La transformation des marqueurs]
Dans les deux langues, les mots qui situent le discours dans le « ici et maintenant » du locuteur doivent être reculés :
\end{propriete}

\begin{center}
\begin{tabularx}{\linewidth}{|X|X|X|}
\hline
\rowcolor{popblueL} Style direct & Español (indirecto) & Français (indirect) \\
\hline
aujourd'hui & \esp{hoy → aquel día} & ce jour-là \\
\hline
demain & \esp{mañana → al día siguiente} & le lendemain \\
\hline
hier & \esp{ayer → el día anterior} & la veille \\
\hline
maintenant & \esp{ahora → entonces} & alors / à ce moment-là \\
\hline
ici & \esp{aquí → allí} & là / là-bas \\
\hline
ce / cet & \esp{este → ese / aquel} & ce...-là \\
\hline
il y a deux jours & \esp{hace dos días → dos días antes} & deux jours plus tôt \\
\hline
\end{tabularx}
\end{center}

\begin{exemplebox}[Observation]
\esp{El periodista dijo: « Publicaré el artículo mañana aquí ».} → \esp{El periodista dijo que publicaría el artículo al día siguiente allí.} (Le journaliste a dit qu'il publierait l'article le lendemain là-bas.)
\end{exemplebox}

\courssub{Différences majeures et pièges}

\textbf{1. Le « si » interrogatif indirect}\par
En français comme en espagnol, le \esp{si} de l'interrogation indirecte ne supporte \textbf{jamais le futur}. En français, le conditionnel est possible (concordance du futur direct) : « il m'a demandé si je \textbf{viendrais} ». En espagnol, après \esp{si} interrogatif introduit par un verbe au passé, on met l'indicatif (imparfait ou plus-que-parfait) ou le conditionnel (concordance du futur), \textbf{jamais le futur} :

\esp{Me preguntó si venía.} (Il m'a demandé si je venais / si je viendrais -- imparfait indicatif.)
\esp{Me preguntó si vendría.} (Il m'a demandé si je viendrais -- conditionnel, concordance du futur.)
Jamais : *\esp{Me preguntó si vendré.}

\textbf{2. Le subjonctif : caché en français, visible en espagnol}\par
C'est la différence la plus redoutable. Le français « cache » souvent le subjonctif derrière une forme identique à l'indicatif ou derrière un infinitif ; l'espagnol, lui, marque \textbf{toujours} le subjonctif, et au passé il exige l'\cle{imparfait du subjonctif} (\esp{viniera}, \esp{hablara}, \esp{tuviera}).

\begin{center}
\begin{tabularx}{\linewidth}{|X|X|X|}
\hline
\rowcolor{popblueL} Français & Español & Remarque \\
\hline
Il a dit qu'il viendrait. & \esp{Dijo que vendría.} & Conditionnel dans les deux langues (déclaration) \\
\hline
Il m'a demandé si je venais. & \esp{Me preguntó si venía.} & Imparfait de l'indicatif (question indirecte) \\
\hline
Il m'a dit de venir. & \esp{Me dijo que viniera.} & \textbf{Subjonctif imparfait obligatoire} (ordre/volonté) \\
\hline
Elle voulait que je l'aide. & \esp{Quería que la ayudara.} & Le français montre le subjonctif ici \\
\hline
Je ne croyais pas qu'il soit venu. & \esp{No creía que hubiera venido.} & Subjonctif plus-que-parfait \\
\hline
\end{tabularx}
\end{center}

\begin{attention}[Erreur n°1 des candidats francophones : le futur non reculé]
Sous la pression de l'examen, beaucoup d'élèves traduisent « il a dit qu'il viendrait » par *\esp{dijo que vendrá} ✗. C'est FAUX. Le verbe introducteur étant au passé (\esp{dijo}), le futur du discours direct doit reculer au \textbf{conditionnel présent} : \esp{dijo que vendría} ✓. Réflexe à acquérir : verbe introducteur au passé = jamais de futur dans la subordonnée.
\end{attention}

\begin{attention}[Le piège de « demain » rapporté]
Au style indirect passé, « demain » ne se traduit JAMAIS par \esp{mañana} : on dit \esp{al día siguiente}. De même « hier » devient \esp{el día anterior} et « aujourd'hui » \esp{aquel día}.
\esp{Dijo que vendría al día siguiente.} (Il a dit qu'il viendrait le lendemain.) — Jamais : *\esp{dijo que vendría mañana}.
\end{attention}

\begin{attention}[Le piège de l'infinitif français]
« Il m'a dit de venir » ne se traduit PAS par *\esp{me dijo de venir} ✗. En espagnol, \esp{decir} ne se construit pas avec l'infinitif pour rapporter un ordre ou une volonté : il faut une subordonnée au \textbf{subjonctif imparfait} introduite par \esp{que} : \esp{me dijo que viniera} ✓. Même chose avec \esp{pedir} : « elle m'a demandé de l'aider » → \esp{me pidió que la ayudara}.
\end{attention}

\begin{popfigure}
\begin{tikzpicture}[node distance=0.6cm and 1.6cm,
  fr/.style={rectangle, rounded corners, draw=popblue, fill=popblueL, align=center, text width=5.2cm, inner sep=5pt},
  es/.style={rectangle, rounded corners, draw=popgreen, fill=popgreenL, align=center, text width=5.2cm, inner sep=5pt},
  fleche/.style={-{Stealth[length=3mm]}, thick, poporange}]
\node[fr] (f1) {Il a dit qu'il \textbf{viendrait}.};
\node[es, right=of f1] (e1) {\esp{Dijo que \textbf{vendría}.}};
\draw[fleche] (f1) -- (e1);
\node[fr, below=of f1] (f2) {Il m'a demandé si je \textbf{venais}.};
\node[es, right=of f2] (e2) {\esp{Me preguntó si \textbf{venía}.}};
\draw[fleche] (f2) -- (e2);
\node[fr, below=of f2] (f3) {Il m'a dit \textbf{de venir}.};
\node[es, right=of f3] (e3) {\esp{Me dijo que \textbf{viniera}.}};
\draw[fleche] (f3) -- (e3);
\node[above=0.15cm of f1, popblue, font=\bfseries] {Français};
\node[above=0.15cm of e1, popgreen!60!black, font=\bfseries] {Español};
\end{tikzpicture}
\legende{Correspondances français/espagnol au style indirect passé : conditionnel, imparfait, subjonctif imparfait.}
\end{popfigure}

\courssub{Méthode de traduction en thème}

\begin{methode}[Traduire une phrase au style indirect passé (français → espagnol)]
\begin{enumerate}
\item \textbf{Repérer le temps du verbe introducteur français.} S'il est au passé (« il a dit », « elle a demandé », « il promit »), la concordance s'impose.
\item \textbf{Identifier le type de subordonnée} :
      \begin{itemize}
      \item Déclaration → \esp{que} + \textbf{indicatif} (imparfait, plus-que-parfait, conditionnel).
      \item Question fermée → \esp{si} + \textbf{indicatif} (imparfait, plus-que-parfait, conditionnel).
      \item Ordre, conseil, volonté (« de + infinitif », « que » + subjonctif) → \esp{que} + \textbf{subjonctif imparfait}.
      \end{itemize}
\item \textbf{Choisir le temps espagnol} dans le tableau de concordance : présent → imparfait ; passé → plus-que-parfait ; futur → conditionnel ; subjonctif présent → subjonctif imparfait.
\item \textbf{Transformer les marqueurs} : demain → \esp{al día siguiente}, ici → \esp{allí}, etc.
\item \textbf{Vérifier} : aucun futur après un verbe introducteur au passé ; subjonctif bien marqué après un ordre ou un souhait ; accord des participes si nécessaire.
\end{enumerate}
\end{methode}

\begin{exoresolu}[Thème guidé]
Énoncé : traduisez en espagnol : « Elle a affirmé qu'elle avait vu l'accident. »
\tcblower
Solution détaillée :
\begin{enumerate}
\item Verbe introducteur : « a affirmé » = passé composé → \esp{afirmó} (prétérit indéfini).
\item Type : déclaration → \esp{que} + indicatif.
\item « avait vu » = plus-que-parfait → l'espagnol conserve le plus-que-parfait de l'indicatif : \esp{había visto}.
\item Traduction : \esp{Afirmó que había visto el accidente.}
\end{enumerate}
Vérification : pas de futur, pas de subjonctif (simple constatation), concordance respectée.
\end{exoresolu}

\begin{exoresolu}[Version guidée]
Énoncé : traduisez en français : \esp{El alcalde prometió que construiría una escuela y preguntó si los vecinos estaban de acuerdo.}
\tcblower
Solution détaillée :
\begin{itemize}
\item \esp{prometió} : prétérit → « a promis » ; verbe introducteur au passé, donc recul des temps en français aussi.
\item \esp{construiría} : conditionnel → « il construirait » (correspondance directe).
\item \esp{preguntó si} : question indirecte → « a demandé si ».
\item \esp{estaban de acuerdo} : imparfait → « étaient d'accord ».
\end{itemize}
Traduction : « Le maire a promis qu'il construirait une école et a demandé si les habitants étaient d'accord. »
\end{exoresolu}

\courssub{Entraînement}

\begin{exercice}[Version : de l'espagnol au français]
Traduisez en français :
\begin{enumerate}
\item \esp{La profesora dijo que el examen sería al día siguiente.}
\item \esp{Mi padre me preguntó si había visto el partido de fútbol.}
\item \esp{El locutor afirmó que la noticia era falsa.}
\item \esp{Ella me pidió que la ayudara con la traducción.}
\end{enumerate}
\end{exercice}

\begin{corrige}[Exercice : version]
\begin{enumerate}
\item La professeure a dit que l'examen aurait lieu le lendemain. (\esp{sería} → conditionnel ; \esp{al día siguiente} → « le lendemain ».)
\item Mon père m'a demandé si j'avais vu le match de football. (\esp{había visto} → plus-que-parfait.)
\item Le présentateur a affirmé que la nouvelle était fausse. (\esp{era} → imparfait.)
\item Elle m'a demandé de l'aider avec la traduction. (\esp{pidió que la ayudara} : le français rend le subjonctif par « de + infinitif ».)
\end{enumerate}
\end{corrige}

\begin{exercice}[Thème : du français à l'espagnol]
Mettez au style indirect passé et traduisez en espagnol :
\begin{enumerate}
\item Il a dit : « Je viendrai demain. »
\item Elle m'a demandé : « As-tu fini tes devoirs ? »
\item Le journaliste a affirmé : « J'ai interviewé le ministre hier. »
\item Ma mère m'a dit : « Rentre tout de suite ! »
\end{enumerate}
\end{exercice}

\begin{corrige}[Exercice : thème]
\begin{enumerate}
\item \esp{Dijo que vendría al día siguiente.} (Futur → conditionnel ; « demain » → \esp{al día siguiente}.)
\item \esp{Me preguntó si había terminado mis deberes.} (Question fermée → \esp{si} ; passé composé → plus-que-parfait.)
\item \esp{El periodista afirmó que había entrevistado al ministro el día anterior.} (« Hier » → \esp{el día anterior} ; \esp{entrevistar a alguien} : \esp{a} personnelle.)
\item \esp{Mi madre me dijo que volviera enseguida.} (Ordre → \esp{que} + subjonctif imparfait \esp{volviera}, jamais l'infinitif.)
\end{enumerate}
\end{corrige}

\begin{aretenir}[L'essentiel de la section]
\begin{itemize}
\item Français et espagnol reculent les temps de la même manière ; le futur devient conditionnel dans les deux langues : \esp{dijo que vendría} = « il a dit qu'il viendrait ».
\item Après un verbe introducteur au passé : jamais de futur, jamais de \esp{mañana} — on dit \esp{al día siguiente}.
\item L'infinitif français (« il m'a dit de venir ») se rend par \esp{que} + subjonctif imparfait : \esp{me dijo que viniera}.
\item En thème : repérer d'abord le temps du verbe introducteur, puis appliquer le tableau de concordance, puis transformer les marqueurs.
\end{itemize}
\end{aretenir}

\newpage

\coursec{Activité d'intégration}

\coursec{Leçon E27 : Traduction — La concordance des temps ; le style direct et indirect}

\courssub{Objectifs de l'activité}

À l'issue de cette activité d'intégration, tu seras capable de :
\begin{itemize}
\item comprendre et prendre en note des déclarations orales au \cle{style direct} (questions et réponses d'une conférence de presse) ;
\item transposer ces déclarations au \cle{style indirect passé} en appliquant correctement la \cle{concordance des temps} à l'indicatif et au subjonctif ;
\item employer avec précision les \cle{verbes introducteurs} (\esp{decir, afirmar, prometer, responder, preguntar, asegurar, anunciar}) et leurs constructions (\esp{decir que}, \esp{preguntar si}, \esp{pedir que} + subjonctif) ;
\item rédiger un court article de presse en espagnol rapportant des paroles, dans un contexte de communication authentique ;
\item évaluer et corriger la production d'un camarade à l'aide d'une grille de relecture fondée sur la concordance des temps.
\end{itemize}

\courssub{Situation}

Tu es journaliste stagiaire au quotidien \emph{La Voz del Litoral}, à Douala. Hier, le ministre de la Santé a tenu une \cle{conférence de presse} (\esp{rueda de prensa}) au Palais des congrès pour annoncer la construction d'un nouvel hôpital de district dans le quartier de Bonabéri. Ton rédacteur en chef t'a envoyé couvrir l'événement : tu as pris des notes pendant que tes confrères posaient leurs questions. Aujourd'hui, tu dois rédiger un \cle{article} de 8 à 10 lignes rapportant les déclarations du ministre et les questions des journalistes, au \cle{style indirect passé}.

Voici les notes que tu as prises pendant la conférence (style direct, tel que prononcé) :

\begin{textebox}[Rueda de prensa : apuntes del periodista]
\textbf{\esp{El ministro de la Salud :}} \esp{«Construiremos un nuevo hospital de distrito en Bonabéri. Las obras empezarán en marzo y terminarán dentro de dos años. El gobierno ya ha firmado el acuerdo con la empresa constructora.»}

\textbf{\esp{Periodista 1 :}} \esp{«¿Cuánto costará el proyecto y quién lo financiará?»}

\textbf{\esp{El ministro :}} \esp{«El proyecto costará menos de lo previsto porque hemos negociado bien. Lo financiará el Estado con la ayuda de socios internacionales.»}

\textbf{\esp{Periodista 2 :}} \esp{«¿Contratarán a médicos cameruneses o extranjeros?»}

\textbf{\esp{El ministro :}} \esp{«Prometo que contrataremos sobre todo a personal camerunés. Además, pido a los jóvenes que se formen en las profesiones sanitarias.»}

\textbf{\esp{Periodista 3 :}} \esp{«¿Está seguro de que las obras terminarán a tiempo?»}

\textbf{\esp{El ministro :}} \esp{«Afirmo que el hospital abrirá sus puertas antes de las próximas elecciones. Confío en que la empresa respete el calendario.»}
\end{textebox}

\textbf{Glossaire :} \esp{las obras} = les travaux ; \esp{el acuerdo} = l'accord ; \esp{lo previsto} = ce qui était prévu ; \esp{contratar} = embaucher ; \esp{formarse} = se former ; \esp{a tiempo} = à temps ; \esp{el calendario} = le calendrier, le planning ; \esp{confiar en que} = espérer, avoir confiance que (+ subjonctif).

\courssub{Consignes}

\begin{enumerate}
\item \textbf{Compréhension et repérage.} Relis les notes de la conférence de presse. Souligne en bleu toutes les \cle{déclarations} du ministre et en rouge toutes les \cle{questions} des journalistes. Pour chaque phrase soulignée, note le temps du verbe conjugué (présent, futur, parfait composé, impératif\ldots).

\item \textbf{Lexique.} Trouve dans les notes l'équivalent espagnol de : « les travaux », « embaucher », « à temps », « moins que prévu », « ouvrir ses portes ». Ajoute ces mots à ton carnet de vocabulaire du module « Médias et communication ».

\item \textbf{Préparation grammaticale.} Le ministre parle au \emph{présent} et au \emph{futur} ; ton article sera écrit \emph{au passé} (\esp{El ministro anunció que\ldots}). À l'aide de la grille de concordance vue en cours, écris sur ton brouillon la transposition des temps suivants : \esp{construiremos} $\rightarrow$ ? ; \esp{empezarán} $\rightarrow$ ? ; \esp{ha firmado} $\rightarrow$ ? ; \esp{costará} $\rightarrow$ ? ; \esp{prometo que contrataremos} $\rightarrow$ ? ; \esp{pido que se formen} $\rightarrow$ ? ; \esp{confío en que respete} $\rightarrow$ ?

\item \textbf{Transposition des questions.} Les questions fermées (réponse oui/non) se rapportent avec \esp{preguntó si} + indicatif ; les questions ouvertes avec \esp{preguntó cuánto / quién / cuándo\ldots} (attention à l'accent sur l'interrogatif). Transpose les trois questions des journalistes au style indirect passé.

\item \textbf{Production écrite.} Rédige ton article de 8 à 10 lignes pour \emph{La Voz del Litoral}. Il doit contenir : un titre, une phrase d'introduction (qui, quoi, où, quand), au moins \textbf{quatre déclarations} rapportées avec des verbes introducteurs variés (\esp{anunciar, afirmar, prometer, asegurar, responder}) et au moins \textbf{deux questions} rapportées (\esp{Un periodista preguntó si\ldots}).

\item \textbf{Relecture en binôme.} Échange ton article avec celui de ton voisin. Corrige-le à l'aide de la grille : \textbf{(a)} chaque verbe introducteur principal est-il au passé ? \textbf{(b)} le temps de la subordonnée est-il bien transposé (imparfait, plus-que-parfait, conditionnel, imparfait de subjonctif) ? \textbf{(c)} les pronoms et possessifs sont-ils adaptés (\esp{yo} $\rightarrow$ \esp{él}, \esp{nuestro} $\rightarrow$ \esp{su}) ? \textbf{(d)} les accents sur les interrogatifs indirects sont-ils présents ?

\item \textbf{Production orale.} Les trois meilleurs articles sont lus à voix haute devant la classe ; les autres élèves doivent identifier, à l'écoute, chaque verbe introducteur et le temps qui le suit.
\end{enumerate}

\courssub{Ressources à mobiliser}

\begin{aretenir}[Grille de concordance (discours rapporté au passé)]
\begin{center}
\begin{tabularx}{\linewidth}{|X|X|}\hline
\rowcolor{popblueL} \textbf{Style direct} & \textbf{Style indirect passé} \\ \hline
Présent : \esp{construyo} & Imparfait : \esp{dijo que construía} \\ \hline
Futur : \esp{construiré} & Conditionnel : \esp{dijo que construiría} \\ \hline
Prétérit indéfini : \esp{construí} & Plus-que-parfait : \esp{dijo que había construido} \\ \hline
Parfait composé : \esp{he firmado} & Plus-que-parfait : \esp{dijo que había firmado} \\ \hline
Présent de subjonctif : \esp{pido que vengan} & Imparfait de subjonctif : \esp{pidió que vinieran} \\ \hline
Impératif : \esp{¡formaos!} & \esp{mandó que se formaran} \\ \hline
\end{tabularx}
\end{center}
\end{aretenir}

\begin{itemize}
\item \textbf{Verbes introducteurs :}
\begin{itemize}
\item \esp{decir / afirmar / asegurar / anunciar / responder / explicar que} + indicatif (présent, imparfait, conditionnel, plus-que-parfait selon la concordance).
\item \esp{prometer que} + indicatif (futur en style direct $\rightarrow$ conditionnel en style indirect passé).
\item \esp{preguntar si} + indicatif (question fermée).
\item \esp{preguntar cuánto / quién / cuándo / dónde / cómo / por qué} + indicatif (question ouverte, accent obligatoire sur l'interrogatif).
\item \esp{pedir / mandar / recomendar / aconsejar / rogar que} + subjonctif (présent $\rightarrow$ imparfait de subjonctif).
\end{itemize}
\item \textbf{Changements de personne et de repères :} \esp{yo} $\rightarrow$ \esp{él} ; \esp{nosotros} $\rightarrow$ \esp{ellos} ; \esp{nuestro} $\rightarrow$ \esp{su} ; \esp{aquí} $\rightarrow$ \esp{allí} ; \esp{mañana} $\rightarrow$ \esp{al día siguiente} ; \esp{ahora} $\rightarrow$ \esp{entonces} ; \esp{este / esta} $\rightarrow$ \esp{ese / esa}.
\item \textbf{Lexique des médias :} \esp{la rueda de prensa} (conférence de presse), \esp{el portavoz} (le porte-parole), \esp{el titular} (le titre), \esp{según el ministro} (selon le ministre), \esp{en declaraciones a la prensa} (dans des déclarations à la presse), \esp{el comunicado} (le communiqué).
\end{itemize}

\begin{attention}[Pièges à éviter]
Ne laisse jamais le futur ou le présent après un verbe introducteur au passé : on écrit \esp{El ministro prometió que contrataría a personal camerunés} (et non \esp{*prometió que contratará}). N'oublie pas l'accent dans les questions indirectes ouvertes : \esp{preguntó cuánto costaría} (et non \esp{*pregunto cuanto costaria}). Enfin, \esp{preguntar si} n'introduit que des questions fermées (réponse oui/non) ; pour une question ouverte, garde le mot interrogatif accentué (\esp{quién, cuánto, cuándo\ldots}). Attention au faux ami \esp{realizar} (faire, exécuter) $\neq$ réaliser (prendre conscience) $\rightarrow$ \esp{darse cuenta de}.
\end{attention}

\courssub{Proposition de production et corrigé commenté}

\begin{exoresolu}[Article modèle pour \emph{La Voz del Litoral}]
Rédige un article de 8 à 10 lignes rapportant la conférence de presse au style indirect passé.
\tcblower
\textbf{\esp{El nuevo hospital de Bonabéri abrirá antes de las elecciones}}

\esp{Ayer, el ministro de la Salud ofreció una rueda de prensa en el Palacio de Congresos de Douala para presentar el proyecto del nuevo hospital de distrito de Bonabéri.}

\esp{El ministro anunció que las obras empezarían en marzo y que terminarían dentro de dos años. Explicó que el gobierno ya había firmado el acuerdo con la empresa constructora.}

\esp{Un periodista preguntó cuánto costaría el proyecto y quién lo financiaría. El ministro respondió que costaría menos de lo previsto porque el Estado había negociado bien, y añadió que lo financiaría el Estado con la ayuda de socios internacionales.}

\esp{Otro periodista preguntó si contratarían a médicos cameruneses o extranjeros. El ministro prometió que contratarían sobre todo a personal camerunés y pidió a los jóvenes que se formaran en las profesiones sanitarias.}

\esp{Finalmente, cuando un colega le preguntó si estaba seguro de que las obras terminarían a tiempo, el ministro afirmó que el hospital abriría sus puertas antes de las próximas elecciones y aseguró que confiaba en que la empresa respetara el calendario.}

\medskip
\textbf{Commentaire des choix de langue :}
\begin{itemize}
\item \textbf{\esp{empezarían}, \esp{terminarían}, \esp{costaría}, \esp{abriría}, \esp{financiaría}, \esp{contratarían}} : futurs du discours direct transposés au \textbf{conditionnel}, conformément à la concordance des temps (futur $\rightarrow$ conditionnel).
\item \textbf{\esp{había firmado}, \esp{había negociado}} : parfaits composés (\esp{ha firmado}, \esp{hemos negociado}) transposés au \textbf{plus-que-parfait} (antériorité par rapport au passé du verbe introducteur).
\item \textbf{\esp{pidió\ldots que se formaran}} et \textbf{\esp{confiaba en que\ldots respetara}} : présents de subjonctif (\esp{se formen}, \esp{respete}) transposés à l'\textbf{imparfait de subjonctif} (règle de concordance : verbe introducteur au passé + subjonctif $\rightarrow$ imparfait de subjonctif).
\item \textbf{\esp{preguntó cuánto costaría}} : question ouverte, mot interrogatif accentué (\esp{cuánto}) ; \textbf{\esp{preguntó si contratarían}} : question fermée introduite par \esp{si}, sans accent.
\item Les pronoms et sujets sont adaptés : \esp{hemos negociado} (nous = le gouvernement) devient \esp{el Estado había negociado} (ou \esp{habían negociado}) ; les verbes introducteurs sont variés (\esp{anunció, explicó, respondió, añadió, prometió, afirmó, aseguró}) pour éviter les répétitions de \esp{dijo}.
\item Le titre de l'article utilise le futur \esp{abrirá} (style journalistique « titre-présent/futur » pour l'actualité immédiate), tandis que le corps de l'article respecte strictement le passé narratif.
\end{itemize}
\end{exoresolu}

\courssub{Bilan de l'activité}

Cette activité t'a permis de mobiliser en situation réelle les trois savoir-faire de la leçon : la \cle{concordance des temps} à l'indicatif (présent $\rightarrow$ imparfait, futur $\rightarrow$ conditionnel, prétérit indéfini/parfait composé $\rightarrow$ plus-que-parfait) et au subjonctif (présent $\rightarrow$ imparfait), la \cle{transposition du style direct au style indirect passé} avec adaptation des pronoms, des possessifs et des repères spatio-temporels, et l'emploi précis des \cle{verbes introducteurs} selon le type d'énoncé (déclaration, question fermée, question ouverte, demande/volonté). Ces compétences sont directement réinvestissables à l'épreuve de \textbf{thème} du baccalauréat, où les phrases à traduire contiennent très souvent un discours rapporté au passé, ainsi qu'en \textbf{compréhension de l'écrit}, où les articles de presse hispanophones rapportent constamment des déclarations officielles. Retiens la démarche rigoureuse : 1) identifier le temps du discours direct, 2) choisir le verbe introducteur adapté au sens, 3) appliquer la grille de concordance, 4) vérifier pronoms, accents, subjonctif et repères spatio-temporels.

\newpage

\coursec{Méthodes et exercices résolus}

\courssub{Fiche méthode 1 : Appliquer la concordance des temps à l'indicatif}

\begin{definition}[La concordance des temps]
La \cle{concordance des temps} (\esp{la concordancia de los tiempos}) est la règle qui impose le temps de la subordonnée en fonction du temps du verbe principal. En espagnol comme en français, quand le verbe principal est au passé, le temps de la subordonnée « recule » d'un cran vers le passé. C'est le cœur du discours rapporté (\esp{estilo indirecto}).
\end{definition}

\begin{methode}[Choisir le temps de la subordonnée à l'indicatif]
Quand le verbe principal est au \textbf{passé} (pretérito indefinido, imperfecto, pluscuamperfecto), le verbe de la subordonnée ne peut pas rester au présent. Procéder en quatre étapes :
\begin{enumerate}
\item \textbf{Repérer le temps du verbe principal} : présent ou passé ? \esp{Dice} (présent) ou \esp{Dijo} (passé) ?
\item \textbf{Identifier le rapport temporel} de la subordonnée par rapport au principal : antériorité (avant), simultanéité (en même temps) ou postériorité (après).
\item \textbf{Appliquer la règle de recul d'un cran} si le principal est au passé :
\begin{center}
\begin{tabularx}{\linewidth}{|X|X|}\hline
\rowcolor{popblueL} Style direct (présent) & Style indirect (passé) \\ \hline
presente \esp{trabaja} & imperfecto \esp{trabajaba} \\ \hline
indefinido / perfecto \esp{trabajó / ha trabajado} & pluscuamperfecto \esp{había trabajado} \\ \hline
futuro \esp{trabajará} & condicional \esp{trabajaría} \\ \hline
imperfecto \esp{trabajaba} & imperfecto (inchangé) \esp{trabajaba} \\ \hline
\end{tabularx}
\end{center}
\item \textbf{Vérifier la cohérence} : relire la phrase entière en espagnol et contrôler que chaque verbe « a reculé » d'un cran.
\end{enumerate}
Réflexe : si le verbe principal est au présent (\esp{dice que...}), la subordonnée garde ses temps normaux ; aucune concordance n'est nécessaire : \esp{Dice que trabaja, que ha trabajado y que trabajará.} « Il dit qu'il travaille, qu'il a travaillé et qu'il travaillera. »
\end{methode}

\begin{aretenir}[Le tableau à connaître par cœur]
\begin{center}
\begin{tabularx}{\linewidth}{|X|X|X|}\hline
\rowcolor{popblueL} Rapport temporel & Principal au présent & Principal au passé \\ \hline
Simultanéité & \esp{dice que viene} & \esp{dijo que venía} \\ \hline
Antériorité & \esp{dice que vino / ha venido} & \esp{dijo que había venido} \\ \hline
Postériorité & \esp{dice que vendrá} & \esp{dijo que vendría} \\ \hline
\end{tabularx}
\end{center}
Le conditionnel espagnol (\esp{vendría}) joue ici le rôle de « futur dans le passé », exactement comme le conditionnel français « viendrait ».
\end{aretenir}

\begin{exoresolu}[Concordance à l'indicatif — thème]
Énoncé : Traduisez en espagnol. « Le journaliste a dit que le match commençait à seize heures et que l'équipe du Cameroun avait gagné la veille. »
\tcblower
Solution détaillée :
\begin{enumerate}
\item Verbe principal : « a dit » = passé composé français, discours rapporté au passé : \esp{dijo} (indefinido de \esp{decir}).
\item « commençait » : simultanéité par rapport à \esp{dijo}. Au style direct, le journaliste avait dit \esp{el partido empieza} (présent). Recul d'un cran : présent $\rightarrow$ imperfecto : \esp{empezaba}.
\item « avait gagné » : antériorité. Style direct : \esp{ganó / ha ganado}. Recul : $\rightarrow$ pluscuamperfecto : \esp{había ganado}.
\item « la veille » : \esp{el día anterior} (et non \esp{ayer}, qui appartient au style direct).
\end{enumerate}
Réponse rédigée : \esp{El periodista dijo que el partido empezaba a las cuatro de la tarde y que la selección de Camerún había ganado el día anterior.}
Conclusion : les deux subordonnées ont reculé d'un cran (présent $\rightarrow$ imparfait ; passé $\rightarrow$ plus-que-parfait), conformément à la concordance après un verbe principal au passé.
\end{exoresolu}

\courssub{Fiche méthode 2 : Appliquer la concordance des temps au subjonctif}

\begin{methode}[Choisir le temps du subjonctif dans la subordonnée]
Après un verbe de volonté, de sentiment ou de doute suivi de \esp{que}, la subordonnée est au subjonctif. Le choix du temps dépend du verbe principal :
\begin{enumerate}
\item \textbf{Principal au présent, au futur, au perfecto ou à l'impératif} : subordonnée au \textbf{presente de subjuntivo} (simultanéité ou postériorité) : \esp{Quiero que vengas.} « Je veux que tu viennes. »
\item \textbf{Principal au passé (imperfecto, indefinido, pluscuamperfecto) ou au condicional} : subordonnée à l'\textbf{imperfecto de subjuntivo} : \esp{Quería que vinieras.} « Je voulais que tu viennes. »
\item \textbf{Antériorité} : utiliser le temps composé correspondant : \esp{Espero que haya llegado.} (présent + perfecto de subjuntivo) « J'espère qu'il est arrivé. » ; \esp{Esperaba que hubiera llegado.} (passé + pluscuamperfecto de subjuntivo) « J'espérais qu'il était arrivé. »
\item \textbf{Vérifier l'irrégularité} des formes en \esp{-ra} : \esp{ser/ir} $\rightarrow$ \esp{fuera, fueras, fuera, fuéramos, fuerais, fueran} ; \esp{haber} $\rightarrow$ \esp{hubiera} ; \esp{saber} $\rightarrow$ \esp{supiera} ; \esp{poder} $\rightarrow$ \esp{pudiera} ; \esp{hacer} $\rightarrow$ \esp{hiciera} ; \esp{decir} $\rightarrow$ \esp{dijera} ; \esp{tener} $\rightarrow$ \esp{tuviera}.
\end{enumerate}
Réflexe de traduction : le subjonctif français « qu'il vînt » (imparfait littéraire) se traduit toujours par l'imperfecto de subjuntivo espagnol, jamais par un passé simple de l'indicatif.
\end{methode}

\begin{propriete}[Rappel : formation de l'imperfecto de subjuntivo]
On part de la 3\textsuperscript{e} personne du pluriel du pretérito indefinido, on retire \esp{-ron} et on ajoute \esp{-ra, -ras, -ra, -ramos, -rais, -ran} (ou la série équivalente en \esp{-se}) : \esp{hablaron} $\rightarrow$ \esp{hablara} ; \esp{comieron} $\rightarrow$ \esp{comiera} ; \esp{vivieron} $\rightarrow$ \esp{viviera}. Toutes les irrégularités de l'indefinido se retrouvent donc ici : \esp{dijeron} $\rightarrow$ \esp{dijera} ; \esp{pusieron} $\rightarrow$ \esp{pusiera} ; \esp{estuvieron} $\rightarrow$ \esp{estuviera}.
\end{propriete}

\begin{exoresolu}[Concordance au subjonctif — thème]
Énoncé : Traduisez en espagnol. « Le ministre a demandé que les journalistes publient la vérité et il était important que tout le monde le sache. »
\tcblower
Solution détaillée :
\begin{enumerate}
\item « a demandé » : verbe de volonté au passé $\rightarrow$ \esp{pidió} (indefinido de \esp{pedir}, verbe à changement radical e $\rightarrow$ i). La subordonnée exige le subjonctif.
\item « publient » : subjonctif français, simultanéité/postériorité par rapport à un principal au passé $\rightarrow$ imperfecto de subjuntivo : \esp{publicaran} (ou \esp{publicasen}).
\item « il était important » : tournure impersonnelle au passé $\rightarrow$ \esp{era importante que} + imperfecto de subjuntivo.
\item « le sache » : \esp{saber} est irrégulier : \esp{supiera}. Sujet « tout le monde » = singulier : \esp{todo el mundo lo supiera}.
\end{enumerate}
Réponse rédigée : \esp{El ministro pidió que los periodistas publicaran la verdad y era importante que todo el mundo lo supiera.}
Conclusion : les deux subordonnées sont à l'imperfecto de subjuntivo parce que les verbes principaux sont au passé ; attention à la forme irrégulière \esp{supiera}.
\end{exoresolu}

\courssub{Fiche méthode 3 : Transposer du style direct au style indirect au passé}

\begin{definition}[Style direct et style indirect]
Le \cle{style direct} (\esp{estilo directo}) reproduit les paroles exactes, entre guillemets : \esp{Ana dijo: «Estoy cansada».} Le \cle{style indirect} (\esp{estilo indirecto}) rapporte les paroles en les intégrant à une subordonnée : \esp{Ana dijo que estaba cansada.} « Ana a dit qu'elle était fatiguée. »
\end{definition}

\begin{methode}[Transformer un discours direct en discours indirect]
Pour rapporter une parole après \esp{dijo que...}, suivre cinq étapes :
\begin{enumerate}
\item \textbf{Supprimer} les guillemets, les deux-points et la majuscule de début de citation ; introduire \esp{que} pour les déclarations.
\item \textbf{Changer les personnes} : première et deuxième personnes $\rightarrow$ troisième personne (ou adaptation logique) : \esp{yo} $\rightarrow$ \esp{él/ella} ; possessifs \esp{mi} $\rightarrow$ \esp{su} ; \esp{nuestro} $\rightarrow$ \esp{su}.
\item \textbf{Faire reculer les temps} d'un cran (voir fiche 1) : présent $\rightarrow$ imparfait ; passé $\rightarrow$ plus-que-parfait ; futur $\rightarrow$ conditionnel ; impératif $\rightarrow$ \esp{que} + imperfecto de subjuntivo (ou \esp{mandó/pidió} + infinitif) : \esp{«Ven aquí»} $\rightarrow$ \esp{Me dijo que fuera allí.} « Il me dit d'y aller. »
\item \textbf{Adapter les marqueurs de temps et de lieu} :
\begin{center}
\begin{tabularx}{\linewidth}{|X|X|}\hline
\rowcolor{popblueL} Style direct & Style indirect \\ \hline
\esp{hoy} (aujourd'hui) & \esp{aquel día / ese día} \\ \hline
\esp{ayer} (hier) & \esp{el día anterior} \\ \hline
\esp{mañana} (demain) & \esp{al día siguiente} \\ \hline
\esp{aquí} (ici) & \esp{allí / allá} \\ \hline
\esp{este / esta} (ce, cette) & \esp{aquel / aquella} \\ \hline
\esp{ahora} (maintenant) & \esp{entonces / en aquel momento} \\ \hline
\end{tabularx}
\end{center}
\item \textbf{Relire} : vérifier la cohérence des personnes, des temps et des adverbes.
\end{enumerate}
\end{methode}

\begin{exoresolu}[Transposition complète — transformation]
Énoncé : Mettez au style indirect, en commençant par \esp{Ana dijo que...} : \esp{« Yo estoy cansada porque trabajé mucho ayer y mañana iré al mercado de Yaoundé con mi hermana. »}
\tcblower
Solution détaillée :
\begin{enumerate}
\item Suppression des guillemets, introduction de \esp{que} : \esp{Ana dijo que...}
\item Personnes : \esp{yo} $\rightarrow$ \esp{ella} (sous-entendu dans la conjugaison) ; \esp{mi hermana} $\rightarrow$ \esp{su hermana}.
\item Recul des temps : \esp{estoy} (présent) $\rightarrow$ \esp{estaba} (imperfecto) ; \esp{trabajé} (indefinido) $\rightarrow$ \esp{había trabajado} (pluscuamperfecto) ; \esp{iré} (futur) $\rightarrow$ \esp{iría} (condicional).
\item Marqueurs : \esp{ayer} $\rightarrow$ \esp{el día anterior} ; \esp{mañana} $\rightarrow$ \esp{al día siguiente}.
\end{enumerate}
Réponse rédigée : \esp{Ana dijo que estaba cansada porque había trabajado mucho el día anterior y que al día siguiente iría al mercado de Yaoundé con su hermana.}
Traduction : « Ana a dit qu'elle était fatiguée parce qu'elle avait beaucoup travaillé la veille et que le lendemain elle irait au marché de Yaoundé avec sa sœur. »
Conclusion : chaque verbe a reculé d'un cran et tous les marqueurs spatio-temporels ont été adaptés ; la répétition de \esp{que} devant la deuxième subordonnée est recommandée en espagnol.
\end{exoresolu}

\courssub{Fiche méthode 4 : Employer les verbes introducteurs}

\begin{methode}[Construire decir, preguntar, afirmar, prometer au discours indirect]
\begin{enumerate}
\item \esp{decir que} + indicatif (déclaration) : \esp{Dijo que vendría.} « Il a dit qu'il viendrait. » Pour un ordre : \esp{decir a alguien que} + subjonctif : \esp{Le dijo que se callara.} « Il lui dit de se taire. »
\item \esp{preguntar} : question totale (réponse oui/non) $\rightarrow$ \esp{preguntó si} + indicatif (jamais d'inversion, jamais de \esp{que}) : \esp{Preguntó si estaba listo.} « Il demanda s'il était prêt. » Question partielle $\rightarrow$ on garde le mot interrogatif avec son accent : \esp{Preguntó dónde vivía.} « Il demanda où il habitait. »
\item \esp{afirmar que} + indicatif : verbe de déclaration forte : \esp{Afirmó que no sabía nada.} « Il affirma qu'il ne savait rien. »
\item \esp{prometer} : deux constructions : \esp{prometer que} + conditionnel (futur reculé) : \esp{Prometió que lo haría.} « Il promit qu'il le ferait. » ; ou \esp{prometer} + infinitif quand le sujet est le même : \esp{Prometió hacerlo.} « Il promit de le faire. »
\end{enumerate}
Réflexe : au discours indirect, \textbf{jamais d'inversion du sujet} après un verbe interrogatif, contrairement au français « demanda-t-il ».
\end{methode}

\begin{attention}[si ou sí ?]
\esp{si} (sans accent) = « si » interrogatif ou hypothétique : \esp{Preguntó si venía.} « Il demanda s'il venait. » \esp{sí} (avec accent) = « oui » ou pronom réfléchi : \esp{Dijo que sí.} « Il dit que oui. » La confusion entre les deux est une faute d'accentuation sanctionnée au Bac.
\end{attention}

\begin{exoresolu}[Verbes introducteurs — thème]
Énoncé : Traduisez en espagnol. a) « Elle m'a demandé si je voulais participer au débat. » b) « Le professeur a demandé où se trouvait la bibliothèque. » c) « Il a promis qu'il écrirait un article sur la paix. »
\tcblower
Solution détaillée :
\begin{enumerate}
\item a) Question totale au style indirect : \esp{preguntar si} + indicatif. « je voulais » : simultanéité, présent reculé $\rightarrow$ imperfecto : \esp{quería}. Réponse : \esp{Me preguntó si quería participar en el debate.}
\item b) Question partielle : on conserve \esp{dónde} avec l'accent (mot interrogatif indirect). « se trouvait » : \esp{estar} ou \esp{encontrarse} à l'imperfecto : \esp{estaba / se encontraba}. Réponse : \esp{El profesor preguntó dónde estaba la biblioteca.} Pas d'inversion du type français ; après le mot interrogatif, l'ordre le plus naturel est ici verbe + sujet.
\item c) \esp{prometer que} + conditionnel (futur reculé) : \esp{escribiría}. Réponse : \esp{Prometió que escribiría un artículo sobre la paz.}
\end{enumerate}
Conclusion : retenir \esp{preguntar si} pour les questions totales, l'accent sur les mots interrogatifs indirects (\esp{dónde, cuándo, cómo, quién}), et le conditionnel après \esp{prometer que} au passé.
\end{exoresolu}

\courssub{Fiche méthode 5 : Traduire une phrase complète au style indirect passé (version et thème)}

\begin{methode}[Stratégie globale de traduction au Bac]
\begin{enumerate}
\item \textbf{Lire toute la phrase} en français et repérer : le verbe introducteur, son temps, le type de subordonnée (déclaration, question, ordre).
\item \textbf{Traduire d'abord le verbe introducteur} au temps juste (souvent indefinido : \esp{dijo, preguntó, afirmó, prometió}).
\item \textbf{Choisir la conjonction} : \esp{que} (déclaration), \esp{si} (question totale), mot interrogatif accentué (question partielle).
\item \textbf{Appliquer la concordance} : faire reculer chaque temps d'un cran ; vérifier les formes irrégulières.
\item \textbf{Adapter} personnes, possessifs et adverbes de temps et de lieu.
\item \textbf{Relire} mentalement à voix haute : accents, accords, ponctuation espagnole (¿ ¡) si la phrase reste interrogative ou exclamative au style direct.
\end{enumerate}
\end{methode}

\begin{exoresolu}[Phrase de type Bac — thème]
Énoncé : Traduisez en espagnol. « Le président a affirmé aux journalistes que la paix était nécessaire et il leur a demandé s'ils publieraient son discours le lendemain. »
\tcblower
Solution détaillée :
\begin{enumerate}
\item Verbes introducteurs : « a affirmé » $\rightarrow$ \esp{afirmó} ; « a demandé » $\rightarrow$ \esp{les preguntó} (pronom pluriel \esp{les} pour « leur »).
\item Première subordonnée : déclaration $\rightarrow$ \esp{que}. « était » : simultanéité, présent reculé $\rightarrow$ imperfecto : \esp{era}. « nécessaire » : \esp{necesaria} (accord féminin avec \esp{la paz}).
\item Deuxième subordonnée : question totale $\rightarrow$ \esp{si} + indicatif. « publieraient » : futur du discours direct (« Publierez-vous... ? ») reculé $\rightarrow$ condicional : \esp{publicarían}.
\item « le lendemain » : \esp{al día siguiente} (style indirect), pas \esp{mañana}.
\end{enumerate}
Réponse rédigée : \esp{El presidente afirmó a los periodistas que la paz era necesaria y les preguntó si publicarían su discurso al día siguiente.}
Conclusion : la phrase combine les trois difficultés du chapitre : verbe introducteur au passé, concordance (imparfait puis conditionnel) et adaptation du marqueur temporel. Une relecture attentive garantit l'accent de \esp{publicarían} et l'accord de \esp{necesaria}.
\end{exoresolu}

\begin{attention}[Les 5 erreurs qui coûtent des points au Bac]
\begin{enumerate}
\item \textbf{Oublier le recul des temps} : écrire \esp{Dijo que viene} au lieu de \esp{Dijo que venía}. Après un verbe principal au passé, le présent de la subordonnée est fautif.
\item \textbf{Traduire « si » interrogatif par \esp{que} ou l'omettre} : « Il a demandé si... » = \esp{Preguntó si...} ; et ne jamais confondre \esp{si} (si) avec \esp{sí} (oui, avec accent).
\item \textbf{Perdre l'accent des mots interrogatifs indirects} : \esp{Preguntó donde vivía} est faux ; il faut \esp{dónde}, \esp{cuándo}, \esp{cómo}, \esp{quién}, même sans point d'interrogation.
\item \textbf{Garder les adverbes du style direct} : traduire « hier » par \esp{ayer} dans un discours rapporté au lieu de \esp{el día anterior} ; « demain » par \esp{mañana} au lieu de \esp{al día siguiente}.
\item \textbf{Calquer le subjonctif français sur l'indicatif espagnol} : après \esp{pidió que}, \esp{era importante que}, l'imperfecto de subjuntivo est obligatoire (\esp{publicaran}, \esp{supiera}) ; un indicatif à cette place coûte le point de grammaire.
\end{enumerate}
\end{attention}

\begin{exercice}[Entraînement — thème]
Traduisez en espagnol :
\begin{enumerate}
\item « Le professeur a dit aux élèves que l'examen aurait lieu le lendemain. »
\item « Ma mère m'a demandé si j'avais acheté du cacao au marché la veille. »
\item « Le capitaine a promis que l'équipe gagnerait le match. »
\item « Il était essentiel que les jeunes écoutent ce message de paix. »
\end{enumerate}
\end{exercice}

\begin{corrige}[Entraînement — thème]
\begin{enumerate}
\item \esp{El profesor dijo a los alumnos que el examen tendría lugar al día siguiente.} (futur « aura lieu » $\rightarrow$ conditionnel \esp{tendría} ; « le lendemain » $\rightarrow$ \esp{al día siguiente}.)
\item \esp{Mi madre me preguntó si había comprado cacao en el mercado el día anterior.} (question totale $\rightarrow$ \esp{si} ; antériorité $\rightarrow$ pluscuamperfecto \esp{había comprado} ; « la veille » $\rightarrow$ \esp{el día anterior}.)
\item \esp{El capitán prometió que el equipo ganaría el partido.} (\esp{prometer que} + conditionnel \esp{ganaría}.)
\item \esp{Era esencial que los jóvenes escucharan ese mensaje de paz.} (tournure impersonnelle au passé $\rightarrow$ imperfecto de subjuntivo \esp{escucharan}.)
\end{enumerate}
\end{corrige}

\newpage

\coursec{Exercices}

\courssub{Je vérifie mes connaissances}

\begin{exercice}[QCM : concordance des temps à l'indicatif]
Choisissez la forme verbale correcte pour compléter chaque phrase au style indirect.
\begin{enumerate}
\item El ministro anunció que el proyecto \underline{\hspace{3cm}} la próxima semana.
      \begin{itemize}
      \item[a)] comenzará
      \item[b)] comenzaría
      \item[c)] comenzaba
      \item[d)] comience
      \end{itemize}
\item Los periodistas confirmaron que el acuerdo \underline{\hspace{3cm}} firmado ayer.
      \begin{itemize}
      \item[a)] ha sido
      \item[b)] había sido
      \item[c)] fue
      \item[d)] sería
      \end{itemize}
\item El entrenador dijo que el equipo \underline{\hspace{3cm}} entrenando desde enero.
      \begin{itemize}
      \item[a)] está
      \item[b)] estaba
      \item[c)] ha estado
      \item[d)] estaría
      \end{itemize}
\item La directora informó que los exámenes \underline{\hspace{3cm}} el lunes siguiente.
      \begin{itemize}
      \item[a)] empiezan
      \item[b)] empezarían
      \item[c)] empezaron
      \item[d)] comenzaran
      \end{itemize}
\end{enumerate}
\end{exercice}

\begin{exercice}[Vrai ou faux justifié : concordance au subjonctif]
Déterminez si la concordance est correcte (V) ou incorrecte (F). Justifiez brièvement en français.
\begin{enumerate}
\item El presidente pidió que los ministros \textbf{presentan} sus dimisiones mañana.
\item La ONU exigió que se \textbf{respeten} los derechos humanos en la zona.
\item El alcalde prometió que \textbf{construiría} el hospital el año que viene.
\item Los padres insistieron en que sus hijos \textbf{estudien} más.
\item El periodista preguntó si el tratado \textbf{entrara} en vigor ya.
\end{enumerate}
\end{exercice}

\begin{exercice}[Vocabulaire : verbes introducteurs et marqueurs]
Associez chaque verbe introducteur à sa construction typique (indicatif, subjonctif, infinitif) et donnez un exemple original en espagnol (contexte camerounais ou hispanophone).
\begin{center}
\begin{tabularx}{\linewidth}{|X|X|X|}
\hline
\rowcolor{popblueL}
Verbe introducteur & Construction & Exemple personnel \\
\hline
decir que & & \\
\hline
preguntar si & & \\
\hline
afirmar que & & \\
\hline
prometer + infinitif & & \\
\hline
ordenar que & & \\
\hline
rogar que & & \\
\hline
\end{tabularx}
\end{center}
\end{exercice}

\begin{exercice}[Conjugaison à compléter : temps du rapporté]
Conjuguez le verbe entre parenthèses au temps approprié dans la subordonnée (style indirect au passé). Le verbe introducteur est au passé.
\begin{enumerate}
\item Dijo que él \underline{\hspace{2cm}} (llegar) al día siguiente.
\item Preguntaron si nosotros \underline{\hspace{2cm}} (conocer) Yaoundé.
\item Afirmó que la cosecha de cacao \underline{\hspace{2cm}} (ser) excelente ese año.
\item Prometió que \underline{\hspace{2cm}} (enviar) el dinero inmediatamente.
\item Ordenó que todos \underline{\hspace{2cm}} (salir) de la sala en ese momento.
\item Rogó que le \underline{\hspace{2cm}} (ayudar) con las maletas.
\end{enumerate}
\end{exercice}

\courssub{Je m'entraîne}

\begin{exercice}[Traduction : phrases au style indirect passé]
Traduisez en espagnol les phrases suivantes en appliquant la concordance des temps et la transformation des marqueurs spatio-temporels.
\begin{enumerate}
\item Le ministre a déclaré qu'il visiterait les écoles du Nord le mois suivant.
\item Les joueurs ont affirmé qu'ils s'entraînaient dur depuis janvier.
\item Le journaliste a demandé si le match avait lieu ce soir-là.
\item La mère a ordonné à son fils de rentrer immédiatement à la maison.
\item Le président a promis qu'il construirait de nouvelles routes l'année suivante.
\item L'enseignant a dit que les examens commenceraient le lundi d'après.
\item L'ambassadeur a demandé si l'accord était déjà signé.
\item Le maire a affirmé que la mairie aiderait les familles sinistrées.
\end{enumerate}
\end{exercice}

\begin{exercice}[Transposition mécanique : style direct → style indirect]
Transformez les phrases suivantes du style direct au style indirect après le verbe introducteur indiqué. Adaptez les pronoms, les temps, les marqueurs temporels et spatiaux.
\begin{enumerate}
\item « Voy a viajar a Guinea Ecuatorial la semana que viene. » → Dijo que…
\item « ¿Llegará el tren a tiempo? » → Preguntó si…
\item « Estamos construyendo un nuevo mercado en Douala. » → Anunciaron que…
\item « ¡Terminad el trabajo ahora! » → Ordenó que…
\item « Mañana firmaré el contrato. » → Prometió que…
\item « ¿Has visto el informe sobre el cacao? » → Quiso saber si…
\item « No puedo ir a la reunión hoy. » → Explicó que…
\item « ¡Cuidado con el tráfico en la autopista! » → Advirtió que…
\end{enumerate}
\end{exercice}

\begin{exercice}[Texte à trous : marqueurs du discours rapporté]
Complétez le discours indirect avec les marqueurs transformés appropriés choisis parmi : \emph{aquel día, al día siguiente, allí, entonces, ese, la semana anterior, el mes siguiente, aquel momento}.
\begin{textebox}[Discours rapporté : conférence de presse]
El portavoz del gobierno declaró en rueda de prensa que \underline{\hspace{2cm}} se había firmado un acuerdo histórico. Añadió que el presidente visitaría la región \underline{\hspace{2cm}} y que se reuniría con los líderes locales \underline{\hspace{2cm}}. Precisó que los fondos llegarían \underline{\hspace{2cm}} y que las obras comenzarían \underline{\hspace{2cm}}. Recordó que \underline{\hspace{2cm}} ya se habían celebrado consultas previas y que \underline{\hspace{2cm}} se tomaría una decisión definitiva.
\end{textebox}
\end{exercice}

\begin{exercice}[Appariement : correspondance temps principale / temps subordonnée]
Reliez chaque temps du verbe introducteur (colonne A) au temps correct de la subordonnée (colonne B) pour une action antérieure, simultanée ou postérieure. Complétez le tableau.
\begin{center}
\begin{tabularx}{\linewidth}{|X|X|X|}
\hline
\rowcolor{popblueL}
Verbe introducteur (passé) & Relation temporelle & Temps de la subordonnée (indicatif) \\
\hline
Dijo que… & Anterioridad & \\
\hline
Dijo que… & Simultaneidad & \\
\hline
Dijo que… & Posterioridad & \\
\hline
Preguntó si… & Anterioridad & \\
\hline
Afirmó que… & Simultaneidad & \\
\hline
Prometió que… & Posterioridad & \\
\hline
\end{tabularx}
\end{center}
\textbf{Consigne :} Donnez un exemple complet en espagnol pour chaque ligne.
\end{exercice}

\courssub{Je me prépare au Bac}

\begin{exercice}[Compréhension et langue : texte de presse original]
Lisez le texte suivant et répondez aux questions en français.
\begin{textebox}[El País (adaptado) — Conferencia sobre paz en Yaoundé]
El Ministro de Relaciones Exteriores anunció ayer que Camerún acogerá el próximo mes una conferencia internacional sobre cultura de la paz. Afirmó que el gobierno había preparado un programa ambicioso y que participarían representantes de más de veinte países africanos y latinoamericanos. Preguntado por los periodistas sobre la seguridad, respondió que las fuerzas del orden garantizarían la protección de todos los delegados. Prometió que se publicaría el calendario definitivo la semana siguiente y rogó a los medios que difundieran la información. Finalmente, ordenó que se reforzaran los controles en los aeropuertos y en las fronteras durante el evento.
\end{textebox}
\textbf{Questions :}
\begin{enumerate}
\item Relevez dans le texte six verbes introducteurs et indiquez pour chacun le temps du verbe principal et le temps du verbe de la subordonnée.
\item Transformez les trois premières phrases du discours indirect au style direct (restituez les paroles supposées du ministre).
\item Expliquez en français les changements de marqueurs temporels opérés entre le style direct hypothétique et le style indirect du texte (\emph{ayer → ? ; el próximo mes → ? ; la semana siguiente → ?}).
\item Traduisez en français le dernier paragraphe (« Preguntado por los periodistas… »).
\end{enumerate}
\end{exercice}

\begin{exercice}[Thème type Bac : 5 phrases au style indirect passé]
Traduisez en espagnol les phrases suivantes. Soignez la concordance des temps, le choix de l'indicatif ou du subjonctif, et la transformation des marqueurs.
\begin{enumerate}
\item Le secrétaire général a annoncé que la conférence s'ouvrirait le lundi suivant.
\item Les organisateurs ont affirmé que tout était prêt depuis la veille.
\item Le chef de la sécurité a demandé si tous les participants avaient reçu leurs badges.
\item Le ministre a promis qu'il recevrait les délégations personnellement à l'aéroport.
\item Le porte-parole a ordonné que personne ne quitte la salle avant la fin de la cérémonie.
\end{enumerate}
\end{exercice}

\begin{exercice}[Version type Bac : paragraphe de presse au style indirect]
Traduisez en français le paragraphe suivant extrait d'un article de presse hispanophone.
\begin{textebox}[La Nación (adaptado) — Cumbre iberoamericana]
El presidente del gobierno español declaró en la cumbre que España aumentaría su cooperación con América Latina en materia educativa. Confirmó que se habían firmado acuerdos con cinco universidades mexicanas y que se negociarían otros con Chile y Argentina el mes siguiente. Advirtió que la crisis climática exigía respuestas conjuntas y pidió que todos los países presentaran propuestas concretas antes de fin de año. Añadió que el rey visitaría la región el año próximo para reforzar los lazos culturales.
\end{textebox}
\textbf{Consigne :} Traduction littéraire soignée, respect des temps (concordance), registre formel. (Environ 80 mots)
\end{exercice}

\begin{exercice}[Sujet de rédaction : expression écrite]
\textbf{Sujet :} « Les médias jouent un rôle essentiel dans la promotion de la culture de la paix. » Rédigez un article de presse (200–250 mots) en espagnol pour un journal scolaire camerounais. Vous rapporterez les propos tenus lors d'une table ronde organisée dans votre lycée sur ce thème.
\textbf{Contraintes obligatoires :}
\begin{itemize}
\item Utilisez au moins 5 verbes introducteurs différents (\emph{decir, afirmar, preguntar, prometer, ordenar, rogar, advertir, explicar…}).
\item Appliquez correctement la concordance des temps (indicatif et subjonctif).
\item Transformez les marqueurs temporels et spatiaux (\emph{hoy → aquel día, mañana → al día siguiente, aquí → allí, ahora → entonces…}).
\item Structurez l'article : titre, chapeau, corps (3 paragraphes), conclusion.
\item Vocabulaire ciblé : \emph{mediación, diálogo, tolerancia, resolución de conflictos, convivencia, fake news, verificación, responsabilidad social}.
\end{itemize}
\end{exercice}

\newpage

\coursec{Corrigés des exercices}

\begin{corrige}[Exercice 1 — QCM : concordance des temps à l'indicatif]
\begin{enumerate}
\item \textbf{b) comenzaría}. \esp{El ministro anunció que el proyecto comenzaría la próxima semana.} Le verbe introducteur est au passé (\esp{anunció}) et l'action rapportée est postérieure : le futur du style direct (\esp{comenzará}) devient \textbf{conditionnel} (\esp{comenzaría}).
\item \textbf{b) había sido}. \esp{Los periodistas confirmaron que el acuerdo había sido firmado ayer.} Antériorité par rapport au moment de la parole : le passé composé (\esp{ha sido}) devient \textbf{plus-que-parfait} (\esp{había sido}).
\item \textbf{b) estaba}. \esp{El entrenador dijo que el equipo estaba entrenando desde enero.} Simultanéité : le présent (\esp{está}) devient \textbf{imparfait} (\esp{estaba}). La périphrase progressive se conserve : \esp{estaba entrenando}.
\item \textbf{b) empezarían}. \esp{La directora informó que los exámenes empezarían el lunes siguiente.} Le présent à valeur de futur du style direct (\esp{empiezan el lunes}) devient \textbf{conditionnel} au style indirect passé. Attention : \esp{comenzaran} (d) est un subjonctif imparfait, impossible après \esp{informar que} qui exprime un fait.
\end{enumerate}
\end{corrige}

\begin{attention}[Piège]
Ne jamais garder le futur après un verbe introducteur au passé : *\esp{dijo que vendrá} est fautif ; il faut \esp{dijo que vendría}.
\end{attention}


\begin{corrige}[Exercice 2 — Vrai ou faux justifié : concordance au subjonctif]
\begin{enumerate}
\item \textbf{F}. Après \esp{pidió que} (passé + volonté), il faut le \textbf{subjonctif imparfait} : \esp{El presidente pidió que los ministros \textbf{presentaran} sus dimisiones al día siguiente.} Le subjonctif présent (\esp{presentan} est d'ailleurs un indicatif !) est impossible ; de plus \esp{mañana} devrait devenir \esp{al día siguiente}.
\item \textbf{F}. Même règle : \esp{exigió que} appelle le subjonctif imparfait : \esp{La ONU exigió que se \textbf{respetaran} los derechos humanos en la zona.}
\item \textbf{V}. \esp{Prometió que construiría} : \esp{prometer que} + fait rapporté postérieur prend le \textbf{conditionnel} (indicatif), pas le subjonctif. La concordance est correcte. (On préférerait \esp{el año siguiente} à \esp{el año que viene}.)
\item \textbf{F}. Après \esp{insistieron en que} au passé : subjonctif imparfait : \esp{Los padres insistieron en que sus hijos \textbf{estudiaran} más.}
\item \textbf{V}. \esp{Preguntó si} + subordonnée interrogative indirecte : l'espagnol admet le subjonctif imparfait (\esp{entrara}) pour une hypothèse postérieure ; l'indicatif imparfait (\esp{entraba}) serait aussi possible pour la simultanéité. La forme proposée est correcte.
\end{enumerate}
\end{corrige}

\begin{corrige}[Exercice 3 — Vocabulaire : verbes introducteurs et marqueurs]
\begin{center}
\begin{tabularx}{\linewidth}{|X|X|X|}
\hline
\rowcolor{popblueL}
Verbe introducteur & Construction & Exemple personnel \\
\hline
decir que & + indicatif & \esp{El profesor dijo que el examen era fácil.} \\
\hline
preguntar si & + indicatif (ou subj. imparfait) & \esp{Preguntó si conocíamos Douala.} \\
\hline
afirmar que & + indicatif & \esp{Afirmó que el cacao camerunés era excelente.} \\
\hline
prometer + infinitif & infinitif (même sujet) & \esp{Prometió estudiar más español.} \\
\hline
ordenar que & + subjonctif & \esp{Ordenó que callaran los alumnos.} \\
\hline
rogar que & + subjonctif & \esp{Rogó que le ayudáramos.} \\
\hline
\end{tabularx}
\end{center}
\textbf{Justification :} les verbes de \textbf{déclaration} (\esp{decir, afirmar, informar}) introduisent un fait → indicatif ; les verbes de \textbf{volonté ou d'influence} (\esp{ordenar, rogar, pedir, exigir}) → subjonctif ; \esp{prometer} suivi du même sujet prend l'infinitif, mais \esp{prometer que} + sujet différent prend l'indicatif (\esp{prometió que vendría}).
\end{corrige}

\begin{corrige}[Exercice 4 — Conjugaison à compléter : temps du rapporté]
\begin{enumerate}
\item \esp{Dijo que él \textbf{llegaría} al día siguiente.} (postériorité → conditionnel)
\item \esp{Preguntaron si nosotros \textbf{conocíamos} Yaoundé.} (simultanéité → imparfait)
\item \esp{Afirmó que la cosecha de cacao \textbf{era} excelente ese año.} (simultanéité → imparfait)
\item \esp{Prometió que \textbf{enviaría} el dinero inmediatamente.} (postériorité → conditionnel)
\item \esp{Ordenó que todos \textbf{salieran} de la sala en ese momento.} (volonté → subjonctif imparfait)
\item \esp{Rogó que le \textbf{ayudaran} / \textbf{ayudasen} con las maletas.} (volonté → subjonctif imparfait)
\end{enumerate}
\end{corrige}

\begin{corrige}[Exercice 5 — Traduction : phrases au style indirect passé]
\begin{enumerate}
\item \esp{El ministro declaró que visitaría las escuelas del Norte el mes siguiente.}
\item \esp{Los jugadores afirmaron que se entrenaban duro desde enero.}
\item \esp{El periodista preguntó si el partido tenía lugar esa noche.}
\item \esp{La madre ordenó a su hijo que volviera inmediatamente a casa.} (ou : \esp{le ordenó volver inmediatamente a casa} — infinitif possible avec complément d'objet)
\item \esp{El presidente prometió que construiría nuevas carreteras el año siguiente.}
\item \esp{El profesor dijo que los exámenes comenzarían el lunes siguiente.}
\item \esp{El embajador preguntó si el acuerdo ya estaba firmado.} (ou \esp{si ya se había firmado el acuerdo} pour l'antériorité)
\item \esp{El alcalde afirmó que el ayuntamiento ayudaría a las familias damnificadas.}
\end{enumerate}
\end{corrige}

\begin{attention}[Points de vigilance]
\esp{ce soir-là} = \esp{esa noche / aquella noche} ; « sinistrées » = \esp{damnificadas} (faux ami : \esp{siniestrada} signifie « sinistre, lugubre ») ; « mairie » = \esp{ayuntamiento}.
\end{attention}


\begin{corrige}[Exercice 6 — Transposition mécanique : style direct → style indirect]
\begin{enumerate}
\item \esp{Dijo que iba a viajar a Guinea Ecuatorial la semana siguiente.} (\esp{voy a} → \esp{iba a} ; \esp{la semana que viene} → \esp{la semana siguiente})
\item \esp{Preguntó si el tren llegaría a tiempo.} (futur → conditionnel ; pas de point d'interrogation au style indirect)
\item \esp{Anunciaron que estaban construyendo un nuevo mercado en Douala.} (présent → imparfait, périphrase conservée)
\item \esp{Ordenó que terminaran el trabajo en ese momento.} (impératif → subjonctif imparfait ; \esp{ahora} → \esp{en ese momento})
\item \esp{Prometió que firmaría el contrato al día siguiente.} (\esp{mañana} → \esp{al día siguiente} ; futur → conditionnel)
\item \esp{Quiso saber si había visto el informe sobre el cacao.} (passé composé → plus-que-parfait)
\item \esp{Explicó que no podía ir a la reunión ese día.} (\esp{puedo} → \esp{podía} ; \esp{hoy} → \esp{ese día})
\item \esp{Advirtió que tuvieran cuidado con el tráfico en la autopista.} (impératif → subjonctif imparfait ; \esp{tuvieran} car le destinataire est pluriel dans \esp{¡Cuidado!} — accepter \esp{tuviera} si contexte singulier)
\end{enumerate}
\end{corrige}

\begin{corrige}[Exercice 7 — Texte à trous : marqueurs du discours rapporté]
\esp{El portavoz del gobierno declaró en rueda de prensa que \textbf{aquel día} se había firmado un acuerdo histórico. Añadió que el presidente visitaría la región \textbf{el mes siguiente} y que se reuniría con los líderes locales \textbf{allí}. Precisó que los fondos llegarían \textbf{la semana siguiente}…}

\textbf{Correction complète (ordre des blancs) :}
\begin{enumerate}
\item \textbf{aquel día} (\esp{hoy} → jour même du discours rapporté)
\item \textbf{el mes siguiente} (\esp{el mes que viene})
\item \textbf{allí} (\esp{aquí} → déplacement spatial)
\item \textbf{al día siguiente} (les fonds arrivent « demain » par rapport au discours)
\item \textbf{entonces} (\esp{ahora} → les travaux commencent « maintenant »)
\item \textbf{la semana anterior} (\esp{la semana pasada} → consultations déjà tenues)
\item \textbf{aquel momento} (\esp{en ese momento} : une décision serait prise sur-le-champ)
\end{enumerate}
(Ordre exact des trous : \esp{aquel día} ; \esp{el mes siguiente} ; \esp{allí} ; \esp{al día siguiente} ; \esp{entonces} ; \esp{la semana anterior} ; \esp{aquel momento}.)
\end{corrige}

\begin{corrige}[Exercice 8 — Appariement : correspondance temps principale / temps subordonnée]
\begin{center}
\begin{tabularx}{\linewidth}{|X|X|X|}
\hline
\rowcolor{popblueL}
Verbe introducteur (passé) & Relation temporelle & Temps de la subordonnée (indicatif) \\
\hline
Dijo que… & Anterioridad & plus-que-parfait : \esp{Dijo que había llegado.} \\
\hline
Dijo que… & Simultaneidad & imparfait : \esp{Dijo que trabajaba.} \\
\hline
Dijo que… & Posterioridad & conditionnel : \esp{Dijo que vendría.} \\
\hline
Preguntó si… & Anterioridad & plus-que-parfait : \esp{Preguntó si habíamos comido.} \\
\hline
Afirmó que… & Simultaneidad & imparfait : \esp{Afirmó que todo iba bien.} \\
\hline
Prometió que… & Posterioridad & conditionnel : \esp{Prometió que pagaría pronto.} \\
\hline
\end{tabularx}
\end{center}
\textbf{Exemples complets acceptés :} toute phrase respectant le schéma \esp{había + participio} (antériorité), imparfait (simultanéité), conditionnel (postériorité). Rappel : le conditionnel espagnol joue ici le rôle de « futur dans le passé », exactement comme en français (« il a dit qu'il viendrait »).
\end{corrige}

\begin{corrige}[Exercice 9 — Compréhension et langue : texte de presse]
\textbf{1. Verbes introducteurs et temps (6 attendus) :}
\begin{itemize}
\item \esp{anunció} (prétérit) + \esp{acogerá} : futur conservé dans le texte journalistique (discours rapporté « frais ») — au style indirect strict : \esp{acogería} (conditionnel).
\item \esp{Afirmó} (prétérit) + \esp{había preparado} (plus-que-parfait, antériorité) et \esp{participarían} (conditionnel, postériorité).
\item \esp{respondió} (prétérit) + \esp{garantizarían} (conditionnel).
\item \esp{Prometió} (prétérit) + \esp{se publicaría} (conditionnel).
\item \esp{rogó} (prétérit) + \esp{difundieran} (subjonctif imparfait, volonté).
\item \esp{ordenó} (prétérit) + \esp{se reforzaran} (subjonctif imparfait).
\end{itemize}
\textbf{2. Restitution au style direct :}
\esp{— Camerún acogerá el próximo mes una conferencia internacional sobre cultura de la paz. El gobierno ha preparado un programa ambicioso y participarán representantes de más de veinte países africanos y latinoamericanos.}\\
(Conditionnel → futur ; plus-que-parfait → passé composé.)
\textbf{3. Marqueurs temporels :} \esp{ayer} (« hier ») reste \esp{ayer} car le journaliste écrit le lendemain — sinon \esp{el día anterior} ; \esp{el próximo mes} (« le mois prochain », style direct) → \esp{el mes siguiente} au style indirect strict ; \esp{la semana siguiente} est déjà la forme rapportée de \esp{la semana que viene}.
\textbf{4. Traduction :} « Interrogé par les journalistes sur la sécurité, il a répondu que les forces de l'ordre garantiraient la protection de tous les délégués. Il a promis que le calendrier définitif serait publié la semaine suivante et a prié les médias de diffuser l'information. Enfin, il a ordonné que les contrôles soient renforcés dans les aéroports et aux frontières pendant l'événement. »\\
\textbf{Barème indicatif (10 pts) :} Q1 : 3 pts (0,5 par verbe correctement analysé) ; Q2 : 2 pts ; Q3 : 2 pts ; Q4 : 3 pts.
\end{corrige}

\begin{corrige}[Exercice 10 — Thème type Bac : 5 phrases au style indirect passé]
\begin{enumerate}
\item \esp{El secretario general anunció que la conferencia se abriría el lunes siguiente.}
\item \esp{Los organizadores afirmaron que todo estaba listo desde la víspera.}
\item \esp{El jefe de seguridad preguntó si todos los participantes habían recibido sus acreditaciones.}
\item \esp{El ministro prometió que recibiría personalmente a las delegaciones en el aeropuerto.}
\item \esp{El portavoz ordenó que nadie saliera de la sala antes del final de la ceremonia.}
\end{enumerate}
\textbf{Barème indicatif (10 pts) :} 2 pts par phrase : 1 pt pour la concordance (conditionnel / imparfait / plus-que-parfait / subjonctif imparfait), 1 pt pour le lexique et les marqueurs (\esp{el lunes siguiente}, \esp{la víspera}, \esp{acreditaciones}).
\end{corrige}

\begin{attention}[Points de vigilance]
« badges » = \esp{acreditaciones} (ou \esp{credenciales}), jamais *\esp{badges} ; « la veille » = \esp{la víspera} ; phrase 5 : subjonctif obligatoire après \esp{ordenar que}, et négation \esp{nadie saliera} sans double négation fautive.
\end{attention}


\begin{corrige}[Exercice 11 — Version type Bac : paragraphe de presse]
\textbf{Traduction proposée :}\\
« Le président du gouvernement espagnol a déclaré lors du sommet que l'Espagne augmenterait sa coopération avec l'Amérique latine dans le domaine de l'éducation. Il a confirmé que des accords avaient été signés avec cinq universités mexicaines et que d'autres seraient négociés avec le Chili et l'Argentine le mois suivant. Il a averti que la crise climatique exigeait des réponses communes et a demandé que tous les pays présentent des propositions concrètes avant la fin de l'année. Il a ajouté que le roi visiterait la région l'année suivante pour renforcer les liens culturels. »\\
\textbf{Barème indicatif (10 pts) :} concordance des temps en français (conditionnel, plus-que-parfait, subjonctif) : 4 pts ; lexique (\esp{cumbre} = sommet, \esp{lazos} = liens, \esp{en materia educativa} = dans le domaine de l'éducation) : 3 pts ; fluidité et registre : 3 pts.
\end{corrige}

\begin{attention}[Faux amis]
\esp{advirtió} = « il a averti » (pas « il a avisé » au sens d'annoncer) ; \esp{exigía} = « exigeait » ; \esp{el año próximo} se rend ici par « l'année suivante » (style rapporté en français).
\end{attention}


\begin{corrige}[Exercice 12 — Sujet de rédaction : expression écrite]
\textbf{Proposition de rédaction modèle (environ 220 mots) :}
\begin{textebox}[Los medios, pilares de la paz — Mesa redonda en el Lycée de Yaoundé]
\textbf{Los medios de comunicación y la cultura de la paz: un compromiso compartido}\\[2pt]
El pasado viernes, nuestro lycée acogió una mesa redonda sobre el papel de los medios en la promoción de la cultura de la paz.\\[4pt]
El director del centro afirmó que los jóvenes pasaban muchas horas frente a las pantallas y explicó que los medios podían fomentar el diálogo o, al contrario, difundir el odio. Preguntado sobre las \emph{fake news}, un periodista invitado respondió que la verificación de las fuentes era una responsabilidad social de todos y advirtió que las noticias falsas destruían la convivencia.\\[4pt]
Una profesora de español declaró que la tolerancia se aprendía también en las redes sociales y pidió que los alumnos respetaran las opiniones ajenas. El delegado de clase prometió que crearía un club de mediación y rogó que todos participaran en la resolución pacífica de los conflictos. Finalmente, el director ordenó que se organizara cada año una semana de la paz en el establecimiento.\\[4pt]
Aquel día comprendimos que informar es un acto de responsabilidad: los medios, cuando sirven a la verdad, construyen la paz.
\end{textebox}
\textbf{Traduction du modèle :} « Vendredi dernier, notre lycée a accueilli une table ronde sur le rôle des médias dans la promotion de la culture de la paix. Le directeur a affirmé que les jeunes passaient de longues heures devant les écrans et a expliqué que les médias pouvaient encourager le dialogue ou, au contraire, diffuser la haine. Interrogé sur les fake news, un journaliste invité a répondu que la vérification des sources était une responsabilité sociale de tous et a averti que les fausses nouvelles détruisaient la convivence. Un professeur d'espagnol a déclaré que la tolérance s'apprenait aussi sur les réseaux sociaux et a demandé que les élèves respectent les opinions d'autrui. Le délégué de classe a promis qu'il créerait un club de médiation et a prié tous de participer à la résolution pacifique des conflits. Enfin, le directeur a ordonné qu'une semaine de la paix soit organisée chaque année dans l'établissement. Ce jour-là, nous avons compris qu'informer est un acte de responsabilité : les médias, quand ils servent la vérité, construisent la paix. »\\
\textbf{Barème indicatif (20 pts) :} respect des contraintes (5 verbes introducteurs variés, marqueurs transformés, structure) : 6 pts ; concordance des temps et modes : 6 pts ; vocabulaire ciblé : 4 pts ; cohérence et correction linguistique générale : 4 pts.
\end{corrige}

\begin{attention}[Points de vigilance]
Vérifier : imparfait pour la simultanéité (\esp{pasaban, era}), conditionnel pour la postériorité (\esp{crearía}), subjonctif imparfait après les verbes de volonté (\esp{respetaran, participaran, se organizara}) ; transformer \esp{hoy} → \esp{aquel día}, \esp{aquí} → \esp{allí / en el centro}.
\end{attention}


\newpage

\coursec{Fiche bilan}

\begin{popfigure}
\begin{tikzpicture}[mindmap, grow cyclic, every node/.style={concept, align=center, font=\small},
root concept/.append style={concept color=popblue, font=\bfseries},
level 1/.append style={level distance=3.4cm, sibling angle=51.4},
level 2/.append style={level distance=2.1cm, font=\scriptsize}]
\node[root concept, align=center]{Concordance des temps\\et style indirect}
  child[concept color=poporange]{ node[align=center]{Indicatif\\(déclaration)}
    child { node[align=center]{pret. imperfecto\\+ imperfecto} }
    child { node[align=center]{futuro $\to$\\condicional} }
  }
  child[concept color=popgreen]{ node[align=center]{Subjonctif\\(volonté, doute)}
    child { node[align=center]{pret. imperfecto\\de subjuntivo} }
    child { node {hubiera / hubiese} }
  }
  child[concept color=poppurple]{ node[align=center]{Transposition\\direct $\to$ indirect}
    child { node[align=center]{que / si /\\interrogatif} }
  }
  child[concept color=poppink]{ node[align=center]{Verbes\\introducteurs}
    child { node[align=center]{decir, preguntar,\\afirmar, prometer} }
  }
  child[concept color=popgold]{ node[align=center]{Changements\\de personne}
    child { node[align=center]{yo $\to$ él/ella\\me $\to$ le/se} }
  }
  child[concept color=popblue!70]{ node[align=center]{Repères\\temporels}
    child { node[align=center]{hoy $\to$ aquel día\\mañana $\to$ al día siguiente} }
  }
  child[concept color=popmuted]{ node[align=center]{Pièges\\francophones}
    child { node[align=center]{pas de futur\\après si} }
  };
\end{tikzpicture}
\legende{Carte mentale de la leçon : concordance des temps et style indirect}
\end{popfigure}

\begin{bilan}
\textbf{Lexique clé (verbes introducteurs et repères)}

\begin{center}
\begin{tabularx}{\linewidth}{|X|X|}\hline
\rowcolor{popblueL} \textbf{Español} & \textbf{Français} \\ \hline
decir (que) & dire (que) \\ \hline
preguntar (si / qué / dónde) & demander (si / quoi / où) \\ \hline
afirmar / asegurar & affirmer / assurer \\ \hline
prometer / jurar & promettre / jurer \\ \hline
contestar / responder & répondre \\ \hline
exclamar / gritar & s'exclamer / crier \\ \hline
añadir / aclarar & ajouter / préciser \\ \hline
ordenar / mandar (+ inf.) & ordonner \\ \hline
aconsejar / prohibir (+ inf. ou subj.) & conseiller / interdire \\ \hline
\end{tabularx}
\end{center}

\textbf{Repères temporels à transformer}

\begin{center}
\begin{tabularx}{\linewidth}{|X|X|}\hline
\rowcolor{popblueL} \textbf{Style direct} & \textbf{Style indirect passé} \\ \hline
hoy / esta mañana & aquel día / esa mañana \\ \hline
ayer & el día anterior / la víspera \\ \hline
mañana & al día siguiente \\ \hline
ahora & entonces / en ese momento \\ \hline
aquí & allí / allá \\ \hline
este / esta & ese / esa, aquel / aquella \\ \hline
la semana pasada & la semana anterior \\ \hline
\end{tabularx}
\end{center}

\textbf{Règles de grammaire à connaître par cœur}

\begin{center}
\begin{tabularx}{\linewidth}{|X|X|}\hline
\rowcolor{popblueL} \textbf{Style direct} & \textbf{Style indirect (verbe au passé)} \\ \hline
presente : \esp{dice: «trabajo»} & imperfecto : \esp{dijo que trabajaba} \\ \hline
pretérito perfecto / indefinido & pluscuamperfecto : \esp{dijo que había trabajado} \\ \hline
futuro : \esp{«trabajaré»} & condicional : \esp{dijo que trabajaría} \\ \hline
presente de subjuntivo & imperfecto de subjuntivo : \esp{pidió que trabajara} \\ \hline
imperativo : \esp{«¡trabaja!»} & infinitivo ou subj. imperfecto : \esp{le mandó trabajar} \\ \hline
\end{tabularx}
\end{center}

\textbf{Méthodes express}
\begin{itemize}
\item \textbf{Étape 1} : repérer le verbe introducteur et son temps (présent ou passé).
\item \textbf{Étape 2} : choisir le connecteur : \esp{que} (déclaration), \esp{si} (question fermée), mot interrogatif accentué \esp{qué, dónde, cuándo} (question ouverte).
\item \textbf{Étape 3} : reculer les temps d'un cran (présent $\to$ imparfait ; futur $\to$ conditionnel ; passé $\to$ plus-que-parfait).
\item \textbf{Étape 4} : adapter personnes, possessifs et repères temporels ; supprimer les guillemets et les deux-points.
\end{itemize}
\end{bilan}

\begin{aretenir}[Checklist avant le Bac]
\begin{itemize}
\item[$\square$] Je connais les trois connecteurs du style indirect : \esp{que}, \esp{si}, mot interrogatif accentué.
\item[$\square$] Je sais que le futur du style direct devient \textbf{conditionnel} au style indirect passé.
\item[$\square$] Je sais que le présent devient \textbf{imparfait de l'indicatif} après un verbe au passé.
\item[$\square$] Je transforme le passé composé / passé simple en \textbf{plus-que-parfait} (\esp{había dicho}).
\item[$\square$] Je transforme le présent de subjonctif en \textbf{imparfait de subjonctif} (\esp{trabajara / trabajase}).
\item[$\square$] Je sais rendre l'impératif par \textbf{l'infinitif} (\esp{le mandó callarse}) ou le subjonctif imparfait.
\item[$\square$] Je change correctement les personnes et les possessifs (\esp{mi} $\to$ \esp{su}, \esp{me} $\to$ \esp{le}).
\item[$\square$] Je transforme les repères : \esp{hoy} $\to$ \esp{aquel día}, \esp{mañana} $\to$ \esp{al día siguiente}.
\item[$\square$] Je n'emploie jamais de futur après \esp{si} interrogatif au style indirect.
\item[$\square$] Je conjugue sans erreur \esp{decir} : \esp{dijo, dijeron, había dicho, diría}.
\item[$\square$] Je vérifie les accents sur les mots interrogatifs : \esp{qué, cómo, dónde, cuándo}.
\item[$\square$] Je relis ma traduction en vérifiant la concordance verbe introducteur $\to$ verbe subordonné.
\end{itemize}
\end{aretenir}

\findecours

\end{document}
